% Standalone source. All manuscript text, references, and printed listings are embedded.
\documentclass[reqno]{amsart}
\usepackage[T1]{fontenc}
\usepackage{lmodern}
\usepackage{amsmath,amssymb,mathtools,mathrsfs}
\usepackage{booktabs,longtable,array}
\usepackage{enumitem}
\usepackage{xcolor}
\usepackage{microtype}
\usepackage[unicode,pdfusetitle,hidelinks]{hyperref}
\usepackage{listings}
\setlength{\emergencystretch}{2em}
\setcounter{tocdepth}{1}
\allowdisplaybreaks[2]
\newtheorem{theorem}{Theorem}[section]
\newtheorem{proposition}[theorem]{Proposition}
\newtheorem{lemma}[theorem]{Lemma}
\newtheorem{corollary}[theorem]{Corollary}
\theoremstyle{definition}
\newtheorem{definition}[theorem]{Definition}
\theoremstyle{remark}
\newtheorem{remark}[theorem]{Remark}
\newtheorem{example}[theorem]{Example}
\newcommand{\C}{\mathbb C}
\newcommand{\R}{\mathbb R}
\newcommand{\N}{\mathbb N}
\newcommand{\D}{\mathbb D}
\newcommand{\HH}{\mathbb H}
\DeclareMathOperator{\Tr}{Tr}
\DeclareMathOperator{\Id}{Id}
\DeclareMathOperator{\supp}{supp}
\DeclareMathOperator{\rank}{rank}
\DeclareMathOperator{\diag}{diag}
\DeclareMathOperator{\osc}{osc}
\DeclareMathOperator{\dist}{dist}


\title[Stability generators, spectral gaps, and computation]{Generators of stability-preserving semigroups,\ spectral gaps, and classical ground-state computation}
\author{Dongsheng Wei}
\address{Independent Researcher}
\email{dongshengwei2025@icloud.com}
\date{}
\subjclass[2020]{Primary 47D06, 30C15; Secondary 15A42, 81P45, 68Q25}
\keywords{stable polynomial, generator classification, spectral gap, Lee--Yang theorem, quantum Hamiltonian, classical algorithm}
\hypersetup{pdftitle={Generators of stability-preserving semigroups, spectral gaps, and classical ground-state computation},pdfauthor={Dongsheng Wei}}
\begin{document}
\begin{abstract}
We classify all linear generators of complex stability-preserving
semigroups on finite boxes of polynomial degrees. Their canonical
differential order is at most two, and their coefficients have
nonnegativity certificates of size independent of the capacities.
For binary variables, preservation is decidable in polynomial bit time
from a supplied sum of one-site and two-site matrices.
Positive coordinate damping makes the actual semigroup kernel strictly
zero-free and gives a sharp, capacity-independent real-part spectral
gap, extending the Hermitian \(2\mu\) gap of Bravyi, Gosset, Liu, and
Wong to arbitrary complex stability-preserving generators.
For bounded-degree local inputs at fixed norm and positive damping,
the resulting analytic domain yields deterministic classical
algorithms for the principal eigenvalue and relative stable product
tests of its eigenvectors. Applied to Suzuki--Fisher Hamiltonians,
these give ground-energy and amplitude queries, product-event
probabilities, and classical randomized sampling of adaptive equatorial
measurement records.
A field perturbation also gives zero-field energy-value approximation
schemes. We further prove an all-temperature Hamiltonian converse
at finite spins, weighted token-graph spectral concavity resolving the
adjacency and signless-Laplacian conjectures of Apte, Parekh, and Sud,
and an obstruction separating connected static preservers from
products of preserving flows.
\end{abstract}
\maketitle
\clearpage
\tableofcontents
\section{Introduction}\label{intro:section}

A polynomial with no zero in a prescribed domain imposes a global
constraint on its coefficients. When the coefficients encode a matrix
or a quantum state, that constraint can determine which interactions
are possible and which spectral quantities can be computed. We study
continuous linear evolutions that preserve all complex stable
polynomials on a finite box of degrees, and use their infinitesimal
structure to obtain spectral bounds and classical algorithms.

The classification of individual stability-preserving operators by their
symbols gives a starting point for this question
\cite{BB09,BB10,BB10II}. A semigroup imposes an additional condition:
every operator $e^{tA}$, including those at arbitrarily small positive
times, must be a preserver. This infinitesimal requirement is especially
restrictive at collisions of roots. It is also useful. A generator
description can identify all admissible interactions, distinguish
physical positivity from complex stability, and supply a spectral
structure that is not apparent from the entries of the operator.

The one-variable theory begins with the multiplier-sequence
classification of P\'olya and Schur \cite{PolyaSchur14}. Constant-coefficient
differential operators and heat-type deformations preserving real zeros
form another classical part of the subject; see Craven and Csordas
\cite{CravenCsordas94}. Our question concerns arbitrary endomorphisms of
finite multivariable boxes, whose variable coefficients and boundary
degenerations must be controlled simultaneously.

Infinitesimal stability preservers have also been studied directly.
Purbhoo constructs a family on multiaffine polynomials indexed by
matrices with real diagonal and nonnegative off-diagonal entries
\cite[Proposition~4.5]{Purbhoo18}. His result gives sufficient
generators. The structural theorem below supplies necessary and
sufficient conditions for arbitrary endomorphisms on every finite
coordinate box.

This paper develops these consequences in the following chain:
\[
 \begin{gathered}
 \text{root-collision constraints}
 \ \Longrightarrow\
 \text{generator structure}\\
 \ \Longrightarrow\
 \text{strict complex spectral control}
 \ \Longrightarrow\
 \text{classical computation}.
 \end{gathered}
\]
The finite-box classification is the structural result. Strict damping
turns it into a spectral and analytic tool, and a compact local input
makes that tool algorithmic. The all-temperature Hamiltonian converse
is a structural application; comparisons between homogeneous sectors
and the distinction between static preservers and generated flows give
further consequences of the same classification.

\subsection{The finite-dimensional structure theorem}

Let $V_\kappa$ be the space of polynomials with coordinatewise degree
bounds $\kappa=(\kappa_1,\ldots,\kappa_n)$. Every linear operator on this
space has a canonical normally ordered differential representation.
The first main theorem describes which such operators generate
stability-preserving semigroups, without assuming positivity, locality,
or a differential order in advance.

In real half-plane coordinates, the answer is second order. If
\[
 A=\sum_{\alpha\leq\kappa}P_\alpha(z)\partial^\alpha
\]
is the canonical expression, preservation forces $P_\alpha=0$ for
$|\alpha|>2$. Each diagonal second-order coefficient is a polynomial
of its own variable and is nonpositive on the real line. Each mixed
coefficient is a polynomial of the corresponding pair of variables and
is nonpositive on the real plane. Together with the condition that
$A$ acts on $V_\kappa$, these conditions are also sufficient.
This is Theorem~\ref{gen:main-real}. The finite-box constraint is
substantial: it couples the highest terms of the lower-order
coefficients to the second-order part.

For complex generators there is one further component.
Every generator is the sum of such a real generator and
\begin{equation}\label{intro:imaginary}
 \mathrm i\left(cI+\sum_i
   \left(q_i(z_i)\partial_i-\frac{\kappa_i}{2}q_i'(z_i)\right)\right),
 \qquad c\in\R,\quad q_i\geq0\ \hbox{on }\R,
\end{equation}
where each $q_i$ is a real polynomial of degree at most two.
Conversely every such sum generates a complex stability-preserving
semigroup (Theorem~\ref{gen:main-complex}).
The operators in \eqref{intro:imaginary} are the infinitesimal
one-variable M\"obius motions pointing into the half-plane.
Thus arbitrary interactions in the imaginary part are excluded.

The proof has two different ingredients. Root-cluster asymptotics force
the order bound and the signs. The local multiple-root mechanism has
an antecedent in the variational analysis of Burke and Overton
\cite[Theorem~1.2 and Lemma~1.4]{BurkeOverton01}.
A finite-dimensional algebraic
decomposition then reduces sufficiency to explicit binary primitives.
The required certificates use Gram matrices of orders at most four,
independently of the numerical capacities. Polarization gives an exact
binary-slot realization, with the binomial factors retained.
These are generator statements; they are not obtained by differentiating
a smooth boundary of the stable set. That boundary is singular, and its
higher-multiplicity strata matter.

For a sum of one-site and two-site matrices on binary variables, the
classification is effectively checkable. Corollary~\ref{gen:recognition}
decides preservation of the entire semigroup in polynomial bit time
from this local description. It checks the sum of the local terms,
allowing cancellations between them, and never constructs the full
exponential-size matrix.

\subsection{Strict damping and a complex spectral gap}

Pass to the unit disk by a Cayley transform. Suppose $A_0$ generates
a semigroup preserving all complex disk-stable polynomials in
$V_\kappa$, and set
\begin{equation}\label{intro:damping}
 A=A_0-2\sum_i\mu_iE_i,\qquad E_i=z_i\partial_i,\qquad \mu_i>0.
\end{equation}
The damping has a simple interpretation: it contracts each polynomial
variable towards the center of the disk. The conclusion is stronger
than preservation of the open disk. For every $t>0$ the actual
binomial-weighted kernel of $e^{tA}$ is nonzero on the closed unit
polydisk, and it is stable in a polydisk of some radius $r_t>1$.
For any fixed positive time, a common strict radius works at all later
times (Theorem~\ref{spec:strict}).

Consequently $A$ has an algebraically simple principal eigenvalue $a_*$
and
\begin{equation}\label{intro:gap}
 \Re a_*-\Re a\geq2\min_i\mu_i
 \quad\text{for every other eigenvalue }a.
\end{equation}
The constant is sharp, and the bound is independent of the capacities
and number of variables (Theorem~\ref{spec:gap}). Neither self-adjointness
nor coefficient positivity is assumed. The principal right and left
objects are strictly stable.

The proof uses the finite-dimensional complex-cone theory of Rugh
\cite{Rugh10}, together with a damping argument and contraction of stable
kernels. The sharp radius-to-spectrum mechanism has a Hermitian
Lee--Yang antecedent in \cite{BGLW26}. Our extension concerns arbitrary
complex finite-box stable generators and their actual kernels.
The corresponding dynamical estimate concerns the oscillation of a
quotient $h/p$ over the polydisk. Its contraction, including
time-dependent flows, is developed in Section~\ref{spec:section}.
This strengthens the spectral conclusion; the algorithms below use
the spectral gap and zero-free kernel, rather than this oscillation estimate.
The projective-metric and oscillation constructions themselves have
antecedents in \cite{Dubois09,GaubertQu15}.

\subsection{Comparing homogeneous sectors}

Coefficient positivity gives another spectral consequence. Suppose a
finite-box generator preserves real stability, has nonnegative
off-diagonal entries, and preserves every homogeneous degree.
The spectral bounds of its degree blocks form a concave sequence
(Theorem~\ref{sector:growth}). This also holds for nonself-adjoint,
reducible blocks. The proof evolves a single polynomial with positive
coefficients, applies Newton's inequalities to its diagonal restriction,
and takes exponential growth rates. The underlying growth-exponent
principle is classical \cite{Branden10Discrete}.

For a nonnegatively weighted graph \(G\), let \(A_k,D_k\) be the
adjacency and degree matrices of its \(k\)-token graph. The generator
classification applies to the whole family
\[
 B_k(a)=A_k+aD_k,\qquad -1\le a\le1.
\]
Its largest eigenvalues satisfy
\[
 2\lambda_{\max}(B_k(a))
 \ge\lambda_{\max}(B_{k-1}(a))+\lambda_{\max}(B_{k+1}(a)),
\]
and are symmetric under \(k\leftrightarrow n-k\).
They are therefore nondecreasing up to half filling. This proves
the adjacency and signless-Laplacian monotonicity assertions of
\cite[Conjectures~5--6]{APS26}, with a stronger weighted concavity
statement and a sharp universal interval for \(a\)
(Theorem~\ref{sector:token}).
The \(a=0\) semigroup is already contained in Purbhoo's construction
\cite{Purbhoo18}; the new conclusion compares its particle-number
sectors. The corresponding EPR interpretation uses the established
token-graph correspondence of \cite{APS26}.

\subsection{A classical algorithm beyond self-adjoint operators}

Here is the algorithmic consequence in a form that specifies its input
model. On $n$ binary variables let $A_0$ be given as a sum of one-site
and two-site matrices with Gaussian-rational entries. Suppose the
interaction graph has degree at most $d$, each local matrix has norm
at most $J$, and $e^{tA_0}$ preserves all complex disk-stable
multiaffine polynomials. The preservation hypothesis concerns the whole generator; individual
displayed summands need not be preservers. It is decidable in
polynomial bit time by Corollary~\ref{gen:recognition}.
For fixed $d,J$, and a fixed rational $\mu>0$, the principal eigenvalue of
\[
 A_0-2\mu\sum_iE_i
\]
can be approximated to absolute complex error $\delta$ by a deterministic
classical algorithm in time
\[
 \operatorname{poly}_{d,J,\mu}
       \bigl(n,\delta^{-1},\text{input bit length}\bigr).
\]
This is Theorem~\ref{alg:complex}. The matrix has exponentially many
entries, but the input is its sparse local description. There is no
claim of a running time polynomial uniformly in $\mu^{-1}$.

The same local input also gives relative access to the principal
eigenvectors. Normalize the right and transpose-left vectors $v,w$
to have vacuum coefficient one. For any tensor product $B=\bigotimes_iB_i$
of Gaussian-rational $2$-by-$2$ matrices with $(B_i)_{00}=1$ and
disk-stable kernels, Theorem~\ref{alg:product-test} computes
\[
 w^{\mathsf T}Bv
 \quad\hbox{and}\quad
 \Tr(B\Pi)=\frac{w^{\mathsf T}Bv}{w^{\mathsf T}v},
 \qquad \Pi=\frac{vw^{\mathsf T}}{w^{\mathsf T}v},
\]
to relative complex error $\epsilon$ in
$\operatorname{poly}_{d,J,\mu}(n,\epsilon^{-1},L_{\rm in})$ time.
Both quantities are nonzero. This is scalar query access; no
ordinary-norm conditioning bound for $\Pi$ is required.

The algorithm uses two distinct analytic regions. A local expansion
at a product vacuum has a size-independent convergence radius.
The perturbation-coefficient algorithm of Bravyi, DiVincenzo, and
Loss \cite{BDL08} supplies these coefficients. A complexification
argument extends the coefficient bounds from Hermitian perturbations
to arbitrary complex ones. The stability gap then continues the same
principal branch through a right half-plane. An explicit disk map
joining these regions turns logarithmically many local coefficients
into a value estimate at the desired fixed damping. We give the
branch-identification argument, the coefficient construction, and the
algorithm in Section~\ref{alg:section}; the detailed rational
bit bounds are collected in Appendix~\ref{arith:section}.

For the state queries, the extra step is to obtain a nonvanishing
scalar test on that same type of domain. The creation-coefficient
bounds in \cite{BDL08} give a uniform vacuum disk for both eigenvectors.
The half-plane continuation and strict disk contraction then show
that $w^{\mathsf T}Bv$ never vanishes. Its logarithm has a one-sided
real-part bound of order $n$ throughout the domain
(Lemma~\ref{alg:test-domain}). Connected coefficients of this logarithm
can be computed from subsystems of logarithmic size, without requiring
the induced subsystems to satisfy the global preservation promise.

Connected-cluster computation combined with analytic continuation is an
established strategy
\cite{PatelRegts17,YYZ22,WildAlhambra23}.
Its point is that complex stability
provides a global analytic domain for a class of sparse operators that
need not be self-adjoint, positive, or Markovian. The computational
statement uses that domain to go beyond the perturbative neighborhood
where the coefficient expansion is initially justified.

\subsection{Quantum consequences and their scope}

For qubits, consider
\begin{equation}\label{intro:sf}
 \begin{gathered}
 H=-\sum_{ij}
 \left(J^z_{ij}Z_iZ_j+
       \sum_{a,b\in\{x,y\}}J^{ab}_{ij}\sigma_i^a\sigma_j^b\right)
   -\sum_i(h_i^xX_i+h_i^yY_i+h_i^zZ_i),
 \\
 J^z_{ij}\geq\|J^\perp_{ij}\|_{\rm op},\quad h_i^z\geq0 ,
 \end{gathered}
\end{equation}
with real coefficients and
$J^\perp_{ij}=(J^{ab}_{ij})_{a,b\in\{x,y\}}$.
This is the phase-rotated Suzuki--Fisher family.
Its Lee--Yang sufficiency, including the full transverse operator-norm
condition, is classical: see \cite[equation~(2.47)]{SuzukiFisher71}
and \cite[Theorem~2]{Dunlop79}. Asano's contemporaneous account already
uses a full matrix generating function and its composition law for the
spin-one-half XXZ subfamily \cite[equation~(1.15), Theorem~1]{AsanoJapanese70}.
The tensor language of \cite{WBG26} is a modern formulation of this
tradition.
The generator theorem gives the converse: up to an additive
scalar, these are precisely the finite-qubit Hamiltonians whose
full Gibbs tensor is Lee--Yang at every temperature
(Theorem~\ref{eq:sf}). In particular no genuine higher-body term is
possible under that full-tensor hypothesis. Ruelle's all-temperature converse
\cite[Theorem 9]{Ruelle10} already forces pair interactions for classical
coefficientwise Boltzmann Lee--Yang polynomials. The object here is the
full matrix tensor of an arbitrary Hermitian Hamiltonian, including
noncommuting interactions.

For rational Pauli coefficients, bounded degree and bounded local
strength, and longitudinal fields bounded below by a fixed $\mu>0$,
Theorem~\ref{alg:sf} gives deterministic classical ground-energy
approximation to arbitrary absolute accuracy.
Theorem~\ref{alg:sf-state} additionally computes nonzero relative
amplitudes against product states
$(|0\rangle+s_i|1\rangle)/\sqrt{1+|s_i|^2}$,
$s_i\in\mathbb Q(i)$ and $|s_i|\le1$,
and relative probabilities of the corresponding product events.
It gives a classical randomized sampler for adaptive equatorial
measurements, with any prescribed total variation error. Each site
is measured at most once; the strategy's own running time and
the bit lengths of its phases are included explicitly.
At nonnegative fields, adding a fixed field of size proportional
to $\varepsilon$ gives an additive $\varepsilon n$ energy-value
approximation scheme (Theorem~\ref{alg:value}).
This covers bounded-degree EPR models on arbitrary graphs and,
after a local unitary transformation, bounded-degree bipartite
Quantum MaxCut. For both unweighted classes, the estimate gives a deterministic
$(1-\varepsilon)$ approximation to the optimum value, with polynomial
running time for every fixed $\varepsilon>0$.

Earlier classical algorithms treat important subclasses of quantum
ferromagnets \cite{BG17,HMS20}. The broader classical computation
question is stated in \cite{BGLW26}; a related interpolation strategy
was announced in \cite{WongTalk23}, with further background in
\cite{WongThesis26}.
Relative \(X\)-basis amplitudes, phases, and
measurement marginals already appear in
\cite[Section~4]{WBG26}, for general Lee--Yang states with a fixed
spatial zero-free radius greater than one and state-dependent
quasipolynomial circuits. Here the local Hamiltonian is the input,
the field is any fixed positive constant, and the complexity is
polynomial; a size-independent spatial radius margin is not assumed.
The new analytic input is the nonvanishing reciprocal-field domain
for the principal-state tests. It is distinct from the thermal
partition-function domains used for product measurements in
\cite[Theorem~15]{YYZ22}.
The zero-field energy-value scheme does not
prepare an approximately optimal quantum state. It does not solve
the unrestricted inverse-polynomial-accuracy EPR problem
\cite{MarwahaSud26}, and it is not a replacement for approximation
algorithms that also construct states \cite{ALMPSS26}.

Positive damping also proves strict zero-freeness of the actual Gibbs
kernel and of its unique ground-state polynomial, addressing the
qualitative radius question in \cite{WBG26}.
The stronger proposed radius $s^{-1/2}$ for a deformed EPR ground
state is false. Appendix~\ref{cx:section} gives an unweighted
six-vertex tree at $s=1/2$, together with an exact rational certificate
of a zero strictly inside that radius. It leaves the qualitative
radius-greater-than-one conclusion intact.

The converse also extends to arbitrary finite local spins
(Section~\ref{spin:section}). The resulting Hamiltonian cone permits
longitudinally dominated pair interactions and the corresponding
on-site quadratic terms, with the coherent binomial normalization
made explicit. This is an equilibrium classification, separate from
classification of completely positive quantum channels.

\subsection{Static preservers and the range of the theory}

The generator cone does not exhaust the connected semigroup of static
preservers. On the multiaffine space, every invertible complex
 disk-stability preserver can be joined to the identity by a path of
preservers. Nevertheless, a nontrivial coordinate permutation is
outside the closure of finite products of admissible one-parameter
flows (Theorem~\ref{glob:main}). This identifies a precise limitation
of passing from an infinitesimal classification to static operators.

The coefficient-positive preservation problem and physical quantum
Markov dynamics have their own generator cones and analytic hypotheses.
They are developed in companion manuscripts
\cite{WeiPositive2026,WeiQuantum2026}. The present paper is confined to
finite boxes and the spectral and computational consequences of
preserving the full complex stable class.

\subsection{Organization}
Section~\ref{pre:section} fixes the normalizations.
Section~\ref{gen:section} proves the generator classification and its
local-input recognition corollary. Section~\ref{spec:section} develops
strictification, the spectral gap, dynamical contraction, and sector
concavity. Section~\ref{alg:section} proves the principal-eigenvalue,
state-query, and measurement-sampling algorithms, including the
zero-field value scheme. Sections~\ref{eq:section} and
\ref{spin:section} classify all-temperature Hamiltonians for qubits
and higher spins. The appendices give the exact-arithmetic details,
the obstruction to flow generation, and the certified radius
counterexample.
\par


\section{Conventions and zero-free kernels}\label{pre:section}

\subsection{Regions, degrees, and preservation}
Write $\HH=\{z\in\C:\Im z>0\}$ and $\D_r=\{z\in\C:|z|<r\}$,
with $\D=\D_1$.
\begin{definition}[Stability and preservation]\label{pre:stability}
A nonzero polynomial is \emph{$\Omega$-stable} if it
has no zero when all its variables lie in $\Omega$.
For a product of different regions the corresponding meaning is used.
A polynomial is \emph{real stable} if its coefficients are real and it is
$\HH$-stable. In one variable this is equivalent to real-rootedness.
A \emph{positive stable polynomial} has nonnegative real coefficients and
is real stable. Positivity throughout refers to coefficients, unless an
operator on a Hilbert space is explicitly said to be positive semidefinite.
An operator \emph{preserves} one of these classes when it sends every
member to another member or to zero.
\end{definition}

Let $\N=\{0,1,2,\ldots\}$. For $\kappa\in\N^n$ put
\[
 V_\kappa=\{f\in\C[z_1,\ldots,z_n]:\deg_{z_i}f\leq\kappa_i\},
 \qquad
 \binom{\kappa}{\alpha}=\prod_i\binom{\kappa_i}{\alpha_i}.
\]
Variables of capacity zero can be deleted. A degree bound is not an
assertion that the polynomial has that degree. All finite-box statements
include degree drops and constants. For a real operator, its action on
complex polynomials always means its complex-linear extension.
Every operator on $V_\kappa$ has a unique canonical expression
\begin{equation}\label{pre:normal-order}
 A=\sum_{\alpha\leq\kappa}P_\alpha(z)\partial^\alpha,
\end{equation}
as an operator on $V_\kappa$. Uniqueness follows by applying the expression
successively to monomials, starting with $1$: the equation for $z^\alpha$
determines $\alpha!P_\alpha$ after all lower terms are known. The phrase
\emph{differential order} refers to this expression.

A finite-dimensional generator $A$ preserves a class when $e^{tA}$
preserves it for all $t\geq0$. All generator and spectral statements here act on finite-dimensional
polynomial boxes. Positive dynamics at infinite capacities and their
analytic generation are studied separately in \cite{WeiPositive2026}.

\begin{definition}[Complete stability preservation]\label{pre:complete}
An operator is \emph{completely stable} if its tensor product with the
identity in arbitrarily many extra polynomial variables preserves the
same complex stable class, for every finite choice of the extra degree
bounds.
\end{definition}
This is a statement about polynomial
stability, distinct from complete positivity of a quantum channel.
Positive stability with nonnegative coefficients and complete complex
stability must also be distinguished.

\subsection{Polarization and the disk kernel}
The normalized polarization of $z_i^a$ into $\kappa_i$ binary slots is
\[
 z_i^a\longmapsto
 \binom{\kappa_i}{a}^{-1}
 e_a(z_{i1},\ldots,z_{i\kappa_i}).
\]
It is symmetric in the slots and specializes to $z_i^a$ on the diagonal.
Polarization preserves stability in the convex circular regions used
here, including open disks and half-planes; diagonalization preserves
it directly. We use the Grace--Walsh--Szeg\H{o} coincidence
theorem in this form; see \cite{BB09,BB10II,BB09Circular}.
The binomial normalization is part of the definition, not an optional
choice of basis.\label{pre:polarization}

For $T:V_\kappa\to V_\lambda$ define its disk kernel by
\begin{equation}\label{pre:kernel}
 K_T(z,w)=T_z\prod_i(1+z_iw_i)^{\kappa_i}
 =\sum_{\alpha\leq\kappa}\binom{\kappa}{\alpha}
       (Tz^\alpha)(z)w^\alpha.
\end{equation}
The notation $z_i$ in the first expression denotes the input variables
on which $T$ acts; its result is a polynomial in the output variables.
The kernel has output degree bound $\lambda$ and input degree bound
$\kappa$. In particular $K_{\Id}(z,w)=\prod_i(1+z_iw_i)^{\kappa_i}$.

For two polynomials of degree at most $\kappa$ in paired variables define
their bilinear contraction by
\begin{equation}\label{pre:contraction}
 [f(u),g(v)]_\kappa
   =\sum_{\alpha\leq\kappa}
      \binom{\kappa}{\alpha}^{-1} f_\alpha g_\alpha,
 \qquad
 f(u)=\sum_\alpha f_\alpha u^\alpha,\quad
 g(v)=\sum_\alpha g_\alpha v^\alpha .
\end{equation}
There is no complex conjugation in this formula. Matrix composition is
exactly contraction of the corresponding kernels:
\[
 K_{ST}(z,w)=[K_S(z,u),K_T(v,w)]_\kappa,
\]
where $\kappa$ is the intermediate capacity.

\begin{lemma}[Disk contraction]\label{pre:grace}
Suppose $f(z,u)$ and $g(v,w)$ are stable in products of disks and
$\deg_{u_i}f,\deg_{v_i}g\leq\kappa_i$. If their paired radii
$r_i,s_i$ satisfy $r_is_i\geq1$, then their contraction
\eqref{pre:contraction} is stable in the remaining disks or is identically
zero. If all $r_is_i>1$, contraction at any point in the remaining
disks is nonzero.
\end{lemma}
\begin{proof}
First consider a multiaffine polynomial
$F(u,v)=a+bu+cv+duv$ that is nonzero on $\D_r\times\D_s$.
Then $a\ne0$ and $F(rt,st)$ is nonzero for $|t|<1$.
If $d\ne0$, the product of its two roots has modulus at least one,
so $|d|rs\le|a|$; if $d=0$, this inequality is immediate.
Thus $a+d\ne0$ when $rs>1$.
The same argument applies after fixing any uncontracted variables
in their respective disks. Consequently the operation
$a+bu+cv+duv\mapsto a+d$ preserves stability under strict
separation of the paired radii. It may be iterated on a general
multiaffine polynomial, without a factorization assumption.

Polarize the paired variables of $f$ and $g$ into $\kappa_i$
slots for coordinate $i$, and apply this operation to each pair
of corresponding slots in their product.
For a fixed multi-index $\alpha$, the matching slot subsets have
cardinality $\binom{\kappa}{\alpha}$, while the two normalized
polarizations contribute the factor
$\binom{\kappa}{\alpha}^{-2}$.
The result is therefore exactly \eqref{pre:contraction}.
This proves strict nonvanishing, including degree deficiencies.
It is the weighted form of circular-domain contraction;
see \cite{AsanoJapanese70,BB09,BB10II}.
For equality of paired radii, replace every paired variable in each input
polynomial by $(1-\varepsilon)$ times that variable, apply the strict
statement, and let $\varepsilon\downarrow0$. Hurwitz's theorem gives
the stable-or-zero conclusion. The argument is valid after fixing
any values of the unpaired variables in their respective disks.
\end{proof}

The finite-box symbol theorem \cite{BB09,BB10II} identifies stability
preservers of rank at least two by the appropriate stable kernel;
the low-rank alternatives are retained whenever they arise.
In particular an invertible disk-stability preserver has a stable
kernel \eqref{pre:kernel}. We do not infer a kernel criterion for
arbitrary positive-coefficient preservers from this theorem.

The Cayley isomorphism is
\begin{equation}\label{pre:cayley}
 (\mathcal C_\kappa f)(x)
   =\prod_i(x_i+\mathrm i)^{\kappa_i}
      f\left(\frac{x_1-\mathrm i}{x_1+\mathrm i},\ldots,
              \frac{x_n-\mathrm i}{x_n+\mathrm i}\right).
\end{equation}
It is a linear isomorphism between the degree boxes, and $f$ is
$\D$-stable if and only if $\mathcal C_\kappa f$ is $\HH$-stable.
Thus the generator classifications in half-plane coordinates and
the spectral results in disk coordinates concern conjugate
finite-dimensional semigroups.

\subsection{Coherent coordinates and physical tensors}
Equip $V_\kappa$ with the inner product for which
$\binom{\kappa}{\alpha}^{1/2}z^\alpha$ is an orthonormal basis.
Equivalently, identify the standard basis $|\alpha\rangle$ of
$\bigotimes_i\C^{\kappa_i+1}$ with
$\binom{\kappa}{\alpha}^{1/2}z^\alpha$. The polynomial of a vector
$\psi$ is
\[
 f_\psi(z)=\sum_{\alpha\leq\kappa}
          \binom{\kappa}{\alpha}^{1/2}\psi_\alpha z^\alpha.
\]
An operator $M$ on this Hilbert space has coherent tensor
\begin{equation}\label{pre:coherent}
 F_M(z,w)=\sum_{\alpha,\beta\leq\kappa}
 \binom{\kappa}{\alpha}^{1/2}\binom{\kappa}{\beta}^{1/2}
 M_{\alpha\beta}z^\alpha w^\beta .
\end{equation}
This is precisely \eqref{pre:kernel} for its action in coherent
coordinates. Bra variables are independent complex variables;
they are not conjugated when testing zero-freeness.
For a positive semidefinite density matrix $\rho$, the assertion
that $F_\rho$ is $\D$-stable is called its \emph{Lee--Yang property}.

For a linear map $\Phi$ between matrix algebras the full coherent
Choi tensor is defined by replacing $M_{\alpha\beta}$ in
\eqref{pre:coherent} with the coefficients
$\Phi(|\gamma\rangle\langle\delta|)_{\alpha\beta}$ and adding
input monomials $u^\gamma v^\delta$ with the corresponding square-root
binomial weights. Its stability is called \emph{full-Choi
Lee--Yang preservation}. For qubits all capacities are one, so
these weights are one. They cannot be omitted for higher spin.
Ordinary physical preservation tests only positive semidefinite
Lee--Yang inputs; it is a weaker condition for isolated maps.

We use the Pauli matrices
\[
 X=\begin{pmatrix}0&1\\1&0\end{pmatrix},\qquad
 Y=\begin{pmatrix}0&-\mathrm i\\\mathrm i&0\end{pmatrix},\qquad
 Z=\begin{pmatrix}1&0\\0&-1\end{pmatrix}.
\]
Thus $E_i=z_i\partial_i$ represents $|1\rangle\langle1|_i$
and $Z_i=I-2E_i$ on a qubit.
For general capacity $E_i=z_i\partial_i$ has eigenvalues
$0,\ldots,\kappa_i$. A \emph{principal eigenvalue} below means an
algebraically simple eigenvalue whose real part is strictly greater
than the real parts of all other eigenvalues; for a discrete operator
it means the corresponding strict modulus dominance.


\section{Generators on a finite polynomial box}
\label{gen:section}

We first classify semigroups acting on the whole stable class, with no
assumption on the signs of coefficients of the polynomials or of the
operators.  The distinction from preservation of the coefficient-positive
stable class is essential: the two problems have different generator cones.

Let \(\kappa=(\kappa_1,\ldots,\kappa_d)\) have finite positive integer
coordinates, and write
\[
 V_\kappa^{\mathbb F}
 =\mathbb F[z_1,\ldots,z_d]_{\deg_{z_i}\leq\kappa_i},
 \qquad \mathbb F\in\{\R,\C\}.
\]
Coordinates of capacity zero are omitted.  If no coordinate remains, the
space consists of constants and every scalar generator has the requisite
preservation property.  Stability and complete preservation have the meanings
of Definitions~\ref{pre:stability} and~\ref{pre:complete}; in this section the
domain of stability is \(\HH^d\).
A real operator acts on \(V_\kappa^\C\) by complexification.

For a polynomial \(q\) in one variable of degree at most two, define
\begin{equation}
 J_{i,\kappa_i}(q)
 =q(z_i)\partial_i-\frac{\kappa_i}{2}q'(z_i).
 \label{gen:eq-J}
\end{equation}
These are endomorphisms of the box.  The three choices \(q=1,z_i,z_i^2\)
give
\[
 J_i^-=\partial_i,\qquad
 J_i^0=z_i\partial_i-\frac{\kappa_i}{2},\qquad
 J_i^+=z_i^2\partial_i-\kappa_i z_i.
\]
Their span is the sitewise \(\mathfrak{sl}_2\) algebra.

Every endomorphism has a unique canonical normal representation
\begin{equation}
 A=\sum_{\alpha\leq\kappa}P_\alpha(z)\partial^\alpha.
 \label{gen:eq-canonical}
\end{equation}
Here the coefficient polynomials are not assumed to belong to the original
box.  The order of \(A\) means the largest \(|\alpha|\) for which its
canonical coefficient \(P_\alpha\) is nonzero.  The uniqueness of this
representation, proved below, is important at small capacities: in particular
\(\partial_i^2\) is absent when \(\kappa_i=1\).

\begin{theorem}[Real generators]
\label{gen:main-real}
Let \(A\in\operatorname{End}_{\R}(V_\kappa^\R)\).  The following conditions
are equivalent.
\begin{enumerate}
\item For every \(t\geq0\), \(e^{tA}\) preserves real stability.
\item For every \(t\geq0\), its complexification preserves complex stability,
also after adjoining arbitrary inert polynomial variables.
\item The canonical representation of \(A\) has order at most two, and its
second-order coefficients satisfy
\begin{align}
 P_{2e_i}(s)&\leq0
 &&\text{for all }s\in\R,\quad \kappa_i\geq2,
 \label{gen:eq-real-unary}\\
 P_{e_i+e_j}(s,t)&\leq0
 &&\text{for all }(s,t)\in\R^2,\quad i<j.
 \label{gen:eq-real-pair}
\end{align}
\end{enumerate}
In the third condition, box preservation forces \(P_{2e_i}\) to depend only
on \(z_i\), with degree at most four, and \(P_{e_i+e_j}\) to depend only on
\(z_i,z_j\), with bidegree at most \((2,2)\).  There are no further
inequalities on the zeroth- and first-order coefficients beyond the
identities already forced by \(A(V_\kappa^\R)\subseteq V_\kappa^\R\).
\end{theorem}

The coefficient \(P_{e_i+e_j}\) in this statement is the full coefficient
of \(\partial_i\partial_j\), with \(i<j\).  Thus a convention writing the
second-order part as \(\sum_{i,j}Q_{ij}\partial_i\partial_j\), with \(Q\)
symmetric, has \(P_{e_i+e_j}=2Q_{ij}\).
We use the full-coefficient convention throughout this section.

\begin{theorem}[Complex generators]
\label{gen:main-complex}
Let \(A\in\operatorname{End}_{\C}(V_\kappa^\C)\).  The following conditions
are equivalent.
\begin{enumerate}
\item For every \(t\geq0\), \(e^{tA}\) preserves complex stability.
\item For every \(t\geq0\), \(e^{tA}\) preserves complex stability after
adjoining arbitrary inert polynomial variables.
\item There are a real generator \(R\) satisfying
Theorem~\ref{gen:main-real}, a real scalar \(c\), and real polynomials
\(q_i\) of degree at most two, nonnegative on \(\R\), such that
\begin{equation}
 A=R+i\left(c\Id+\sum_{i=1}^dJ_{i,\kappa_i}(q_i)\right).
 \label{gen:eq-complex-decomposition}
\end{equation}
\end{enumerate}
Equivalently, the canonical representation has order at most two, all its
second-order coefficients are real and satisfy
\eqref{gen:eq-real-unary}--\eqref{gen:eq-real-pair}, and every
\(\Im P_{e_i}\) is nonnegative on \(\R^d\).  In this formulation box
preservation forces
\(\Im P_{e_i}(z)=q_i(z_i)\) with \(\deg q_i\leq2\), and forces the
corresponding imaginary zeroth-order term in
\eqref{gen:eq-complex-decomposition}.
\end{theorem}

The two theorems identify ordinary and complete preservation for
\emph{semigroups starting at the identity}.  We will prove this by constructing
the generators.  It is not an assertion that an arbitrary isolated real
stability preserver preserves complex stability or admits every tensor
extension.

\subsection{Canonical differential operators and box preservation}
\label{gen:subsec-box}

We record the algebraic restrictions separately from the stability
inequalities.  All statements in this subsection hold over either \(\R\)
or \(\C\).

\begin{lemma}[Canonical normal representation]
\label{gen:canonical}
Every linear endomorphism of \(V_\kappa^\mathbb F\) has a unique
representation \eqref{gen:eq-canonical}.
\end{lemma}

\begin{proof}
Apply the proposed representation successively to monomials in increasing
total degree.  The constant monomial determines \(P_0=A1\).
Once \(P_\beta\) is known for all \(|\beta|<|\alpha|\), the equation on
\(z^\alpha\) determines
\[
 P_\alpha
 =\frac1{\alpha!}
 \left(Az^\alpha-
 \sum_{\substack{\beta\leq\alpha\\\beta\neq\alpha}}
 P_\beta(z)(\alpha)_\beta z^{\alpha-\beta}\right).
\]
Here \((\alpha)_\beta=\prod_i\alpha_i!/(\alpha_i-\beta_i)!\).
No derivative whose index is not bounded by \(\alpha\) acts nontrivially on
\(z^\alpha\).  The recursion therefore proves existence and uniqueness
on the monomial basis.
\end{proof}

Descriptions of differential operators on finite polynomial modules by
$\mathfrak{sl}_2$ enveloping algebras are classical in quasi-exact
solvability \cite{Turbiner92,Turbiner94}; we include the following
decomposition to make the box conventions and small-capacity degeneracies
explicit.

\begin{proposition}[Operators of order at most two on a box]
\label{gen:box-decomposition}
Every box endomorphism of canonical order at most two is a linear
combination of \(\Id\), the operators \(J_i^-,J_i^0,J_i^+\), and their
products of length two, including products at the same site.
Every endomorphism of order at most one has a unique expression
\begin{equation}
 c\Id+\sum_{i=1}^dJ_{i,\kappa_i}(q_i),
 \qquad c\in\mathbb F,\quad q_i\in\mathbb F[z_i]_{\leq2}.
 \label{gen:eq-first-order}
\end{equation}
Consequently the second-order coefficients have the variable dependence
and degree bounds stated in Theorem~\ref{gen:main-real}.  The vector space
of endomorphisms of order at most two has dimension
\begin{equation}
 1+3d+5\#\{i:\kappa_i\geq2\}+9\binom d2.
 \label{gen:eq-order-two-dimension}
\end{equation}
\end{proposition}

\begin{proof}
Write \(p_{\alpha,\beta}\) for the coefficient of \(z^\beta\) in
\(P_\alpha\), and group terms according to the integer monomial shift
\(\nu=\beta-\alpha\).  For input exponent
\(x\in\prod_i\{0,\ldots,\kappa_i\}\), the coefficient of \(z^{x+\nu}\) is
the restriction to this grid of
\begin{equation}
 F_\nu(x)=
 \sum_{\substack{|\alpha|\leq2,\ \alpha\leq\kappa\\
                  \alpha+\nu\geq0}}
 p_{\alpha,\alpha+\nu}(x)_\alpha.
 \label{gen:eq-shift-polynomial}
\end{equation}
This polynomial has total degree at most two and degree at most
\(\kappa_i\) in each \(x_i\).  We use the polynomial
\eqref{gen:eq-shift-polynomial} also when \(x+\nu\) is outside the exponent
grid.

Suppose first that \(\nu_i<0\).  Every summand has
\(\alpha_i\geq-\nu_i\), so \(F_\nu\) is divisible by
\((x_i)_{-\nu_i}\).  If \(-\nu_i>\kappa_i\), there are no summands and
\(F_\nu=0\).  In particular, whenever a grid point would give a negative
output exponent, this falling factorial makes \(F_\nu\) vanish there.

If \(\nu_i>0\), box preservation gives vanishing on every overflow face
\[
 x_i\in\{\kappa_i-\nu_i+1,\ldots,\kappa_i\}
          \cap\{0,\ldots,\kappa_i\}.
\]
To justify the divisibility as a polynomial identity, fix one such value
of \(x_i\).  On the remaining tensor grid the value of \(F_\nu\) is zero:
either the output has a negative coordinate, in which case the preceding
falling-factorial argument applies, or it is a genuine monomial outside
the box, whose coefficient must vanish.  Tensor-product interpolation,
using the coordinate degree bounds in
\eqref{gen:eq-shift-polynomial}, shows that the restriction to this face
is identically zero as a polynomial in all remaining variables.  Thus
the corresponding linear factor in \(x_i\) divides \(F_\nu\).
If \(\nu_i>\kappa_i\), the same argument gives vanishing on all
\(\kappa_i+1\) coordinate values and hence \(F_\nu=0\).

For every remaining nonzero shift polynomial, these factors give
\begin{equation}
 F_\nu=D_\nu G_\nu,\qquad
 D_\nu(x)=\prod_i(x_i)_{\nu_i^-}
                       (\kappa_i-x_i)_{\nu_i^+},
 \qquad
 \deg G_\nu\leq2-|\nu|_1.
 \label{gen:eq-shift-factor}
\end{equation}
The factors belonging to distinct coordinate values are coprime.
Thus \(F_\nu\neq0\) implies \(|\nu|_1\leq2\).

Put \(E_i=z_i\partial_i\).  For \(\nu=0\), every possible \(F_\nu\)
is realized by a combination of
\(\Id,E_i,E_i^2,E_iE_j\).
For \(\nu=\pm e_i\), the required constant or affine factor \(G_\nu\)
is realized by \(J_i^\pm\) and \(J_i^\pm E_j\).
For \(\nu=\pm2e_i\) use \((J_i^\pm)^2\); these shifts can be nonzero
only if \(\kappa_i\geq2\).
For \(\nu=\pm e_i\pm e_j\), with \(i\neq j\), use the corresponding
products \(J_i^\pm J_j^\pm\).
The sign of the raising coefficient
\(x_i-\kappa_i=-(\kappa_i-x_i)\) is absorbed in the scalar multiplier.
This accounts for every possibility in \eqref{gen:eq-shift-factor}
on the input grid, and hence for every matrix entry of \(A\).

This argument also covers binary sites.  For example, if
\(\kappa_i=1\), the canonical coordinate degree of \(F_\nu\) in
\(x_i\) is at most one; any inadmissible quadratic dependence has
already been excluded before the realization step.
Some displayed products may then be linearly dependent or zero, which
does not affect the spanning assertion.

Repeating the same argument with total degree at most one leaves only
a scalar and the three sitewise operators at each coordinate.
They are linearly independent: the off-diagonal shifts separate the
raising and lowering terms, while an affine polynomial in the diagonal
exponents vanishing on the full tensor grid has every coefficient zero.
This proves \eqref{gen:eq-first-order}, including uniqueness.

The first-order coefficients of \(J_i^-,J_i^0,J_i^+\) span
\(1,z_i,z_i^2\).  Products at one site therefore have principal
coefficients spanning all polynomials of degree at most four.
Products at two different sites span all biforms of bidegree at most
\((2,2)\).  At a site with \(\kappa_i\geq2\), these are canonical
second-order coefficients; at a binary site that derivative is absent.
The kernel of this principal-coefficient map is exactly the
order-at-most-one space.  Its dimension is \(1+3d\), and the respective
principal-coefficient spaces have dimensions five and nine.
Surjectivity gives \eqref{gen:eq-order-two-dimension}.
\end{proof}

The proposition says that the operators of canonical order at most two
form the image of the degree-at-most-two part of the enveloping algebra
of the sitewise \(\mathfrak{sl}_2\)'s.  This is an algebraic statement
about the whole box.  Stability will impose precisely the signs of
its principal coefficients and, in the complex case, an inward sign
on the imaginary first-order fields.

\subsection{Root clusters and infinitesimal necessity}
\label{gen:subsec-roots}

The relevant localization uses an isolated root cluster and is valid
even when the ambient polynomial degree changes.  No full-degree
tangent formula is applied to a lower-degree polynomial.
The underlying multiple-root contact principle is classical:
after rotation of the half-plane, it is contained in the polynomial
abscissa tangent theorem of Burke and Overton
\cite[Theorem~1.2 and Lemma~1.4]{BurkeOverton01}.
We give a direct moment proof to record the signs and the local
factorization needed for the operator argument.

\begin{lemma}[Lower-half-plane root clusters]
\label{gen:cluster-moments}
Let \(p_t(s)\), \(t\geq0\), be a curve of complex polynomials of uniformly
bounded degree whose coefficients are right differentiable at zero.
Suppose
\[
 p_0(s)=s^m g(s),\qquad m\geq1,\qquad g(0)\neq0,
\]
and that \(p_t\) has no zero in \(\HH\) for all sufficiently small \(t\).
Let \(C_t\) be the monic factor consisting of the \(m\) roots converging
to zero.  Its coefficients are right differentiable, and
\begin{equation}
 \left.\frac{d}{dt}C_t(s)\right|_{t=0}
 =c_1s^{m-1}+c_2s^{m-2},
 \qquad
 \Im c_1\geq0,\quad c_2\in\R,\quad c_2\leq0.
 \label{gen:eq-cluster-variation}
\end{equation}
When \(m=1\), the \(c_2\) term is absent.  If every \(p_t\) is real
rooted, then \(c_1\) is real as well.
\end{lemma}

\begin{proof}
Choose a small circle about zero containing no other zero of \(p_0\).
Coefficient continuity and Rouch\'e's theorem isolate exactly \(m\)
roots inside it for small \(t\).  If these roots are
\(\lambda_1(t),\ldots,\lambda_m(t)\), their power sums are
\[
 M_k(t)=\sum_{j=1}^m\lambda_j(t)^k
       =\frac1{2\pi i}\int s^k\frac{p_t'(s)}{p_t(s)}\,ds,
       \qquad k\geq1,
\]
where the contour is the chosen circle.  The integrands are right
differentiable with a denominator bounded away from zero.
Thus each \(M_k\) is right differentiable and \(M_k(t)=O(t)\).
Newton's identities give the same differentiability for the coefficients
of \(C_t\).

Write \(\lambda_j=a_j+ib_j\).  The roots lie in the closed lower
half-plane, so \(b_j\leq0\), and
\[
 \sum_j|b_j|=-\Im M_1(t)=O(t),\qquad
 \sum_jb_j^2=O(t^2).
\]
Moreover,
\[
 \sum_ja_j^2=\Re M_2(t)+\sum_jb_j^2=O(t).
\]
Consequently
\[
 \Im M_2(t)=2\sum_ja_jb_j=O(t^{3/2})=o(t),
\]
and, for \(k\geq3\),
\[
 |M_k(t)|
 \leq\sum_j|\lambda_j(t)|^k
 \leq\left(\sum_j|\lambda_j(t)|^2\right)^{k/2}
 =O(t^{k/2})=o(t).
\]
The second inequality uses the ordinary finite counting sum over
the roots.

Write
\(C_t=s^m+\gamma_1(t)s^{m-1}+\cdots+\gamma_m(t)\).
At \(t=0\), all \(\gamma_j\) and \(M_j\) vanish.
Differentiating Newton's identities gives
\(\gamma_j'(0)=-M_j'(0)/j\).
Hence every derivative below the coefficient of \(s^{m-2}\)
vanishes.  The first coefficient has nonnegative imaginary part
because \(\Im M_1(t)\leq0\).
For the second coefficient, \(M_2'(0)\) is real, and
\[
 \Re M_2'(0)=\lim_{t\downarrow0}
       \frac{\sum_ja_j(t)^2-\sum_jb_j(t)^2}{t}\geq0.
\]
This proves \eqref{gen:eq-cluster-variation}.  Real roots make
\(M_1\) real.
\end{proof}

\begin{proposition}[Order and sign from product contacts]
\label{gen:contact-necessity}
Suppose \(A\in\operatorname{End}_{\C}(V_\kappa^\C)\) generates a
complex stability-preserving semigroup.  Then \(A\) has canonical
order at most two, every second-order canonical coefficient is real
and nonpositive on \(\R^d\), and
\(\Im P_{e_i}\geq0\) on \(\R^d\).
If \(A\) is real and its semigroup is only assumed to preserve real
stability, the order bound and the nonpositivity of its
second-order coefficients still follow.
\end{proposition}

\begin{proof}
For \(0\neq\alpha\leq\kappa\) and \(y\in\R^d\), use the stable
product
\[
 f_{\alpha,y}(z)=\prod_i(z_i-y_i)^{\alpha_i}.
\]
In the canonical representation, every derivative at \(z=y\)
vanishes except for \(\partial^\alpha\), and therefore
\begin{equation}
 (Af_{\alpha,y})(y)=\alpha!P_\alpha(y).
 \label{gen:eq-contact-isolation}
\end{equation}
Choose \(v\in\R_{>0}^d\) and restrict the polynomial curve to
\(z=y+sv\).  Its initial value is
\[
 \left(\prod_iv_i^{\alpha_i}\right)s^{|\alpha|}.
\]
Its coefficients are differentiable and its degree is bounded by
\(|\kappa|\).  It has no zero for \(s\in\HH\), because such an \(s\)
puts every active coordinate in \(\HH\).
The restricted polynomial cannot be identically zero, since that
would contradict multivariate stability.  Invertibility of
\(e^{tA}\) also excludes the zero multivariate output.

Apply Lemma~\ref{gen:cluster-moments}.  Factoring the restriction
as \(C_t(s)G_t(s)\), its constant first variation is
\[
 G_0(0)\left.\frac{d}{dt}C_t(0)\right|_{t=0}.
\]
The term involving the derivative of \(G_t\) vanishes because \(C_0(0)=0\).
Here \(G_0(0)=\prod_iv_i^{\alpha_i}>0\).
For \(|\alpha|\geq3\), the constant variation is zero, so
\eqref{gen:eq-contact-isolation} gives \(P_\alpha(y)=0\).
For \(|\alpha|=2\), it is real and nonpositive, giving the
corresponding assertion for \(P_\alpha(y)\).
For \(\alpha=e_i\), its imaginary part is nonnegative.

Since \(y\) is arbitrary, polynomial identity proves the vanishing
of all coefficients of order greater than two.  A complex polynomial
real-valued on all of \(\R^d\) has real coefficients; this proves
reality of the principal coefficients.  For a real semigroup,
the same inputs and restrictions are real rooted, so the real
case follows from the same lemma.  Proposition~\ref{gen:box-decomposition}
then supplies the asserted variable dependence.
\end{proof}

The moment calculation identifies the obstruction to higher differential
order: one-sided control of the imaginary root depths, together with
the second root moment, eliminates every higher first-order cluster
moment.  Positivity of coefficients of the input polynomial is not
used; every real contact point is available.

\subsection{Binary apolar generators and sums of squares}
\label{gen:subsec-binary}

The sufficiency proof starts in two binary variables.
For
\(f=f_{00}+f_{10}s+f_{01}t+f_{11}st\in V_{(1,1)}^\R\), put
\[
 \Delta(f)=f_{10}f_{01}-f_{00}f_{11}.
\]
A nonzero real \(f\) is stable if and only if \(\Delta(f)\geq0\).
Indeed, if \(f=a+bs+ct+dst\) and \(b+dt\neq0\), its zero in \(s\)
is \(-(a+ct)/(b+dt)\), whose imaginary part equals
\(-\Delta(f)\Im t/|b+dt|^2\).
The cases in which \(f\) is independent of \(s\) follow directly.

\begin{lemma}[Low-degree nonnegative polynomials]
\label{gen:sos}
Every nonnegative real univariate polynomial of degree at most four is a sum
of squares of real polynomials of degree at most two.
Every nonnegative real polynomial of bidegree at most \((2,2)\)
is a sum of squares of real bilinear polynomials.
\end{lemma}

\begin{proof}
For a nonzero nonnegative univariate polynomial, every real root has
even multiplicity and its nonreal roots occur in conjugate pairs.
A positive real quadratic is a sum of two squares of linear
polynomials.  Repeated use of
\[
 (a^2+b^2)(c^2+d^2)=(ac-bd)^2+(ad+bc)^2
\]
gives the first assertion, with square degrees at most two.
The zero polynomial presents no exception.

For the second assertion let \(R(s,t)\geq0\) and homogenize it to
\[
 \widetilde R(x,y,u)=u^4R(x/u,y/u).
\]
The bidegree bound makes this a ternary quartic.  It is nonnegative
when \(u\neq0\), and hence everywhere by continuity.
Hilbert's theorem on nonnegative ternary quartics writes it as a sum
of squares of real quadratic forms; we use the version including
all degenerate nonnegative forms \cite{PRSS04}.
The coefficients of \(x^4\) and \(y^4\) in \(\widetilde R\) are zero.
They are sums of the squares of the \(x^2\) and \(y^2\) coefficients
of the quadratic forms, respectively.  Thus every form is a linear
combination of \(xy,xu,yu,u^2\).  Setting \(u=1\) yields bilinear
squares.
\end{proof}

\begin{lemma}[Apolar rank-one primitive]
\label{gen:binary-primitive}
For a real bilinear polynomial \(F=a+bs+ct+dst\), define
\[
 \ell_F(f)=d f_{00}-c f_{10}-b f_{01}+a f_{11},
 \qquad
 B_Ff=-F\ell_F(f).
\]
The full mixed canonical coefficient of \(B_F\) is \(-F^2\).
Moreover
\begin{align}
 B_F^2&=2\Delta(F)B_F,
 \label{gen:eq-apolar-square}\\
 \Delta(f+kB_Ff)
 &=\Delta(f)+k\bigl(1+k\Delta(F)\bigr)\ell_F(f)^2
 \qquad(k\in\R).
 \label{gen:eq-apolar-discriminant}
\end{align}
For all \(t\geq0\), \(e^{tB_F}\) preserves real stability.
\end{lemma}

\begin{proof}
For any endomorphism \(B\) of the binary box, its mixed coefficient is
\begin{equation}
 P_B(s,t)=B(st)-sB(t)-tB(s)+stB(1).
 \label{gen:eq-binary-coefficient}
\end{equation}
The four evaluations of \(B_F\) give \(P_{B_F}=-F^2\).
Also
\(\ell_F(F)=2ad-2bc=-2\Delta(F)\), which gives
\eqref{gen:eq-apolar-square}.
The polarization of the quadratic form \(\Delta\) satisfies
\(D\Delta_f[F]=-\ell_F(f)\).  Expanding
\(\Delta(f-kF\ell_F(f))\) proves
\eqref{gen:eq-apolar-discriminant}.

If \(D_F=\Delta(F)\), equation~\eqref{gen:eq-apolar-square} gives
\[
 e^{tB_F}=\Id+k_tB_F,\qquad
 k_t=
 \begin{cases}
 (e^{2tD_F}-1)/(2D_F),&D_F\neq0,\\
 t,&D_F=0.
 \end{cases}
\]
For \(t\geq0\), \(k_t\geq0\) and
\[
 1+k_tD_F=\frac{1+e^{2tD_F}}2>0.
\]
Thus \eqref{gen:eq-apolar-discriminant} preserves the inequality
\(\Delta(f)\geq0\).
The exponential is invertible, so a nonzero input does not become
the zero polynomial.  These formulas also include \(F=0\) and
\(D_F=0\).
\end{proof}

For use below, we make explicit the passage from real to complex
preservation in this binary semigroup.

\begin{lemma}[Orientation at the identity]
\label{gen:binary-completeness}
Let \(B\) be a real endomorphism of \(V_{(1,1)}^\R\).
Then \(e^{tB}\) preserves real stability for all \(t\geq0\) if and
only if its mixed coefficient \(P_B\) is nonpositive on \(\R^2\).
In this case every \(e^{tB}\) is a complete complex stability
preserver.
\end{lemma}

\begin{proof}
For necessity, take \(f=(s-x)(t-y)\).  It has \(\Delta(f)=0\).
Differentiating \(\Delta(e^{\tau B}f)\) at \(\tau=0\), and using
\eqref{gen:eq-binary-coefficient}, gives \(-P_B(x,y)\).
Real preservation implies that this derivative is nonnegative.

Conversely, Lemma~\ref{gen:sos} gives
\(-P_B=\sum_\ell F_\ell^2\).
The operator \(B-\sum_\ell B_{F_\ell}\) has zero mixed coefficient
and hence has canonical order at most one.
Proposition~\ref{gen:box-decomposition} writes it as a real scalar
plus two sitewise operators \(J_{i,1}(q_i)\), with real quadratic
\(q_i\), without a sign restriction on these two quadratics.

For clarity, the real sitewise flows preserve complex stability in
both time directions.  If \(q(z)=q_0+q_1z+q_2z^2\), let
\[
 M_\tau=\exp\left[
 \tau\begin{pmatrix}q_1/2&q_0\\-q_2&-q_1/2\end{pmatrix}
 \right]
 =\begin{pmatrix}a_\tau&b_\tau\\c_\tau&d_\tau\end{pmatrix}.
\]
This is a real matrix of determinant one.  The polynomial action
\begin{equation}
 U_\tau f(z)
 =(c_\tau z+d_\tau)^\kappa
 f\left(\frac{a_\tau z+b_\tau}{c_\tau z+d_\tau}\right)
 \label{gen:eq-real-mobius}
\end{equation}
has infinitesimal generator \(J_\kappa(q)\).
The quotient maps \(\HH\) to itself, since its imaginary part is
\(\Im z/|c_\tau z+d_\tau|^2\), and the denominator has no zero there.
Formula~\eqref{gen:eq-real-mobius} is polynomial after cancellation
term by term and holds for every real \(\tau\).
It also works with inert variables.

Lemma~\ref{gen:binary-primitive} and the finite-dimensional Lie
product formula now show that \(e^{\tau B}\) preserves real
stability.  Products of preserving factors converge coefficientwise;
Hurwitz's theorem and invertibility of the limiting exponential
exclude a zero output.

It remains to justify complex preservation of the apolar factors
and of the resulting binary semigroup.  We use the finite-box
classification of Borcea and Br\"and\'en
\cite[Theorems~1.1 and~1.2]{BB09}.
A real stability preserver of rank greater than two has a stable
plus symbol or a stable reflected-minus symbol.  For the binary
box these are
\[
 G_T^+(s,t;u,v)=T\bigl((s+u)(t+v)\bigr),\qquad
 G_T^-(s,t;u,v)=T\bigl((s-u)(t-v)\bigr),
\]
where \(T\) acts only in \(s,t\).
Every \(e^{\tau B}\) has rank four.
At \(\tau=0\) the minus symbol is \((s-u)(t-v)\), which is not
stable: it vanishes when \(s=u\in\HH\).
The cone consisting of stable polynomials and zero is closed
in this fixed coefficient space.  Continuity in \(\tau\) therefore
excludes the minus alternative for all sufficiently small
nonnegative \(\tau\).  The plus symbol is then stable, and the
complex classification gives complex preservation.

More explicitly, adjoining an identity operator on an auxiliary
box of capacity \(\lambda\) multiplies the plus symbol by
\(\prod_j(\xi_j+\eta_j)^{\lambda_j}\).  This remains stable.
The same complex classification proves preservation on every
enlarged box.  Finally, divide an arbitrary time into sufficiently
small equal pieces and compose.  This proves complete preservation
for every time and completes the lemma.
\end{proof}

This use of the symbol theorem is limited and precise.
The real rank-two exceptional case is not invoked: the binary
pair operator has rank four at every time.
A single binary site is handled directly by
\eqref{gen:eq-real-mobius}.

\subsection{Polarization and the proof of the real classification}
\label{gen:subsec-real-proof}

Let \(\mathcal P_\kappa\) be normalized polarization into
\(N=|\kappa|\) binary variables, grouped into \(\kappa_i\) slots at
site \(i\):
\[
 \mathcal P_\kappa(z_i^r)
 =\frac{e_r(z_{i,1},\ldots,z_{i,\kappa_i})}
             {\binom{\kappa_i}{r}}.
\]
Let \(\mathcal D_\kappa\) identify all slots belonging to each site.
The two maps are inverse between \(V_\kappa^\C\) and the subspace
symmetric within each group of slots.
Normalized polarization preserves upper-half-plane stability, and
diagonalization preserves it in the reverse direction
\cite[Proposition~2.4]{BB09}, with the normalization fixed in the preliminaries.

We will need the exact derivative factors
\begin{align}
 \mathcal D_\kappa\partial_{i,a}\mathcal P_\kappa f
 &=\frac{\partial_i f}{\kappa_i},
 \label{gen:eq-polar-one}\\
 \mathcal D_\kappa\partial_{i,a}\partial_{j,b}
                          \mathcal P_\kappa f
 &=\frac{\partial_i\partial_j f}{\kappa_i\kappa_j},
 &&i\neq j,
 \label{gen:eq-polar-cross}\\
 \mathcal D_\kappa\partial_{i,a}\partial_{i,b}
                          \mathcal P_\kappa f
 &=\frac{\partial_i^2f}{\kappa_i(\kappa_i-1)},
 &&a\neq b.
 \label{gen:eq-polar-same}
\end{align}
The last formula is used only for \(\kappa_i\geq2\).
They follow by differentiating the elementary symmetric polynomial
on a monomial and then diagonalizing.

\begin{proof}[Proof of Theorem~\ref{gen:main-real}]
The implication from complete complex preservation to real
preservation is immediate, since the operator is real.
Necessity of the order and sign conditions is
Proposition~\ref{gen:contact-necessity}.
We prove sufficiency constructively.

Write \(a_i=P_{2e_i}\) and \(b_{ij}=P_{e_i+e_j}\).
For each \(i<j\), apply Lemma~\ref{gen:sos} to obtain
\[
 -b_{ij}(s,t)=\sum_\ell F_{ij,\ell}(s,t)^2
\]
with real bilinear \(F_{ij,\ell}\).
The binary operator
\(B_{ij}=\sum_\ell B_{F_{ij,\ell}}\) has mixed coefficient
\(b_{ij}\) and generates a complete preserving semigroup.
On the full slot space, place a copy of \(B_{ij}\) on each pair
consisting of one slot at site \(i\) and one at site \(j\), and
sum these \(\kappa_i\kappa_j\) copies.
Every summand is a complete generator, so the sum is one by the
Lie product formula.
It commutes with permutations of slots at either site and therefore
preserves the symmetric subspace.

Descend this operator by \(\mathcal D_\kappa\) and
\(\mathcal P_\kappa\).
Equation~\eqref{gen:eq-polar-cross} shows that its full mixed
coefficient is exactly \(b_{ij}(z_i,z_j)\).
No other second-order coefficient is created.
Indeed, after writing each binary summand in normal order,
zeroth- and first-order terms remain of those respective orders
after diagonalization, and the only second-order derivative
involves its two selected slots.

For a site with \(\kappa_i\geq2\), write
\[
 -a_i(z)=\sum_\ell p_{i,\ell}(z)^2,\qquad
 p_{i,\ell}(z)=a_\ell z^2+b_\ell z+c_\ell,
\]
using the univariate part of Lemma~\ref{gen:sos}.
Set
\[
 F_{i,\ell}(s,t)=a_\ell st+\frac{b_\ell}{2}(s+t)+c_\ell.
\]
These forms are symmetric in their two variables and satisfy
\(F_{i,\ell}(z,z)=p_{i,\ell}(z)\).
The symmetric binary operator
\(2\sum_\ell B_{F_{i,\ell}}\) has mixed coefficient
\[
 P_i(s,t)=-2\sum_\ell F_{i,\ell}(s,t)^2,
 \qquad P_i(z,z)=2a_i(z).
\]
Symmetry of the operator follows directly from symmetry of both
\(F_{i,\ell}\) and its apolar functional.
Sum copies of this operator over all unordered pairs of distinct
slots at site \(i\).
It is again a complete generator and commutes with all slot
permutations at that site.
Its descended same-site principal coefficient, by
\eqref{gen:eq-polar-same}, is
\[
 \binom{\kappa_i}{2}
 \frac{P_i(z_i,z_i)}{\kappa_i(\kappa_i-1)}
 =a_i(z_i).
\]
This is the reason for the factor two in the binary operator.
There is no same-site construction when \(\kappa_i=1\).

Let \(A_{\mathrm{pr}}\) be the sum of all descended generators.
For each summand, permutation invariance implies that its
exponential on the symmetric subspace is intertwined by
polarization with the exponential of the descended operator.
Polarization, complete preservation on the full slot space,
and diagonalization therefore show that \(A_{\mathrm{pr}}\)
generates a complete preserving semigroup on the original box.
It has exactly the same second-order canonical coefficients as
\(A\).

The difference \(A-A_{\mathrm{pr}}\) has order at most one.
By Proposition~\ref{gen:box-decomposition}, it is a real scalar
plus a sum of real sitewise \(J_{i,\kappa_i}(q_i)\).
Formula~\eqref{gen:eq-real-mobius} proves complete preservation
for each of their flows, with no sign restriction on \(q_i\).
A final Lie product formula proves complete preservation for
\(e^{tA}\).
The argument takes place on any fixed enlarged finite box,
and the limiting exponential is invertible there, so Hurwitz's
zero alternative never occurs for a nonzero stable input.
This proves the theorem.
\end{proof}

\subsection{Imaginary inward fields and the complex classification}
\label{gen:subsec-complex-proof}

The remaining complex generators have a particularly explicit
one-site realization.

\begin{lemma}[Imaginary parabolic primitive]
\label{gen:imaginary-primitive}
Let \(a,b\in\R\) and \(\kappa\geq1\).
The generator \(iJ_\kappa((az+b)^2)\) has the semigroup
\begin{equation}
 (U_tf)(z)=d_t(z)^\kappa f(\phi_t(z)),\qquad
 d_t(z)=1-iat(az+b),\qquad
 \phi_t(z)=\frac{z+ibt(az+b)}{1-iat(az+b)}.
 \label{gen:eq-imaginary-flow}
\end{equation}
For every \(t\geq0\), this is a complete complex stability preserver.
\end{lemma}

\begin{proof}
The corresponding fractional-linear matrix is
\[
 M_t=\Id+it
 \begin{pmatrix}ab&b^2\\-a^2&-ab\end{pmatrix}.
\]
The displayed nonidentity matrix has square zero, so
\(M_{t+s}=M_tM_s\) and \(\det M_t=1\).
Denominators in \eqref{gen:eq-imaginary-flow} cancel on each
monomial of degree at most \(\kappa\), giving a polynomial
operator and an exact semigroup.
Differentiating at zero gives
\[
 \left.\frac{d}{dt}U_tf\right|_{t=0}
 =i\bigl((az+b)^2f'-\kappa a(az+b)f\bigr)
 =iJ_\kappa((az+b)^2)f.
\]

For \(z\in\HH\) and \(t\geq0\), direct multiplication by the
conjugate denominator gives
\begin{equation}
 \Im\phi_t(z)
 =\frac{\Im z+t|az+b|^2}{|1-iat(az+b)|^2}>0.
 \label{gen:eq-imaginary-height}
\end{equation}
When \(a\neq0\) and \(t>0\), the denominator has its only zero at
\(-b/a-i/(a^2t)\), below the real axis.
When \(a=0\), it is identically one and
\(\phi_t(z)=z+ib^2t\).
Thus the substitution maps \(\HH\) to itself and the prefactor
does not vanish there.
The same argument with all other variables held inert proves
complete preservation.  The cases \(t=0\) and \(a=b=0\) give
the identity.
\end{proof}

\begin{proof}[Proof of Theorem~\ref{gen:main-complex}]
Assume ordinary complex preservation and decompose \(A=R+iS\)
by entrywise real and imaginary parts in the ordinary monomial
basis.  Both \(R\) and \(S\) are real box endomorphisms.
The canonical recursion in Lemma~\ref{gen:canonical} is real
linear, so it commutes with this decomposition.

By Proposition~\ref{gen:contact-necessity}, \(A\) has order
at most two and its second-order coefficients are real and
nonpositive.  The same coefficients are therefore the
principal coefficients of \(R\), which satisfies
Theorem~\ref{gen:main-real}.  Meanwhile \(S\) has order at most
one.  Proposition~\ref{gen:box-decomposition} gives
\[
 S=c\Id+\sum_iJ_{i,\kappa_i}(q_i),
 \qquad c\in\R,\quad q_i\in\R[z_i]_{\leq2}.
\]
Its first-order coefficients are exactly
\(q_i=\Im P_{e_i}\); the product-contact test makes them
nonnegative on \(\R\).  This proves the necessity of
\eqref{gen:eq-complex-decomposition}.

Conversely, every real quadratic nonnegative on \(\R\) is a
sum of squares of real linear polynomials.
For \(q(z)=az^2+bz+c\) with \(a>0\), complete the square:
\[
 q(z)=\left(\sqrt a\,z+\frac b{2\sqrt a}\right)^2
            +c-\frac{b^2}{4a}.
\]
The last constant is nonnegative.
If \(a=0\), nonnegativity forces \(b=0\) and \(c\geq0\).
Lemma~\ref{gen:imaginary-primitive} and the Lie product formula
therefore give complete preservation for each \(iJ(q_i)\).
The generator \(R\) has the same property by the real theorem,
and \(ic\Id\) contributes a nonzero scalar phase.
One further product formula proves complete preservation for
their sum.  Invertibility of the finite-dimensional exponential
again excludes the zero limit.
Complete preservation implies ordinary preservation, finishing
the equivalence.
\end{proof}

In particular, every complex second-order coefficient in an
admissible generator is forced to be real.  The imaginary part
has no pair term: it consists entirely of sitewise inward
quadratic fields and a scalar.
The sign in \eqref{gen:eq-imaginary-flow} is fixed by the
upper-half-plane convention: \(i\partial\) shifts the argument
upward and moves polynomial zeros downward.

\subsection{Finite certificates, quotient cones, and exact lifts}
\label{gen:subsec-certificates}

The construction records both a finite membership certificate
and a representation by local binary generators.

\begin{corollary}[Gram certificates and lineality]
\label{gen:gram-cones}
In Theorem~\ref{gen:main-real}, the principal inequalities are
equivalent to the existence of real positive semidefinite matrices
\(G_i\) of size three and \(G_{ij}\) of size four such that
\begin{align}
 -P_{2e_i}(s)
  &=(1,s,s^2)G_i(1,s,s^2)^{\top},
  &&\kappa_i\geq2,
 \label{gen:eq-gram-unary}\\
 -P_{e_i+e_j}(s,t)
  &=(1,s,t,st)G_{ij}(1,s,t,st)^{\top},
  &&i<j.
 \label{gen:eq-gram-pair}
\end{align}
The extra inequalities in Theorem~\ref{gen:main-complex} have
size-two certificates
\begin{equation}
 q_i(s)=(1,s)K_i(1,s)^{\top},\qquad K_i\succeq0.
 \label{gen:eq-gram-imaginary}
\end{equation}
These matrix sizes do not depend on the numerical capacities.

Let \(\mathcal N_4,\mathcal N_{2,2},\mathcal N_2\) denote,
respectively, the cones of nonnegative real univariate quartics
of degree at most four, nonnegative real biforms of bidegree
at most \((2,2)\), and nonnegative real univariate quadratics.
The real generator cone has lineality
\[
 \mathfrak g_\kappa
 =\R\Id+\sum_iJ_{i,\kappa_i}(\R[z_i]_{\leq2}),
\]
and its quotient by this lineality is naturally isomorphic to
\begin{equation}
 \prod_{i:\kappa_i\geq2}\mathcal N_4
 \ \times\!
 \prod_{i<j}\mathcal N_{2,2}.
 \label{gen:eq-real-quotient}
\end{equation}
The complex generator cone, regarded as a real cone, has
lineality
\[
 \C\Id+\sum_iJ_{i,\kappa_i}(\R[z_i]_{\leq2}),
\]
and its quotient is the product in
\eqref{gen:eq-real-quotient} times \(\prod_i\mathcal N_2\).
\end{corollary}

\begin{proof}
The square decompositions in Lemma~\ref{gen:sos} give
\eqref{gen:eq-gram-unary} and \eqref{gen:eq-gram-pair} by
taking sums of outer products of the coefficient vectors.
Conversely a real positive semidefinite matrix is a sum
of rank-one positive semidefinite matrices, yielding the
required square decomposition.
The quadratic case is the same argument with the vector
\((1,s)\).

The negative-principal-coefficient map has kernel precisely
the order-at-most-one space, by
Proposition~\ref{gen:box-decomposition}.
The polarization construction realizes every tuple of
nonnegative principal polynomials, so it is onto the product
in \eqref{gen:eq-real-quotient}.
A real \(A\) and its negative are both admissible if and only
if all their second-order coefficients vanish.
This gives exactly the stated real lineality and quotient.

For complex generators also record the imaginary first-order
polynomials \(q_i\).
Their nonnegativity is independent of the real principal
coefficients, as
\eqref{gen:eq-complex-decomposition} shows.
Both signs of a generator can be admissible only if all these
\(q_i\) and all second-order coefficients vanish.
The remaining imaginary scalar is unrestricted.
This proves the complex assertion.
\end{proof}

\begin{corollary}[Exact binary-slot lift]
\label{gen:slot-lift}
For every generator \(A\) in Theorem~\ref{gen:main-complex},
there is a complete stability-preserving generator
\(\widetilde A\) on the full multiaffine space with
\(|\kappa|\) variables such that
\begin{equation}
 \widetilde A\mathcal P_\kappa=\mathcal P_\kappa A,
 \qquad
 e^{t\widetilde A}\mathcal P_\kappa
          =\mathcal P_\kappa e^{tA}\quad(t\geq0).
 \label{gen:eq-exact-lift}
\end{equation}
It commutes with every permutation of slots within a site.
It can be chosen as a sum of a scalar, one-slot real
M\"obius generators, one-slot imaginary parabolic generators,
and nonnegative multiples of the rank-one apolar generators
\(B_F\) on pairs of slots.
\end{corollary}

\begin{proof}
For the real second-order part use exactly the slot sums
constructed in the proof of
Theorem~\ref{gen:main-real}.
They commute with within-site permutations and descend to the
chosen principal-coefficient realization.
Lift each remaining sitewise drift by
\[
 J_{i,\kappa_i}(q)
 \longmapsto
 \sum_{a=1}^{\kappa_i}
 \left(q(z_{i,a})\partial_{i,a}
                    -\frac12q'(z_{i,a})\right).
\]
Equation~\eqref{gen:eq-polar-one} shows that its diagonalization
on a polarized input is precisely \(J_{i,\kappa_i}(q)\).
The lifted sum is permutation invariant, so this equality
upgrades to intertwining with \(\mathcal P_\kappa\).
The same lift applies to the imaginary quadratics, which are
sums of the parabolic primitives of
Lemma~\ref{gen:imaginary-primitive}.
The scalar lifts as the same scalar identity.

All summands generate complete preserving flows on the full
slot space, hence so does their sum.
This proves the first relation in
\eqref{gen:eq-exact-lift}; exponentiation gives the second.
\end{proof}

The lift is a generator intertwining, not merely a static
polarization followed by an unrelated preserving map.
It is not asserted to have a nonnegative coefficient matrix.
Consequently it imposes no coefficient-positive lifting
principle on a different preservation problem.
Likewise, the signed generation statements proved above are
finite dimensional; they contain no infinite-capacity domain assertion.


\subsection{Exact recognition from local matrices}
\label{gen:recognition-subsection}

The small nonnegativity certificates also give an exact recognition
algorithm when the operator is supplied by its local terms.

\begin{corollary}[Polynomial-time local recognition]
\label{gen:recognition}
Let \(A\) be specified as a finite sum of one-site and two-site operators
on \((\C^2)^{\otimes n}\), with Gaussian-rational entries in the
coefficient basis \(1,z_i\) at each site. There is a deterministic
algorithm, polynomial in \(n\) and the total input bit length, deciding
whether \(e^{tA}\) preserves every nonzero complex disk-stable
multiaffine polynomial for all \(t\ge0\).
The test is on the sum \(A\); the displayed summands need not individually
generate preserving semigroups. No degree bound on the interaction
graph is required.
\end{corollary}

\begin{proof}
In the basis \(1,z\), the Cayley map from disk to half-plane coordinates is
\[
 C=\begin{pmatrix}i&-i\\1&1\end{pmatrix},\qquad
 (Cf)(x)=(x+i)f\left(\frac{x-i}{x+i}\right).
\]
Thus \(B=C^{\otimes n}A(C^{-1})^{\otimes n}\) is the conjugated
half-plane generator. Conjugating a local term requires only a matrix
of order two or four, and preserves its support. All entries remain
Gaussian rational with polynomial bit length.

For a transformed two-site term \(U\), with variables \(x,y\), set
\[
 \begin{aligned}
 p_0&=U1,&p_x&=U(x)-xU1,&p_y&=U(y)-yU1,\\
 p_{xy}&=U(xy)-yU(x)-xU(y)+xyU1.
 \end{aligned}
\]
Checking the four monomials \(1,x,y,xy\) proves
\(U=p_0+p_x\partial_x+p_y\partial_y+p_{xy}\partial_x\partial_y\).
The one-site formula is the same with the \(y\)-terms omitted.
Add these expressions and merge equal monomials. Canonical uniqueness
shows that the result is exactly
\[
 B=P_0+\sum_i P_i\partial_i+\sum_{i<j}P_{ij}\partial_i\partial_j.
\]
It has order at most two automatically. Each \(P_{ij}\) depends only
on its two variables and has bidegree at most \((2,2)\); all remaining
coefficients also have sparse polynomial representations with only a
constant number of contributions per input term.

By Theorems~\ref{gen:main-real} and \ref{gen:main-complex}, acceptance
is equivalent to the following finite tests on the merged coefficients:
\begin{enumerate}
\item each \(P_{ij}\) is real and \(P_{ij}(s,t)\le0\) for every
\((s,t)\in\R^2\);
\item each \(\Im P_i\) is a polynomial \(q_i(z_i)\) of degree at most
two, and \(q_i(s)\ge0\) for every \(s\in\R\);
\item \(\Im P_0+\tfrac12\sum_iq_i'\) is a real constant.
\end{enumerate}
Indeed the real part of \(B\) is a real box endomorphism with the
required principal signs, while the last two conditions give exactly
its allowed imaginary part. Conversely the complex classification
forces all three conditions. The Cayley transformation is invertible,
and so is \(e^{tA}\); stable-or-zero preservation therefore gives
preservation of the nonzero stable class here.

The identity tests require rational arithmetic and merging only
polynomially many sparse records. Forming common denominators by
products shows that the bit lengths of all merged coefficients are
polynomial in the input length. After clearing denominators by a
positive multiplier, each remaining sign test is a real sentence
in at most two variables, with one polynomial of total degree at
most four. Fixed-dimensional, fixed-degree real quantifier elimination
has polynomial bit complexity \cite[Chapter~14]{BasuPollackRoy06}.
There are only polynomially many such sentences. Weak inequalities
include zero coefficients, repeated roots, and degree drops. This
proves the claimed bit bound without expanding the global matrix.
\end{proof}

Efficient tests for bivariate real stability and static univariate
real-rootedness preservers were obtained by Raghavendra, Ryder, and
Srivastava \cite{RaghavendraRyderSrivastava17}. Here the finite-box
generator classification reduces recognition from a compact many-site
local description to fixed-dimensional nonnegativity tests.


\section{Spectral contraction and sector growth}\label{spec:section}

We first study arbitrary complex stability-preserving
semigroups on a finite polynomial space.  No self-adjointness, positivity of
coefficients, or differential representation of the generator is assumed.
Positive damping of polynomial degree has two consequences: it improves the
zero-free domain of the actual evolution, and it gives a sharp contraction
estimate in a norm determined by stable polynomials.
The last subsection returns to real half-plane stability and assumes
coefficient positivity and preservation of homogeneous degree.
It compares spectral bounds between sectors without damping.

\subsection{Weighted kernels and the stable cone}

Zero-capacity coordinates may be omitted.  If no coordinate remains,
the kernel and eigenvector assertions reduce to nonzero constants.
For the gap statements assume that at least one coordinate remains.
Fix \(d\ge1\) and \(\kappa\in\N_{>0}^d\), write
\[
 V_\kappa=\{p\in\C[z_1,\ldots,z_d]:\deg_{z_i}p\le\kappa_i\},
 \qquad b_\alpha=\binom\kappa\alpha,
\]
and let \(\D_r^d=\{z:|z_i|<r\text{ for all }i\}\).
Let \(\mathcal C_\kappa\) consist of zero and all polynomials in
\(V_\kappa\) that do not vanish on \(\D_1^d\).  We call a polynomial
\emph{strictly stable} if it does not vanish on \(\overline{\D_1^d}\).
For \(T\in\operatorname{End}_{\C}(V_\kappa)\), use the kernel
of \eqref{pre:kernel}, namely
\begin{equation}\label{spec:kernel}
 K_T(z,w)=T_z\prod_{i=1}^d(1+z_iw_i)^{\kappa_i}
 =\sum_{\alpha\le\kappa}b_\alpha(Tz^\alpha)(z)w^\alpha.
\end{equation}
The corresponding coefficient pairing is bilinear:
\begin{equation}\label{spec:pairing}
 [f,g]_\kappa=\sum_{\alpha\le\kappa}b_\alpha^{-1}f_\alpha g_\alpha.
\end{equation}
In particular, \([K_T(z,\cdot),p]_\kappa=Tp(z)\).  There is no
complex conjugation in this identity.  In the orthonormal coherent basis
\(\sqrt{b_\alpha}z^\alpha\), the matrix kernel has one square-root
binomial weight on each of its two legs.

We shall use the classical disk contraction theorem, Lemma~\ref{pre:grace},
in the following form.
If two nonzero stable coefficient tensors have matched variables of the
same capacities, their contraction by \eqref{spec:pairing} is stable or
zero whenever each pair of contracted radii has product at least one.
If every such product is greater than one, the contraction is nonzero.
The radii of all uncontracted variables are retained.  This is the
Grace--Asano contraction principle, with the finite-capacity formulation
obtained by polarization; see \cite{BB09,BB10II,WBG26}.  Indeed, a
coefficient of multidegree \(\alpha\) is divided by \(b_\alpha\)
under normalized polarization, and summation over its \(b_\alpha\)
equal subsets leaves exactly the reciprocal weight in
\eqref{spec:pairing}.  Applied to two kernels, this pairing gives the
kernel of their operator product.

\begin{lemma}\label{spec:kernel-criterion}
An invertible operator \(T\) preserves \(\mathcal C_\kappa\) if and
only if \(K_T\) is zero-free on \(\D_1^d\times\D_1^d\).
If \(K_T\) is zero-free on \(\D_r^d\times\D_r^d\), where \(r>1\),
then \(Tp\) is nonzero and zero-free on \(\D_r^d\) for every
\(p\in\mathcal C_\kappa\setminus\{0\}\), without an invertibility
assumption on \(T\) in this last assertion.
\end{lemma}
\begin{proof}
For fixed \(w\in\D_1^d\), the polynomial
\(\prod_i(1+z_iw_i)^{\kappa_i}\) is stable and nonzero.  Its image
under an invertible stability preserver is stable and nonzero.  Evaluating
at any \(z\in\D_1^d\) proves the forward implication.  Conversely,
contract the input variables of \(K_T\) against a stable polynomial.
The contraction theorem gives a stable polynomial or zero; invertibility
excludes zero on a nonzero input.  For the final assertion the contracted
radii have product \(r>1\), so the strict contraction theorem itself
excludes zero and retains output radius \(r\).
\end{proof}

\begin{lemma}\label{spec:cone-interior}
The set \(\mathcal C_\kappa\) is a closed, complex-scalar-invariant
cone containing no complex two-dimensional subspace.  Its interior consists
exactly of the strictly stable polynomials.  Every invertible operator
preserving this cone maps its interior into its interior.
\end{lemma}
\begin{proof}
Closedness follows from Hurwitz's theorem on the connected polydisk.
Every complex two-dimensional subspace contains a nonzero polynomial
annihilated by evaluation at zero, and such a polynomial is not stable.
A polynomial nonzero on the closed polydisk remains so under sufficiently
small coefficient perturbations, by compactness.  Conversely, a boundary
zero \(\zeta\) of a stable polynomial moves into the open polydisk for
the arbitrarily small perturbation \(p((1+\varepsilon)z)\), at
\(z=\zeta/(1+\varepsilon)\).  Zero is not an interior point, since
arbitrarily small nonzero multiples of \(z_1\) do not belong to the
cone.  Finally, an invertible linear operator is open, so its image of
the cone interior is an open subset of the cone.
\end{proof}

\subsection{Strictification by degree damping}

Put \(E_i=z_i\partial_i\), \(E=\sum_iE_i\), and \(P_0p=p(0)\).
The following theorem improves the zero-free domain of the exponential
itself, rather than that of an auxiliary splitting operator.

\begin{theorem}[Degree-damping strictification]\label{spec:strict}
Suppose \(A_0\in\operatorname{End}_{\C}(V_\kappa)\) satisfies
\(e^{tA_0}\mathcal C_\kappa\subseteq\mathcal C_\kappa\) for
\(t\ge0\).  Let \(\mu_i>0\) and
\[
 A=A_0-2\sum_{i=1}^d\mu_i E_i.
\]
For every \(t>0\), the kernel \(K_{e^{tA}}\) is nonzero on the
closed unit polydisk in all \(2d\) variables.  Consequently it has a
zero-free radius \(r_t>1\).  The same radius works for the image of
every nonzero unit-disk-stable input polynomial.

Moreover, for each \(t_0>0\), one radius \(r>1\) works simultaneously
for the kernels \(K_{e^{tA}}\), \(t\ge t_0\), and for their action
on all nonzero unit-disk-stable inputs.
\end{theorem}

\begin{proof}
First consider uniform damping \(A=A_0-2\mu E\).  Fix \(t>0\)
and allow \(\mu\) to vary in the complex half-plane \(\Re\mu>0\).
For \(s\ge0\),
\[
 e^{-2s\mu E}p(z)=p(e^{-2s\mu}z)
\]
preserves disk stability.  Thus every Lie product
\[
 \left(e^{tA_0/m}e^{-2t\mu E/m}\right)^m
\]
preserves stability and is invertible.  The finite-dimensional product
formula, followed by Hurwitz and invertibility of the limiting
exponential, proves preservation by
\(T_t(\mu)=e^{t(A_0-2\mu E)}\).  Lemma
\ref{spec:kernel-criterion} gives
\begin{equation}\label{spec:open-complex-field}
 K_{T_t(\mu)}(z,w)\ne0
 \quad\text{if }z,w\in\D_1^d,\quad\Re\mu>0.
\end{equation}
Each kernel coefficient is an entire function of \(\mu\).

We next establish a real large-damping limit, valid without normality.
Let \(a_{00}=(A_01)(0)\).  In the Euclidean norm on monomial
coefficients, \(E\) is diagonal, nonnegative, and has the constants
as its zero eigenspace.  Hence, for real \(\mu>0\),
\[
 \|e^{-2s\mu E}\|\le1,
 \qquad e^{-2s\mu E}\longrightarrow P_0\quad(s>0).
\]
Expand \(T_t(\mu)\) in its Dyson series about \(-2\mu E\).
The term of order \(k\ge1\) is the integral over
\(0<s_1<\cdots<s_k<t\) of
\[
 e^{-2\mu(t-s_k)E}A_0e^{-2\mu(s_k-s_{k-1})E}
 \cdots A_0e^{-2\mu s_1E}.
\]
Every time interval is positive almost everywhere.  The integrand
therefore converges to \(a_{00}^kP_0\), and its norm is bounded by
\(\|A_0\|^k\), independently of \(\mu\).  Dominated convergence
on the simplex gives the limit \(t^ka_{00}^kP_0/k!\).  These integrated
terms are uniformly bounded by the summable sequence
\(t^k\|A_0\|^k/k!\).  Including the order-zero term and passing
the limit through the sum proves
\begin{equation}\label{spec:large-damping}
 T_t(\mu)\longrightarrow e^{ta_{00}}P_0
 \quad\text{in operator norm as real }\mu\longrightarrow+\infty.
\end{equation}
This is a finite-dimensional instance of strong-coupling limits such as
\cite[Theorem~1]{Burgarth19}; the displayed argument fixes its
normalization directly.  Since \(K_{P_0}=1\), the limiting kernel is
the nonzero constant \(e^{ta_{00}}\).

Fix any \((z,w)\in\overline{\D_1^d}\times\overline{\D_1^d}\).
For \(0<\rho<1\), the functions
\[
 f_\rho(\mu)=K_{T_t(\mu)}(\rho z,\rho w)
\]
are holomorphic and nonvanishing on \(\Re\mu>0\).  Their limit as
\(\rho\uparrow1\) is locally uniform in that half-plane, because
the kernel has finitely many entire coefficient functions.  Hurwitz
therefore says that \(K_{T_t(\mu)}(z,w)\), as a function of
\(\mu\), is identically zero or nowhere zero.  The limit
\eqref{spec:large-damping} excludes the first alternative.  This proves
closed-polydisk nonvanishing at the desired real \(\mu>0\).
Compactness then supplies a radius \(r_t>1\).  Lemma
\ref{spec:kernel-criterion} gives the common input-output assertion.

For nonuniform rates, let \(\mu=\min_i\mu_i>0\).  The operator
\(A_0-2\sum_i(\mu_i-\mu)E_i\) generates a stable semigroup by the
same product-formula argument.  Apply the uniform result to this
operator and the remaining damping \(-2\mu E\).

Finally fix \(t_0>0\), and choose \(r>1\) for
\(B=e^{t_0A/2}\).  For every \(t\ge t_0\),
\[
 e^{tA}=B e^{(t-t_0)A}B.
\]
The middle kernel has radius one, including at \(t=t_0\), and the
outer kernels have radius \(r\).  The two kernel contractions pair
radii whose products are \(r>1\), retaining radius \(r\) on both
external legs.  A further contraction with any stable input proves the
uniform input assertion.
\end{proof}

\begin{remark}\label{spec:strict-scope}
The radius can depend on the capacity, generator, and initial positive
time.  The theorem gives neither a dimension-independent lower bound for
\(r-1\) nor an improvement on inert ancillary variables.  Complex
scalar shifts of the generator do not affect its conclusions.  No
normalization by a trace is used.
\end{remark}

\subsection{Dominant eigenvectors and their dual certificates}

We use the following finite-dimensional complex-cone theorem of Rugh
\cite[Theorem~8.4]{Rugh10}: if a nontrivial closed complex-scalar-invariant cone
contains no complex plane and a linear map sends its nonzero part into
its interior, then the map has an algebraically simple nonzero
eigenvalue of strictly largest modulus.  Convexity of the cone and
self-adjointness of the map are not hypotheses.

\begin{proposition}\label{spec:perron}
Under the hypotheses of Theorem~\ref{spec:strict}, \(A\) has an
algebraically simple eigenvalue \(a_\star\) uniquely largest in real
part.  Choose its right eigenvector \(v\) and left eigenfunctional
\(\ell\) so that \(\ell(v)=1\).  Both \(v\) and
\begin{equation}\label{spec:dual-polynomial}
 Q_\ell(w)=\sum_{\alpha\le\kappa}b_\alpha\ell(z^\alpha)w^\alpha
\end{equation}
are zero-free on a common polydisk of radius greater than one.
Furthermore \(\ell(p)\ne0\) for every nonzero stable \(p\), and
\(e^{-ta_\star}e^{tA}p\to\ell(p)v\).
\end{proposition}
\begin{proof}
For each fixed \(t>0\), Theorem~\ref{spec:strict} and Lemma
\ref{spec:cone-interior} verify all the hypotheses of Rugh's theorem
for \(T=e^{tA}\); the cone contains all constant polynomials
and is therefore nontrivial. Its simple dominant eigenspace is invariant under
\(A\), because \(A\) commutes with \(T\).  Thus it is an
\(A\)-eigenline, say for \(a_\star\).  Spectral mapping gives
strict real-part dominance: another generator eigenvalue with equal
real part would give the same modulus for the exponential.  If two
distinct generator eigenvalues exponentiated to the same value, that
value would have algebraic multiplicity at least two.  A nontrivial
Jordan block at \(a_\star\) would also give a nontrivial exponential
Jordan block, since \(t e^{ta_\star}\ne0\).  Both alternatives
contradict algebraic simplicity for \(T\).

Fix a kernel radius \(r>1\) for \(T\), and put
\(\Lambda=e^{ta_\star}\).  Finite-dimensional spectral theory gives
\[
 \Lambda^{-m}T^m\longrightarrow P=v\otimes\ell,
 \qquad \ell(v)=1.
\]
Every power has kernel radius \(r\), because its internal contraction
radii have product \(r^2>1\).  The kernel transform is injective, so
\(K_P\) is not the zero polynomial.  Hurwitz yields
\[
 K_P(z,w)=v(z)Q_\ell(w)\ne0\qquad(z,w\in\D_r^d).
\]
Both factors consequently have radius \(r\).  The strict weighted
contraction of \(Q_\ell\) with a nonzero unit-stable polynomial is
\(\ell(p)\), so it is nonzero.  The strict real-part separation and
the finite-dimensional spectral expansion then give the asserted
convergence, including when the remaining spectrum has Jordan blocks.
\end{proof}

\subsection{Sharp spectral bounds}

We first record the scalar estimate used in both the spectral and
dynamical arguments.  It is the diameter form of the Schwarz lemma
used in \cite[Theorem~2]{BGLW26}.

\begin{lemma}\label{spec:schwarz}
If \(g\) is holomorphic on \(\D_1^d\), continuous on its closure,
and \(0\le q\le1\), then
\[
 \operatorname{diam}g(q\overline{\D_1^d})
 \le q\operatorname{diam}g(\overline{\D_1^d}).
\]
The same conclusion holds for a coordinatewise dilation all of whose
moduli are at most \(q\).
\end{lemma}
\begin{proof}
The cases \(q=1\), zero diameter, and \(q=0\) are immediate.
For \(x,y\in q\overline{\D_1^d}\), put
\(\rho=\max(\|x\|_\infty,\|y\|_\infty)\).  If \(\rho>0\),
\[
 f(\zeta)=
 \frac{g(\zeta x/\rho)-g(\zeta y/\rho)}
 {\operatorname{diam}g(\overline{\D_1^d})}
\]
is holomorphic on the unit disk, bounded in modulus by one, and
vanishes at zero.  Schwarz at \(\zeta=\rho<1\) bounds
\(|g(x)-g(y)|\) by \(\rho\) times the denominator.  Take the
supremum over \(x,y\).  A coordinatewise dilation has its image in
\(q\overline{\D_1^d}\).
\end{proof}

\begin{theorem}[Complex radius bound]\label{spec:radius-gap}
Suppose \(r>1\) and \(K_T\) is zero-free on
\(\D_r^d\times\D_r^d\).  Then \(T\) has an algebraically simple
nonzero eigenvalue \(\Lambda_\star\) of strictly largest modulus,
with an eigenvector of radius \(r\).  Every other eigenvalue satisfies
\begin{equation}\label{spec:static-ratio}
 |\Lambda/\Lambda_\star|\le r^{-2}.
\end{equation}
No invertibility, Hermiticity, or diagonalizability is assumed.
\end{theorem}
\begin{proof}
The kernel maps the nonzero stable cone into its interior by Lemma
\ref{spec:kernel-criterion}.  Rugh's theorem therefore supplies the
simple dominant eigenvalue.  Its rank-one spectral projection is
nonzero on some interior cone vector, since a proper linear hyperplane
cannot contain an open set.  Normalized iterates of that vector converge
to a nonzero eigenvector \(\psi\) in the closed cone.  Applying \(T\)
once improves \(\psi\) to radius \(r\).

Let \(\phi\) be an eigenvector of a different nonzero eigenvalue,
and put \(g=\phi/\psi\), \(\alpha=\Lambda/\Lambda_\star\).
The function \(g\) is nonconstant and holomorphic on \(\D_r^d\).
For \(1/r<s<r\), we claim
\[
 \alpha g(\D_r^d)\subseteq g(\D_s^d).
\]
Otherwise, for some \(v\in\D_r^d\), the nonzero polynomial
\(f=\phi-\alpha g(v)\psi\) would be stable on \(\D_s^d\), but
\(Tf=\Lambda(\phi-g(v)\psi)\) would vanish at \(v\).
This contradicts strict contraction, since \(rs>1\).

For \(1/r<s<t<r\), let
\(\Delta_u=\operatorname{diam}g(\overline{\D_u^d})\).
These diameters are finite and positive.  The inclusion gives
\(|\alpha|\Delta_t\le\Delta_s\).  Applying Lemma
\ref{spec:schwarz} after rescaling radius \(t\) gives
\(\Delta_s\le(s/t)\Delta_t\).  Letting \(s\downarrow1/r\)
and \(t\uparrow r\) proves \eqref{spec:static-ratio}.  Zero
eigenvalues satisfy the inequality trivially.
\end{proof}

The quantitative mechanism in this proof is the ratio-squeezing argument
of \cite{BGLW26}; the complex-cone theorem supplies the eigenvector
without a positive-semidefinite hypothesis.

\begin{theorem}[Sharp degree-damping gap]\label{spec:gap}
Under the hypotheses of Theorem~\ref{spec:strict}, set
\(\mu=\min_i\mu_i\).  The dominant generator eigenvalue satisfies
\begin{equation}\label{spec:real-gap}
 \Re a_\star-\Re a\ge2\mu
 \qquad(a\in\operatorname{spec}(A),\ a\ne a_\star).
\end{equation}
The constant is sharp for every positive finite capacity vector.
\end{theorem}
\begin{proof}
For small \(h>0\), put
\[
 D_h=e^{-h\sum_i\mu_iE_i},\qquad
 T_h=D_he^{hA_0}D_h.
\]
The kernel of \(e^{hA_0}\) has radius one by Lemma
\ref{spec:kernel-criterion}.  If \(q_i=e^{-h\mu_i}\), direct
coefficient calculation gives
\[
 K_{T_h}(z,w)=K_{e^{hA_0}}(q_1z_1,\ldots,q_dz_d,
                              q_1w_1,\ldots,q_dw_d).
\]
It therefore has radius \(e^{h\mu}\).  Theorem
\ref{spec:radius-gap} gives a simple dominant eigenvalue
\(\Lambda_\star(h)\) and bounds every other eigenvalue by
\(e^{-2h\mu}|\Lambda_\star(h)|\).

Now \(B_h=(T_h-I)/h\to A\) in norm.  Its eigenvalues are exactly
\(b(h)=(\Lambda(h)-1)/h\), with algebraic multiplicities.  These
multisets are bounded and converge to the spectrum of \(A\).  Uniformly
for bounded complex \(b\),
\[
 h^{-1}\log|1+hb|=\Re b+O(h).
\]
The modulus inequality consequently implies
\(\Re b_\star(h)-\Re b(h)\ge2\mu+O(h)\).
Choose a sequence along which \(b_\star(h)\) converges.  After removing
this one eigenvalue from each multiset, the remaining eigenvalues
converge, with multiplicity, to the remaining copies of the spectrum
of \(A\).  The limiting inequality separates all those copies by
\(2\mu\).  In particular, the distinguished limit is algebraically
simple and is the unique maximal-real-part eigenvalue, hence is
\(a_\star\).  This proves \eqref{spec:real-gap} without
differentiable eigenvalue labels or a normal-matrix perturbation bound.
For \(A_0=0\), the eigenvalues are \(-2\sum_i\mu_i\alpha_i\);
a degree-one monomial at a site of minimum rate attains equality.
\end{proof}

\subsection{Contraction along the actual evolution}

For strictly stable \(p\) and arbitrary \(h\in V_\kappa\), define
\begin{equation}\label{spec:osc-definition}
 \operatorname{Osc}_p(h)
 =\operatorname{diam}\{h(z)/p(z):z\in\overline{\D_1^d}\}.
\end{equation}
It is a complex seminorm with kernel \(\C p\): homogeneity and the
triangle inequality follow from the diameter definition, and zero
diameter means that the rational function is constant.

\begin{lemma}\label{spec:range-exclusion}
If \(T\) is an invertible stability preserver, \(p\) is strictly
stable, and \(h\in V_\kappa\), then \(Tp\) is strictly stable and
\[
 (Th/Tp)(\D_1^d)\subseteq(h/p)(\D_1^d).
\]
In particular \(\operatorname{Osc}_{Tp}(Th)\le
\operatorname{Osc}_p(h)\).
\end{lemma}
\begin{proof}
Strictness follows from Lemma~\ref{spec:cone-interior}.  If \(w\)
does not belong to the input ratio image, then \(h-wp\) is a nonzero
stable polynomial.  Its image \(Th-wTp\) is nonzero and stable,
excluding \(w\) from the output ratio image.  Taking diameters proves
the assertion; the open and closed polydisks give the same diameter by
continuity and density.
\end{proof}

\begin{theorem}[Moving-denominator contraction]\label{spec:oscillation}
Let \(A_0\) generate a stability-preserving semigroup, and let
\(A=A_0-2\sum_i\mu_iE_i\), with \(\mu_i\ge\mu\ge0\).
For every strictly stable \(p\), every \(h\in V_\kappa\), and
\(t\ge0\),
\begin{equation}\label{spec:osc-bound}
 \operatorname{Osc}_{e^{tA}p}(e^{tA}h)
 \le e^{-2\mu t}\operatorname{Osc}_p(h).
\end{equation}
The constant is sharp.
\end{theorem}
\begin{proof}
The damping operator \(D_\tau=e^{-2\tau\sum_i\mu_iE_i}\) acts
by a coordinatewise dilation of maximum modulus \(e^{-2\mu\tau}\).
Lemma~\ref{spec:schwarz}, applied to \(h/p\), gives the corresponding
oscillation contraction.  Each \(e^{\tau A_0}\) is nonexpansive by
Lemma~\ref{spec:range-exclusion}.  Iterating along
\((e^{tA_0/m}D_{t/m})^m\) therefore gives exactly the factor
\(e^{-2\mu t}\), with all intermediate denominators strictly stable.

The products converge to \(e^{tA}\), which is invertible and preserves
stability by the product formula and Hurwitz.  Its image of \(p\)
is strictly stable.  The minimum modulus of that limiting polynomial
on the closed polydisk is positive, so coefficient convergence gives
uniform convergence of the ratios there.  Diameters are continuous
under uniform convergence, proving \eqref{spec:osc-bound}.  Equality
holds for pure damping with \(p=1\) and \(h=z_j\) at a minimum-rate
site.
\end{proof}

\begin{corollary}[Nonautonomous evolution]\label{spec:nonautonomous}
On each compact time interval, suppose \(A_0(t)\) and \(\mu_i(t)\)
are bounded and piecewise continuous, with finitely many continuity
pieces, each \(A_0(t)\) a stability generator and \(\mu_i(t)\ge0\).
Let \(U(t,s)\) solve the linear evolution equation with generator
\(A_0(t)-2\sum_i\mu_i(t)E_i\).  Then, for strictly stable \(p\),
\[
 \operatorname{Osc}_{U(t,s)p}(U(t,s)h)
 \le \exp\!\left(-2\int_s^t\min_i\mu_i(u)\,du\right)
       \operatorname{Osc}_p(h).
\]
\end{corollary}
\begin{proof}
Approximate on each continuity piece by step functions taking actual
values of the coefficients.  They converge in \(L^1\).  Their frozen
evolutions satisfy Theorem~\ref{spec:oscillation}, including when the
minimum rate is zero.  Multiplying the bounds gives the exponential
of the integral of the stepwise minimum rates.  These integrals converge
to the stated integral, since the minimum of finitely many real
coordinates is Lipschitz.  The integral-equation estimate for finite
matrix ODEs gives uniform convergence of the propagators.  The limiting
fundamental matrix is invertible and preserves stability by Hurwitz;
Lemma~\ref{spec:cone-interior} gives a strictly stable limiting
denominator.  Uniform convergence of the ratios passes the estimate
to the limit, as in the autonomous proof.
\end{proof}

\begin{corollary}[Adapted norm and resolvent]\label{spec:resolvent}
Assume positive damping, and normalize \(v,\ell\) as in
Proposition~\ref{spec:perron}.  On \(W=\ker\ell\), set
\(\|h\|_v=\operatorname{Osc}_v(h)\) and
\(B=(A-a_\star I)|_W\).  Then \(\|\cdot\|_v\) is a norm and
\[
 \|e^{tB}\|_{v\to v}\le e^{-2\mu t},\qquad
 \|(\zeta I-B)^{-1}\|_{v\to v}
 \le\frac1{\Re\zeta+2\mu}\quad(\Re\zeta>-2\mu).
\]
Every eigenvalue of \(B\) with real part \(-2\mu\) is semisimple.
\end{corollary}
\begin{proof}
The only possible null vectors of the seminorm are multiples of \(v\),
and \(\ell(v)=1\).  Thus it is a norm on \(W\).  The left
eigenfunctional identity makes \(W\) invariant.  Since
\(e^{tA}v=e^{ta_\star}v\), Theorem~\ref{spec:oscillation} gives
the first estimate, including the complex phase of the scalar.
For \(\Re\zeta>-2\mu\), the absolutely convergent integral
\(\int_0^\infty e^{-t\zeta}e^{tB}\,dt\) is the inverse of
\(\zeta I-B\), by integration of its derivative.  Integrating the
norm estimate gives the stated bound.  A nontrivial Jordan block at a
boundary eigenvalue would produce unbounded polynomial growth after
multiplication by \(e^{2\mu t}\), contradicting the first estimate
and norm equivalence on its finite-dimensional generalized eigenspace.
\end{proof}

\begin{example}[No uniform strict radius]\label{spec:no-uniform-radius}
Fix \(\mu,t>0\).  On \(V_1\), in the coefficient basis \((1,z)\),
let \(X\) and \(Z\) be the usual Pauli matrices and put
\[
 A_{0,R}=RX,\qquad
 A_R=A_{0,R}-2\mu E=
 \begin{pmatrix}0&R\\R&-2\mu\end{pmatrix},\qquad R>0.
\]
The undamped flow preserves disk stability, also with arbitrary inert
variables.  Indeed, writing \(c=\cosh(Rs)\), \(b=\sinh(Rs)\),
its action is the explicit weighted substitution
\[
 (e^{sA_{0,R}}p)(z)=(c+bz)
 p\!\left(\frac{b+cz}{c+bz}\right).
\]
The prefactor is nonzero on the disk and the fraction is a disk
automorphism, since \(c^2-b^2=1\).
Put \(\omega_R=\sqrt{R^2+\mu^2}\).  The identity
\((A_R+\mu I)^2=\omega_R^2I\) gives
\[
 e^{tA_R}=e^{-\mu t}\left[
 \cosh(\omega_Rt)I+
 \frac{\sinh(\omega_Rt)}{\omega_R}(A_R+\mu I)\right].
\]
Thus the image of the fixed input \(1\) has its unique zero at
\[
 z_R=-\frac{\omega_R\coth(\omega_Rt)+\mu}{R}.
\]
For every finite \(R\), \(|z_R|>1\), but \(|z_R|\to1\) as
\(R\to\infty\).  Since \(K_{e^{tA_R}}(z,0)=e^{tA_R}1(z)\),
the same zero obstructs a common kernel radius depending only on
\(\mu,t\).  This failure already occurs for Hermitian generators
on one binary variable.  The dominant eigenvector has its zero at
modulus \((\omega_R+\mu)/R\), which also tends to one.
\end{example}

\begin{example}[Unbounded mixing remainder at a fixed gap]\label{spec:nonnormal}
Fix \(\mu>0\), and on \(V_1\) put
\[
 G_R=R(X+iZ),\qquad
 B_R=G_R-2\mu E=
 \begin{pmatrix}iR&R\\R&-iR-2\mu\end{pmatrix},\qquad R>0.
\]
Here \(G_R^2=0\) and \(G_R^*Z+ZG_R=0\).  Its stability-preserving
action is explicit: with \(a=1+iRs\) and \(b=Rs\),
\[
 (e^{sG_R}p)(z)=(a+bz)
 p\!\left(\frac{b+\overline a z}{a+bz}\right),\qquad
 |a+bz|^2-|b+\overline a z|^2=1-|z|^2.
\]
The denominator is nonzero in the disk, and the fraction maps the
disk to itself.  The formula also proves preservation with inert
variables.

Choose the square root \(\omega_R=\sqrt{\mu^2+2i\mu R}\) with
\(x_R=\Re\omega_R>0\).  The two eigenvalues of \(B_R\) are
\(a_\pm=-\mu\pm\omega_R\), and
\[
 x_R=\sqrt{\frac{\mu\sqrt{\mu^2+4R^2}+\mu^2}{2}}\ge\mu.
\]
The dominant spectral projection and normalized exponential are
\[
 P_R=\frac12\left(I+\frac{B_R+\mu I}{\omega_R}\right),\qquad
 S_R(t)=e^{-ta_+}e^{tB_R}
       =P_R+e^{-2\omega_Rt}(I-P_R).
\]
Their off-diagonal entries give
\[
 \|P_R\|_2\ge\frac{R}{2|\omega_R|},\qquad
 |[S_R(t)]_{10}|
 =\frac{R}{2|\omega_R|}|1-e^{-2\omega_Rt}|
 \ge\frac{R}{2|\omega_R|}(1-e^{-2x_Rt}).
\]
Since \(|\omega_R|=(\mu^4+4\mu^2R^2)^{1/4}\), both quantities
diverge as \(R\to\infty\), the second for every fixed \(t>0\).

To keep the actual spectral gap fixed, pass to \(V_{(1,1)}\) and set
\[
 \mathcal A_R=G_R\otimes I-2\mu(E_1+E_2)
             =B_R\otimes I+I\otimes(-2\mu E).
\]
The undamped part is still a joint stability generator by the displayed
weighted substitution.  Its four eigenvalues are
\(a_+,a_-,a_+-2\mu,a_--2\mu\).  Hence its dominant eigenvalue
is \(a_+\), its real-part gap is exactly \(2\mu\), and its
dominant projector is \(\mathcal P_R=P_R\otimes P_0\).
For the fixed strictly stable input \(p(z_1,z_2)=1+z_2/2\), define
\(\mathcal R_R(t)=e^{-ta_+}e^{t\mathcal A_R}-\mathcal P_R\).
The coefficient of \(z_1z_2\) is exactly
\[
 [z_1z_2]\mathcal R_R(t)p
 =\frac12e^{-2\mu t}\frac{R}{2\omega_R}
   (1-e^{-2\omega_Rt}).
\]
The projector term contains no \(z_2\), which proves this identity
directly from the tensor exponential.  Its modulus tends to infinity
for every fixed \(\mu,t>0\).  Thus even after subtraction of the
dominant projector there is no coefficient-norm mixing bound depending
only on the damping, time, and exact spectral gap, in a fixed two-variable
box and on one fixed strictly stable input.  Equivalence of fixed
finite-dimensional norms gives the same obstruction for every fixed
coefficient norm.  This is consistent with Corollary~\ref{spec:resolvent},
whose norm depends on the dominant polynomial.  The example uses linear
normalization by \(e^{ta_+}\); it makes no claim about nonlinear
normalization by an evolving coefficient.
\end{example}

\begin{remark}\label{spec:projective-attribution}
For the stable cone, the forbidden-pencil projective distance of
\cite{Dubois09} is
\(\operatorname{osc}_{\D_1^d}\log|q/p|\).  Its real infinitesimal
form is \(\operatorname{osc}\Re(h/p)\); maximizing over a scalar
phase of \(h\) gives \eqref{spec:osc-definition}.  Thus the metric
and the general complex-cone contraction principle are inherited.
The estimates above specify the exact degree-damping rate for arbitrary
tangent vectors along the actual, possibly nonautonomous evolution.
\end{remark}


\subsection{Sector growth and token-graph spectral concavity}
\label{sector:section}

We now return to real upper-half-plane stability.
There is a second spectral consequence when coefficient positivity and
total degree are preserved. It does not require degree damping or
self-adjointness. A single stable evolution compares the exponential
growth rates of all homogeneous sectors.

\begin{theorem}[Concavity of sector growth]\label{sector:growth}
Let \(V_\kappa\) be a finite real coordinate box and
\(N=\sum_i\kappa_i\). Suppose that \(A\) has nonnegative off-diagonal
entries in the monomial basis, preserves every homogeneous subspace
\(V_{\kappa,k}\), and that \(e^{tA}\) preserves nonnegative real
stability for every \(t\ge0\). Write
\[
 s_k=\max\{\Re\lambda:\lambda\in
             \operatorname{spec}(A|_{V_{\kappa,k}})\},
 \qquad 0\le k\le N.
\]
Then
\[
 2s_k\ge s_{k-1}+s_{k+1}\qquad(1\le k<N).
\]
Neither irreducibility nor symmetry of \(A\) is assumed.
\end{theorem}

\begin{proof}
Start at the stable polynomial
\(f_0(z)=\prod_i(1+z_i)^{\kappa_i}\).
Every allowed monomial has a strictly positive coefficient.
Choose \(\alpha\) such that \(A+\alpha I\) is entrywise nonnegative.
Then \(e^{tA}=e^{-\alpha t}e^{t(A+\alpha I)}\) is entrywise
nonnegative and has a strictly positive diagonal. Consequently
\(f_t=e^{tA}f_0\) still has every allowed coefficient strictly positive.
Its diagonal restriction is a degree-\(N\) polynomial
\[
 P_t(x)=f_t(x,\ldots,x)=\sum_{k=0}^N c_k(t)x^k,
 \qquad c_k(t)>0,
\]
whose roots are all strictly negative.

Newton's inequalities give
\begin{equation}\label{sector:newton}
 \left(\frac{c_k(t)}{\binom Nk}\right)^2
 \ge
 \frac{c_{k-1}(t)}{\binom N{k-1}}
 \frac{c_{k+1}(t)}{\binom N{k+1}}.
\end{equation}
Here is the particular form needed. The derivative
\(R=P_t^{(k-1)}\) has \(m=N-k+1\) negative roots, so
\(R(x)=R(0)\prod_{j=1}^m(1+u_jx)\) with \(u_j>0\).
Cauchy's inequality gives
\[
 \sum_{i<j}u_iu_j
 \le\frac{m-1}{2m}\left(\sum_i u_i\right)^2.
\]
The first three coefficients of \(R\) are
\((k-1)!c_{k-1}\), \(k!c_k\), and
\((k+1)!c_{k+1}/2\). Substitution gives
\eqref{sector:newton}, including its binomial normalization.

Let \(A_k=A|_{V_{\kappa,k}}\), and let \(v_k\) be the coefficient
vector of the degree-\(k\) part of \(f_0\). Degree preservation gives
\[
 c_k(t)=\mathbf1^{\mathsf T}e^{tA_k}v_k.
\]
Since every entry of \(v_k\) is positive, this quantity is bounded
above and below by fixed positive multiples of the sum of all
entries of \(e^{tA_k}\). For an entrywise nonnegative \(m_k\)-by-\(m_k\)
matrix \(M\), that sum lies between \(\|M\|_1\) and
\(m_k\|M\|_1\), where \(\|\cdot\|_1\) is the induced maximum
column-sum norm. Spectral mapping gives
\(\|e^{tA_k}\|_1\ge e^{ts_k}\), and Jordan normal form gives
\[
 \|e^{tA_k}\|_1
 \le C_k(1+t^{m_k-1})e^{ts_k}\qquad(t\ge0)
\]
for a fixed constant \(C_k\). Therefore
\[
 \lim_{t\to\infty}t^{-1}\log c_k(t)=s_k.
\]
Take logarithms in \eqref{sector:newton}, divide by \(t\), and
let \(t\to\infty\). The fixed binomial factors disappear, proving
the assertion. For \(N\le1\) there is no inequality to check.
\end{proof}

The proof uses coefficient sums, not
\(\operatorname{Tr}(e^{tA_k})\). No real-rooted generating function
for the latter traces is asserted. The passage from coefficient
log-concavity to concavity of growth exponents is the one-variable
tropical-stability principle
\cite[Theorem~5 and the discussion following it]{Branden10Discrete}.
The proof above establishes the needed exponential version directly,
including nonsymmetric blocks and their Jordan factors.
The preservation theorem supplies the input needed to apply it.

Token-graph spectral theory includes the nesting of Laplacian spectra
proved by Dalf\'o and coauthors \cite{DalfoEtAl21}, and the study of
adjacency spectral radii and new Laplacian spectral layers by Reyes,
Dalf\'o, and Fiol \cite{ReyesDalfoFiol24}. The monotonicity assertions
addressed below are those for adjacency and signless-Laplacian largest
eigenvalues in \cite[Conjectures~5--6]{APS26}. The new-layer Laplacian
quantities in \cite[Definition~3.1 and Conjecture~3.3]{ReyesDalfoFiol24}
are different spectral invariants.

Let \(G\) be a finite simple graph on \([n]\), with nonnegative edge
weights \(w_{ij}\). Its \(k\)-token graph has the \(k\)-subsets as
vertices. The edge from \(S\) to \(S-\{i\}+\{j\}\) has weight
\(w_{ij}\) whenever \(i\in S\), \(j\notin S\), and \(ij\) is an
edge of \(G\). Denote its adjacency and weighted degree matrices by
\(A_k(G)\) and \(D_k(G)\). The sectors \(k=0,n\) are singletons
with zero matrices.

\begin{theorem}[Token-graph spectral concavity]\label{sector:token}
For every \(a\in[-1,1]\), put
\[
 \Lambda_k(a)=\lambda_{\max}(A_k(G)+aD_k(G)).
\]
For every nonnegatively weighted \(G\),
\[
 2\Lambda_k(a)\ge\Lambda_{k-1}(a)+\Lambda_{k+1}(a)
 \qquad(1\le k<n),\qquad
 \Lambda_k(a)=\Lambda_{n-k}(a).
\]
In particular the sequence is nondecreasing up to half filling.
The interval \([-1,1]\) is the exact interval of real parameters
for which this concavity holds for every such graph.
\end{theorem}

\begin{proof}
On the binary polynomial box define, for every unordered edge,
\[
 \mathcal B_{ij}=w_{ij}\bigl[
 (z_j+az_i)\partial_i+(z_i+az_j)\partial_j
 -(z_i^2+2az_iz_j+z_j^2)\partial_i\partial_j\bigr].
\]
It annihilates \(1,z_iz_j\) and sends
\(z_i\) to \(w_{ij}(az_i+z_j)\) and
\(z_j\) to \(w_{ij}(z_i+az_j)\).
Thus \(\mathcal B=\sum_{ij}\mathcal B_{ij}\) preserves total
degree, is a real symmetric Metzler matrix, and restricts on the
degree-\(k\) sector to \(A_k(G)+aD_k(G)\).
Its mixed second-order coefficient is
\[
 -w_{ij}\bigl[(s+at)^2+(1-a^2)t^2\bigr]\le0.
\]
Theorem~\ref{gen:main-real} therefore gives stability preservation.
Theorem~\ref{sector:growth} yields concavity, since for a real
symmetric matrix its spectral bound is its largest eigenvalue.

Complementation of subsets gives a weight-preserving isomorphism
between the \(k\)- and \((n-k)\)-token graphs, proving symmetry.
Set \(\delta_k=\Lambda_k-\Lambda_{k-1}\). Concavity says that
the \(\delta_k\) are nonincreasing, while symmetry gives
\(\delta_{n-k+1}=-\delta_k\). If \(k\le(n+1)/2\), it follows that
\(\delta_k\ge-\delta_k\), proving monotonicity.

For sharpness, take the unweighted star \(K_{1,3}\). Its two-token
graph is a six-cycle. On the equal-leaf subspace, the one-token
matrix is
\(\left(\begin{smallmatrix}3a&\sqrt3\\\sqrt3&a\end{smallmatrix}\right)\);
its orthogonal complement has eigenvalue \(a\). Thus
\[
 \Lambda_1(a)=2a+\sqrt{a^2+3},\qquad
 \Lambda_2(a)=2a+2.
\]
For \(a>1\), \(\Lambda_1(a)>\Lambda_2(a)\); symmetry then
contradicts concavity at \(k=2\). For \(a<-1\), a single edge
has \(\Lambda_0=\Lambda_2=0\) and \(\Lambda_1=a+1<0\), again
violating concavity. At \(a=1\), the star has
\(\Lambda_1=\Lambda_2=4\), so connectedness does not imply
strict monotonicity.
\end{proof}

For completeness, the local preservation used here has a short
verification independent of the generator classification. For
\(\tau\ge0\), set \(b=e^{a\tau}\sinh\tau\) and
\(c=e^{a\tau}\cosh\tau\). The unweighted edge exponential fixes
\(1,xy\) and sends \(x,y\) to \(cx+by,bx+cy\).
Its algebraic symbol is
\[
 xy+uv+b(xu+yv)+c(xv+yu).
\]
Writing this symbol as \(\zeta^{\mathsf T}M\zeta/2\) in the
order \(\zeta=(x,y,u,v)^{\mathsf T}\), one has
\[
 M=\begin{pmatrix}
 0&1&b&c\\
 1&0&c&b\\
 b&c&0&1\\
 c&b&1&0
 \end{pmatrix}.
\]
Its eigenvalues are
\[
 1+e^{(a+1)\tau},\quad 1-e^{(a+1)\tau},\quad
 -1-e^{(a-1)\tau},\quad -1+e^{(a-1)\tau}.
\]
Exactly one is positive when \(-1\le a\le1\). The quadratic
form is positive on every strictly positive real vector \(Y\).
A zero at \(X+\mathrm iY\) would imply
\(X^{\mathsf T}MY=0\) and
\(X^{\mathsf T}MX=Y^{\mathsf T}MY>0\), contradicting the
one-positive-eigenvalue property. Hence the symbol is stable.
The classical symbol theorem \cite{BB09} gives complete
preservation for each edge. The finite-dimensional product formula
and coefficientwise closure give preservation for their sum;
invertibility of the limiting exponential excludes zero outputs.

Taking \(a=1\) and \(a=0\) gives, respectively, monotonicity of
the largest eigenvalues of the signless Laplacian and the adjacency
matrix, as conjectured in \cite[Conjectures~5--6]{APS26}.
Theorem~\ref{sector:token} proves the stronger weighted concavity
throughout the displayed interpolation interval.
The adjacency-preserving generator at \(a=0\) is already contained
in Purbhoo's construction \cite[Proposition~4.5]{Purbhoo18}.
The contribution here is the cross-sector spectral conclusion.
The bounds in Conjectures~1--4 of \cite{APS26} were settled by
\cite{BBKL26} and are not claimed here.

In the occupation basis, the same block matrices arise from
\[
 \mathsf B_a(G)=\frac12\sum_{ij}w_{ij}
 \bigl[X_iX_j+Y_iY_j+a(I-Z_iZ_j)\bigr].
\]
Thus its maximum energy over fixed particle number is concave in
that number. The negative Hamiltonian has the corresponding convex
sector ground-state energies. At \(a=1\), the common one-qubit
rotation \(R=e^{-\mathrm i\pi X/4}\), with
\(RYR^*=Z\) and \(RZR^*=-Y\), identifies this operator with the
EPR Hamiltonian in \cite{APS26}. There is a
maximum-energy state in a half-filled sector in the rotated basis.
This locates a maximizing sector; it does not give an efficient
algorithm within that exponentially large sector. For \(a<0\)
the quantity in Theorem~\ref{sector:token} is the largest
eigenvalue, not in general the spectral radius.


\section{Classical computation of principal eigenvalues and states}
\label{alg:section}

The spectral theorem of the preceding section has an algorithmic consequence
for operators with a bounded-degree local description.  A positive damping
rate selects a holomorphic eigenvalue throughout a complex half-plane.
This domain can be joined to a disk on which perturbation coefficients are
efficiently computable.  An explicit change of variable then reduces
evaluation at any fixed positive damping rate to logarithmically many
coefficients.

The coefficient algorithm is due to Bravyi, DiVincenzo, and Loss
\cite{BDL08}.  Our use of it includes a complexification argument, so that
every invocation of their algorithm remains within its Hermitian
hypotheses.  Neither an extension of Hermitian eigenvalue ordering to
nonnormal matrices nor numerical continuation of eigenvectors is needed.
We first estimate the principal eigenvalue.  Starting in
Section~\ref{alg:states-subsection}, we then compute stable product
tests of its eigenvectors, including relative ground-state amplitudes,
product-event probabilities, and adaptive equatorial measurement
records.  These scalar queries do not require a bound on the condition
number of the spectral projection.

\subsection{Local input and the approximation theorem}
\label{alg:input-subsection}

Identify the coefficient space
\[
 \mathcal P_n
 =\{p\in\C[z_1,\ldots,z_n]:\deg_{z_i}p\le1\}
 \cong(\C^2)^{\otimes n},
 \qquad z^\alpha\longleftrightarrow|\alpha\rangle.
\]
We use the Euclidean norm in this coefficient basis and the induced
operator norm.  Let \(\mathcal S_n\) be the set of nonzero polynomials in
\(\mathcal P_n\) that do not vanish in \(\D^n\), where
\(\D=\{z\in\C:|z|<1\}\).  The degree operator is
\begin{equation}
 E=\sum_{i=1}^n |1\rangle\langle1|_i
   =\sum_{i=1}^n z_i\partial_{z_i}.
 \label{alg:degree}
\end{equation}
All local matrices below are supplied in the displayed tensor-product
basis.  A Gaussian rational means an element of \(\mathbb Q(i)\).

\begin{theorem}[Principal-eigenvalue approximation]
\label{alg:complex}
Fix \(d\in\mathbb N\cup\{0\}\) and positive rational numbers \(J,\mu\).
For \(n\ge1\), suppose that the input specifies
\begin{equation}
 A_0=\sum_{i=1}^n K_i+
       \sum_{\{i,j\}\in\mathcal E}K_{ij},
 \label{alg:local-input}
\end{equation}
where the graph \((\{1,\ldots,n\},\mathcal E)\) has maximum degree at
most \(d\), every local matrix has Gaussian-rational entries and operator
norm at most \(J\), and
\begin{equation}
 e^{tA_0}\mathcal S_n\subseteq\mathcal S_n
 \qquad(t\ge0).
 \label{alg:promise}
\end{equation}
Let \(a_\star\) be the eigenvalue of
\(A=A_0-2\mu E\) with largest real part.  This eigenvalue is unique
and algebraically simple.  Given a positive rational \(\delta\), a
deterministic classical algorithm returns
\(\widehat a\in\mathbb Q(i)\) such that
\[
 |\widehat a-a_\star|\le\delta.
\]
For fixed \(d,J,\mu\), its bit complexity is polynomial in \(n\),
the input bit length, and \(\delta^{-1}\).
\end{theorem}

Only the sum \(A_0\) is assumed to satisfy
\eqref{alg:promise}; the displayed local terms need not individually
generate preserving semigroups.  The theorem uses the supplied local
representation. Corollary~\ref{gen:recognition} verifies its
preservation hypothesis in polynomial bit time. The local representation
itself is part of the input.
No positivity, Hermiticity, stochastic normalization, or diagonalizability
is assumed.

We will give explicit choices of all continuation parameters.  In
particular, put
\begin{equation}
 N=2n,\qquad D=d+1,\qquad
 \rho=\frac{2^{-18}}{DJ},\qquad c=\frac1\mu.
 \label{alg:parameters}
\end{equation}
The number \(p\) of perturbation coefficients required by the proof
satisfies
\begin{equation}
 p=O\!\left(
   \left(1+\frac{c}{\rho}\right)^2
   \log\!\left(2+\frac{C_{d,J,\mu}n}{\delta}\right)
 \right),
 \label{alg:order-bound}
\end{equation}
for a constant \(C_{d,J,\mu}\) independent of \(n,\delta\).
The arithmetic operation count is
\(N\exp(O_D(p))\), with a polynomial bit overhead.
In particular, after substitution the coefficient computation takes
\[
 N\left(2+\frac{C_{d,J,\mu}n}{\delta}\right)^%
 {O_D((1+2^{18}DJ/\mu)^2)}
\]
arithmetic operations, up to polynomial factors in the coefficient
order. The constant implicit in \(O_D\) retains the graph-degree
dependence of the perturbation algorithm.
The large constants in \eqref{alg:parameters} are inherited from
the perturbation estimate.  They are sufficient for a complexity
statement, not a claim of practical efficiency.  In particular,
Theorem~\ref{alg:complex} does not assert polynomial dependence
on \(\mu^{-1}\).

\subsection{Perturbation coefficients and complexification}
\label{alg:perturbation-subsection}

We first state precisely the published perturbation input.
For a graph on \(N\) qubits with maximum degree at most \(D\ge1\), set
\begin{equation}
 H_0=2\sum_{u=1}^N |1\rangle\langle1|_u,
 \qquad \Omega=|0\rangle^{\otimes N}.
 \label{alg:unperturbed}
\end{equation}
For any matrix \(V\), the eigenvalue of \(H_0+xV\) issuing from the
simple eigenvalue \(0\) has a convergent germ at \(x=0\), which we write
\begin{equation}
 F_V(x)=\sum_{k\ge1}e_k(V)x^k.
 \label{alg:energy-germ}
\end{equation}

\begin{proposition}[Bravyi--DiVincenzo--Loss]
\label{alg:bdl}
Let \(W=\sum_{\{u,v\}}W_{uv}\) be Hermitian, with each
\(W_{uv}\) Hermitian and of norm at most \(L>0\).
For \(H_0\) as in \eqref{alg:unperturbed}, put
\[
 R_L=\frac{2^{-16}}{DL}.
 \]
Then
\begin{equation}
 |e_k(W)|\le2^{-15}N R_L^{-k}\qquad(k\ge1).
 \label{alg:bdl-bound}
\end{equation}
For every \(p\ge1\), the coefficients
\(e_1(W),\ldots,e_p(W)\) can be computed in
\(N\exp(O_D(p))\) arithmetic operations.
\end{proposition}

These are Lemma~6 and Theorem~3 of \cite{BDL08}, specialized to
unperturbed gap \(2\).  Their Theorem~2 also gives convergence of the
energy expansion throughout the corresponding high-field disk.
The coefficient algorithm is formal: computing the first \(p\)
coefficients does not require the eventual evaluation point to lie in
that disk.  We use the coefficient bound and algorithm exactly in this
Hermitian form.  The required rational-input bit implementation is
specified below.

\begin{lemma}[Complex perturbation coefficients]
\label{alg:complexification}
Let \(V=\sum_{\{u,v\}}V_{uv}\), where the local matrices are arbitrary
complex matrices with \(\|V_{uv}\|\le J\).
The germ \eqref{alg:energy-germ} extends to a holomorphic eigenvalue
function on
\[
 |x|<\rho,\qquad \rho=\frac{2^{-18}}{DJ},
 \]
and its coefficients satisfy
\begin{equation}
 |e_k(V)|\le2^{-15}N\rho^{-k}\qquad(k\ge1).
 \label{alg:complex-bound}
\end{equation}
If the entries of the local matrices are Gaussian rational, for each
integer \(p\ge1\) its first \(p\) coefficients can be computed exactly using
\(N\exp(O_D(p))\) arithmetic operations on Gaussian rationals.
\end{lemma}

\begin{proof}
Let \(Q=\Id-|\Omega\rangle\langle\Omega|\), and let \(S\) be the
inverse of \(H_0\) on \(Q(\C^2)^{\otimes N}\), extended by zero on
\(\Omega\).  Normalize the eigenvector germ by
\[
 v(x)=\Omega+\sum_{k\ge1}v_kx^k,\qquad
 \langle\Omega,v_k\rangle=0.
 \]
Writing \(v_0=\Omega\), coefficient comparison in
\((H_0+xV)v(x)=F_V(x)v(x)\) gives
\begin{align}
 e_k(V)&=\langle\Omega,Vv_{k-1}\rangle,
 \label{alg:formal-energy}\\
 v_k&=-S Q\left(
        Vv_{k-1}-\sum_{j=1}^{k-1}e_j(V)v_{k-j}
       \right).
 \label{alg:formal-vector}
\end{align}
The vacuum coordinate functional in these formulas is fixed; it does
not conjugate any variable matrix entry of \(V\).
Induction shows that \(v_k\) and \(e_k(V)\) are homogeneous
polynomials of degree \(k\) in the entries of \(V\), with no conjugate
variables.  The same recursion also proves uniqueness of the formal
normalized eigenpair.  Thus the coefficients returned by
Proposition~\ref{alg:bdl} are the restrictions of these polynomials
to Hermitian inputs.

Decompose the local terms as
\[
 B_{uv}=\frac{V_{uv}+V_{uv}^*}{2},\qquad
 C_{uv}=\frac{V_{uv}-V_{uv}^*}{2i},
 \]
and write \(V=B+iC\) for the resulting sums.
Both \(B_{uv}\) and \(C_{uv}\) are Hermitian of norm at most \(J\).
For real \(\theta\), the local perturbation
\[
 W_\theta=B\cos\theta+C\sin\theta
 =\frac12 e^{-i\theta}V+\frac12 e^{i\theta}V^*
 \]
is Hermitian, with each local norm at most \(2J\).
Since \(e_k\) is homogeneous of degree \(k\), the Fourier coefficient
of frequency \(-k\) in \(e_k(W_\theta)\) is
\(2^{-k}e_k(V)\).  Indeed, a monomial of degree \(k\) contributes
that frequency only when its \(V\)-term is selected in all \(k\)
factors.  Consequently
\begin{equation}
 e_k(V)=\frac{2^k}{2\pi}
       \int_0^{2\pi} e^{ik\theta}e_k(W_\theta)\,d\theta.
 \label{alg:fourier}
\end{equation}
Apply \eqref{alg:bdl-bound} with \(L=2J\).  The corresponding
radius is \(2^{-17}/(DJ)\), so \eqref{alg:fourier} gives
\[
 |e_k(V)|
 \le 2^k\,2^{-15}N
          \left(\frac{2^{-17}}{DJ}\right)^{-k}
 =2^{-15}N\rho^{-k}.
 \]
This proves normal convergence on compact subsets of \(|x|<\rho\).
Near zero the sum is the actual eigenvalue germ, by
\eqref{alg:formal-energy}--\eqref{alg:formal-vector} and the
simple eigenvalue of \(H_0\).  The identity theorem applied to
\[
 \det\bigl(H_0+xV-F_V(x)\Id\bigr)
 \]
therefore proves the eigenvalue identity on the entire disk.
This step does not assert simplicity there or make any ordering
claim about the other eigenvalues.

For the coefficient algorithm, take the distinct rational nodes
\[
 s_\ell=-1+\frac{2\ell}{p},\qquad 0\le\ell\le p.
 \]
At each node, \(B+s_\ell C\) is a Hermitian local perturbation
with norm bound \(2J\).  Apply
Proposition~\ref{alg:bdl} at each of these \(p+1\) inputs.
For \(k\le p\), the function
\[
 s\longmapsto e_k(B+sC)
 \]
is a polynomial of degree at most \(k\).
Its values at the nodes determine its value at \(s=i\), namely
\(e_k(V)\), by exact rational interpolation.
The factor \(p+1\) in the number of calls and the polynomial
interpolation cost are absorbed in \(N\exp(O_D(p))\).
Thus the integral in \eqref{alg:fourier} is used only for a bound;
the algorithm uses no numerical quadrature.
\end{proof}

\begin{remark}
\label{alg:complexification-scope}
The factor four between \(R_J\) and \(\rho\) has two sources:
the Hermitian family has local norm bound \(2J\), and extracting
the extremal Fourier coefficient contributes \(2^k\).
The argument transfers the coefficient estimate, not a Hermitian
spectral-ordering or spectral-simplicity theorem.  This distinction
is necessary for arbitrary nonnormal inputs.
\end{remark}

One-site perturbations can be incorporated without altering this
published input.  Attach a private auxiliary qubit \(i'\) to each
physical site \(i\), and represent a one-site term \(V_i\) by
\(V_i\otimes\Id_{i'}\) on the edge \(\{i,i'\}\).
Give every auxiliary qubit the unperturbed term
\(2|1\rangle\langle1|_{i'}\).  The resulting operator is exactly
\begin{equation}
 \widetilde H(x)
  =(2E+xV)\otimes\Id_{\mathrm{aux}}
   +\Id_{\mathrm{phys}}\otimes
        2\sum_{i'}|1\rangle\langle1|_{i'}.
 \label{alg:auxiliary}
\end{equation}
Its vacuum eigenvalue germ and all its coefficients agree with
those of the physical operator.
The number of qubits is \(N=2n\), the maximum degree is at most
\(D=d+1\), and the local norm bound is unchanged.
The identity on the auxiliary factor is important: the auxiliary
qubit is not coupled to the physical perturbation.
This construction is only a device for using the coefficient
algorithm; it imposes no stability assumption on the auxiliary
description.

\subsection{A global eigenvalue branch}
\label{alg:branch-subsection}

\begin{lemma}[Half-plane continuation]
\label{alg:half-plane}
Under \eqref{alg:promise}, the matrix
\[
 A(h)=A_0-2hE
 \]
has an algebraically simple eigenvalue \(a(h)\), uniquely maximizing
real part, for every \(h\) with \(\Re h>0\).
The function \(a\) is holomorphic on that half-plane.
For \(V=-A_0\), the function
\begin{equation}
 F(x)=-x\,a(1/x),\qquad \Re x>0,
 \label{alg:reciprocal}
\end{equation}
is a holomorphic eigenvalue of \(2E+xV\).
It agrees with the vacuum germ of
Lemma~\ref{alg:complexification} near the positive real origin.
\end{lemma}

\begin{proof}
Write \(h=u+iv\), where \(u>0\).
The imaginary Euler flow acts on polynomials by
\[
 e^{-2itvE}p(z_1,\ldots,z_n)
   =p(e^{-2itv}z_1,\ldots,e^{-2itv}z_n),
 \]
and therefore preserves \(\mathcal S_n\).
The finite-dimensional product formula gives
\[
 e^{t(A_0-2ivE)}
   =\lim_{m\to\infty}
       \bigl(e^{tA_0/m}e^{-2itvE/m}\bigr)^m.
 \]
For each stable input, every approximating output is stable.
The limit is stable or zero by Hurwitz's theorem, and it cannot
be zero on a nonzero input because the limiting exponential is
invertible.  Hence \(A_0-2ivE\) also generates a preserving
semigroup.

Theorem~\ref{spec:gap}, applied with damping \(u\), supplies an
algebraically simple eigenvalue of
\((A_0-2ivE)-2uE=A(h)\) whose real part exceeds those of all the
others by at least \(2u\).  At each point, the implicit function
theorem for the characteristic polynomial makes this eigenvalue
locally holomorphic.  Spectral continuity preserves its strict
real-part inequality in a sufficiently small neighborhood.
Thus the locally defined functions agree wherever their
neighborhoods overlap: they select the same uniquely distinguished
eigenvalue.  This proves the global holomorphic assertion.

The reciprocal map takes the right half-plane into itself, so
\eqref{alg:reciprocal} is holomorphic there.  The matrix identity
\[
 -x A(1/x)=2E+xV
 \]
proves the eigenvalue assertion.  For positive real \(x\), the
scalar \(-x\) reverses real-part order, so \(F(x)\) is the unique
eigenvalue of \(2E+xV\) with smallest real part.  We do not make
this ordering assertion for nonreal \(x\).

It remains to identify the high-field germ without assuming
Hermiticity of \(V\).  Choose positive \(x\) so small that
\(x\|V\|<1/4\).
Normality of \(2E\) implies that every eigenvalue of
\(2E+xV\) is within distance \(1/4\) of
\(\{0,2,\ldots,2n\}\).  To see the inclusion, outside these disks
the resolvent of \(2E\) has norm at most \(4\), so its
product with \(xV\) has norm less than \(1\).
The Neumann series therefore makes the perturbed resolvent
invertible, including when \(V=0\).
The Riesz projection inside \(|z|=1\) has rank one:
the contour stays in the resolvent set throughout
\(2E+\tau xV\), \(0\le\tau\le1\), and its rank at \(\tau=0\)
is one.  Hence the eigenvalue near zero is precisely the vacuum
branch.  All other eigenvalues have real part greater than one.
The vacuum branch is therefore the uniquely minimal real-part
eigenvalue, proving its agreement with \(F(x)\) on a positive
real interval.
\end{proof}

Apply the identity theorem on the connected overlap
\(\{\Re x>0\}\cap\{|x|<\rho\}\).
Lemma~\ref{alg:half-plane} and
Lemma~\ref{alg:complexification}, with the auxiliary construction
when needed, yield one holomorphic function on
\begin{equation}
 \Omega_\rho=\{x\in\C:\Re x>0\}\cup\{x\in\C:|x|<\rho\}.
 \label{alg:glued-domain}
\end{equation}
It remains an eigenvalue of \(\widetilde H(x)\) throughout this
domain, by \eqref{alg:auxiliary} and the characteristic-polynomial
identity.  In particular,
\begin{equation}
 |F(x)|\le\|\widetilde H(x)\|
          \le2N+|x|W,\qquad W=NDJ.
 \label{alg:norm-bound}
\end{equation}
The displayed \(W\) is a convenient rational upper bound on the
sum of all local norms.  The scalar estimate
\(|\lambda|\le\|T\|\) for an eigenvalue of a matrix \(T\)
does not require normality.
Finally,
\begin{equation}
 a_\star=-\mu F(1/\mu).
 \label{alg:target}
\end{equation}

\subsection{Effective continuation and exact arithmetic}
\label{alg:continuation-subsection}

The next lemma gives the continuation step independently of its
operator-theoretic application.

\begin{lemma}[An explicit continuation map]
\label{alg:map}
Let \(\rho,c>0\), and suppose that \(F\) is holomorphic on
\(\Omega_\rho\).
Define
\begin{equation}
 \eta=\frac{\rho^2}{4c+2\rho},\qquad
 R=c+\eta,\qquad q=\frac cR,\qquad
 b=\frac{R^2-c^2}{R},\qquad
 \phi(w)=\frac{bw}{1-qw}.
 \label{alg:map-parameters}
\end{equation}
Then \(0<q<1\), \(\phi(\D)\subset\Omega_\rho\),
\(\phi(0)=0\), and \(\phi(q)=c\).
If \(|F(x)|\le L_0+L_1|x|\) on \(\Omega_\rho\), the function
\(G=F\circ\phi\) satisfies
\[
 |G(w)|\le M:=L_0+L_1(c+R)\qquad(w\in\D).
 \]
Writing \(F(x)=\sum_{j\ge0}f_jx^j\) at zero and
\(G(w)=\sum_{k\ge0}g_kw^k\), one has \(g_0=f_0\) and
\begin{equation}
 g_k=\sum_{j=1}^k
        f_j b^j q^{k-j}\binom{k-1}{j-1}
       \qquad(k\ge1).
 \label{alg:coefficient-transform}
\end{equation}
For every \(p\ge0\),
\begin{equation}
 \left|F(c)-\sum_{k=0}^p g_kq^k\right|
    \le\frac{Mq^{p+1}}{1-q}.
 \label{alg:tail}
\end{equation}
\end{lemma}

\begin{proof}
The definitions give \(q\in(0,1)\), \(\eta\le\rho/2\), and
\[
 0<R^2-c^2=2c\eta+\eta^2<\rho^2.
 \]
If \(|x-c|<R\) and \(\Re x\le0\), then
\[
 |x|^2=|x-c|^2+2c\Re x-c^2<R^2-c^2<\rho^2.
 \]
Thus the disk of center \(c\) and radius \(R\) lies in
\(\Omega_\rho\).
The identity
\[
 \phi(w)=c+R\frac{w-q}{1-qw}
 \]
shows that \(\phi\) maps \(\D\) biholomorphically onto this
disk.  Direct substitution gives its values at \(0\) and \(q\).
Since \(|\phi(w)|<c+R\), the bound on \(G\) follows.

For \(j\ge1\), expansion at zero gives
\[
 \phi(w)^j=b^j w^j(1-qw)^{-j}
  =b^j\sum_{k\ge j}
       \binom{k-1}{j-1}q^{k-j}w^k.
 \]
Only \(j\le k\) contributes to the coefficient of \(w^k\),
which proves \eqref{alg:coefficient-transform}.
Applying Cauchy's estimate on circles with radii increasing to
one gives \(|g_k|\le M\).
The geometric tail at \(w=q\) now proves \eqref{alg:tail}.
No holomorphic extension to the boundary of \(\D\) is used.
\end{proof}

\begin{proof}[Proof of Theorem~\ref{alg:complex}]
The existence, uniqueness, and simplicity of \(a_\star\) follow
from Theorem~\ref{spec:gap}.
Use \(V=-A_0\), the auxiliary construction
\eqref{alg:auxiliary}, and the parameters
\eqref{alg:parameters}.
Lemmas~\ref{alg:complexification} and
\ref{alg:half-plane} give the function \(F\) on
\(\Omega_\rho\), with \eqref{alg:norm-bound}.
Apply Lemma~\ref{alg:map} with
\[
 L_0=2N,\qquad L_1=W,\qquad
 M=2N+(c+R)W.
 \]
All the parameters in \eqref{alg:map-parameters} are rational.

Choose the smallest integer \(p\ge0\) for which
\begin{equation}
 \mu M q^{p+1}\le\delta(1-q).
 \label{alg:stopping}
\end{equation}
The algorithm finds this integer by rational multiplication and
comparison; evaluation of logarithms is unnecessary.
If \(p>0\), compute \(e_1(V),\ldots,e_p(V)\) by
Lemma~\ref{alg:complexification}, and then compute
\(g_0,\ldots,g_p\) by
\eqref{alg:coefficient-transform}.
Here \(g_0=F(0)=0\).
Return the Gaussian rational
\begin{equation}
 \widehat a=-\mu\sum_{k=0}^p g_kq^k.
 \label{alg:output}
\end{equation}
For \(p=0\), no coefficient call is made and the same formula
returns zero.
Equations \eqref{alg:target}, \eqref{alg:tail}, and
\eqref{alg:stopping} give
\(|\widehat a-a_\star|\le\delta\).

We first bound the number of arithmetic operations.
A useful exact identity is
\begin{equation}
 \frac1{1-q}=\frac R\eta
     =1+\frac{4c^2}{\rho^2}+\frac{2c}{\rho}.
 \label{alg:contraction-parameter}
\end{equation}
Since \(-\log q\ge1-q\), the stopping rule implies
\[
 p=O\!\left(
   \frac1{1-q}
   \log\!\left(2+\frac{\mu M}{\delta(1-q)}\right)
 \right).
 \]
The bound \eqref{alg:order-bound} follows, because
\(M=O_{d,J,\mu}(n)\).
For fixed \(d,J,\mu\), this is
\(p=O_{d,J,\mu}(\log(2+n/\delta))\).
The coefficient algorithm uses
\(N\exp(O_D(p))\) arithmetic operations.
Interpolation, the triangular transform, and evaluation in
\eqref{alg:output} add polynomially many operations in \(p\).

Appendix~\ref{alg:energy-bit-details} proves that the exact coefficient
computation, interpolation, and evaluation have polynomial bit overhead.
Together with \eqref{alg:order-bound}, this completes the proof.
\end{proof}

\subsection{Suzuki--Fisher ground energies}
\label{alg:sf-subsection}

The preceding result applies to the Suzuki--Fisher class of
Theorem~\ref{eq:sf}.  We keep its sign convention:
the longitudinal term is
\(-\sum_i\mu_i^z Z_i\), with \(\mu_i^z\ge0\).
The Lee--Yang theorem gives disk stability preservation by the
Gibbs semigroup \(e^{-tH}\); see
\cite{SuzukiFisher71,WBG26,BGLW26}.

\begin{theorem}[Fixed positive longitudinal field]
\label{alg:sf}
Fix a maximum degree \(d\), a positive rational bound \(J\) on
the absolute values of the Pauli coefficients, and a positive
rational \(\mu\).
Let \(H\) be a Suzuki--Fisher Hamiltonian with rational Pauli
coefficients on an \(n\)-vertex graph of maximum degree at most
\(d\), and suppose
\[
 \mu_i^z\ge\mu\qquad(1\le i\le n).
 \]
For every positive rational \(\delta\), its ground energy
\(E_0(H)\) can be approximated to absolute error at most
\(\delta\) by a deterministic classical algorithm polynomial in
\(n\), the input bit length, and \(\delta^{-1}\).
The polynomial may depend on \(d,J,\mu\).
\end{theorem}

\begin{proof}
Define
\[
 V=H+\mu\sum_i Z_i.
 \]
The longitudinal fields of \(V\) are \(\mu_i^z-\mu\ge0\);
all its interaction inequalities are unchanged.
Thus \(V\) is again Suzuki--Fisher, and
\(A_0=-V\) satisfies \eqref{alg:promise}.
Its local matrices have Gaussian-rational entries.
Combining the finitely many Pauli terms on each edge or site
gives a local norm bound depending only on \(J,\mu\), and leaves
the degree bound unchanged.
For example, sums of the absolute Pauli coefficients give
explicit rational norm upper bounds, so the algorithm need not
compute exact matrix norms.

Since \(2E=n\Id-\sum_i Z_i\),
\begin{equation}
 A_0-2\mu E=-H-\mu n\Id.
 \label{alg:sf-shift}
\end{equation}
The matrix on the right is Hermitian.  Its eigenvalue with
largest real part is therefore
\(-E_0(H)-\mu n\).
Apply Theorem~\ref{alg:complex} with error \(\delta\).
If it returns \(\widehat a\), return the real rational
\[
 \widehat E=-\Re\widehat a-\mu n.
 \]
Equation \eqref{alg:sf-shift} gives the required error bound
and the asserted complexity.
\end{proof}

\begin{remark}[Hermitian coefficient bound]
\label{alg:hermitian-radius}
For this application, the complexification step can be omitted:
the perturbation \(V\) is Hermitian.  Applying
Proposition~\ref{alg:bdl} directly gives the larger guaranteed
disk \(2^{-16}/(DJ')\), where \(J'\) is the local norm bound
after the auxiliary construction.  The same continuation proof
then applies to the ground-energy branch of \(2E+xV\).
The factor-four improvement in this disk changes constants,
not the fixed-parameter complexity conclusion.
\end{remark}

Theorem~\ref{alg:sf} includes nonstoquastic interactions and
does not require conservation of total degree.  Its computational
output and parameter dependence differ from several earlier
results.  Bravyi and Gosset \cite{BG17} give randomized
algorithms for a two-axis stoquastic family on arbitrary graphs.
Harrow, Mehraban, and Soleimanifar \cite[Section~7 of the full version]{HMS20}
give an all-temperature quasipolynomial partition-function
algorithm for a number-conserving XXZ subclass.  Their
fugacity-polynomial reduction uses the commutation of the
interaction with total magnetization.
Bravyi, Gosset, Liu, and Wong \cite{BGLW26} give a polynomial-time
quantum ground-energy algorithm for the full Suzuki--Fisher
family, including zero fields and inverse-polynomial absolute
accuracy.  Their sharp Hermitian gap and quantum algorithm are
not consequences newly claimed here.
Theorem~\ref{alg:sf} gives a deterministic classical conclusion
on bounded-degree inputs at every fixed positive field.
It does not replace the unrestricted quantum theorem.

\subsection{Energy-value approximation at zero field}
\label{alg:value-subsection}

A fixed positive field can also be used as a controlled
perturbation when the requested error is proportional to the
number of sites.

\begin{theorem}[Additive energy-value approximation]
\label{alg:value}
Fix a maximum degree \(d\), a bound \(J\) on the absolute
Pauli coefficients, and a positive rational \(\varepsilon\).
For every rational Suzuki--Fisher Hamiltonian on an
\(n\)-vertex graph satisfying these bounds, a deterministic
classical algorithm polynomial in \(n\) and the input bit length
returns \(\widehat E\in\mathbb Q\) such that
\[
 |\widehat E-E_0(H)|\le\varepsilon n.
 \]
The longitudinal fields may vanish.
The polynomial may depend on \(d,J,\varepsilon\).
The output is an energy value; no approximating state is
asserted.
\end{theorem}

\begin{proof}
It suffices to consider \(0<\varepsilon\le1\), since a smaller
accuracy parameter also meets a larger requested tolerance.
Put \(\nu=\varepsilon/4\) and
\[
 H_\nu=H-\nu\sum_i Z_i.
 \]
This is a Suzuki--Fisher Hamiltonian whose longitudinal fields
are all at least the fixed positive number \(\nu\).
Its local coefficients are bounded in terms of \(J,\varepsilon\).
The variational characterization of the lowest eigenvalue gives
\[
 |E_0(H_\nu)-E_0(H)|
 \le\|H_\nu-H\|\le\nu n.
 \]
Use Theorem~\ref{alg:sf} to estimate \(E_0(H_\nu)\) with
absolute error at most \(\varepsilon n/2\).
The returned value differs from \(E_0(H)\) by at most
\[
 \frac{\varepsilon n}{4}+\frac{\varepsilon n}{2}
   =\frac{3\varepsilon n}{4}
   \le\varepsilon n.
 \]
For fixed \(d,J,\varepsilon\), the field \(\nu\) and all local
bounds are fixed.  The runtime in Theorem~\ref{alg:sf} is
therefore polynomial in the input size and \(n\).
\end{proof}

We spell out the consequence for Quantum MaxCut and EPR
Hamiltonians, keeping value approximation separate from state
construction.
For a graph \(G=(V,\mathcal E)\) with nonnegative rational weights, let
\[
 Q_G=\sum_{\{i,j\}\in\mathcal E}
      w_{ij}|\Psi^-\rangle\langle\Psi^-|_{ij},
 \qquad
 |\Psi^-\rangle=\frac{|01\rangle-|10\rangle}{\sqrt2},
 \]
and write \(\operatorname{OPT}(G)=\lambda_{\max}(Q_G)\).
Similarly, define the EPR objective
\[
 P_G=\sum_{\{i,j\}\in\mathcal E}
      w_{ij}|\Phi^+\rangle\langle\Phi^+|_{ij},
 \qquad
 |\Phi^+\rangle=\frac{|00\rangle+|11\rangle}{\sqrt2}.
 \]

\begin{corollary}[EPR and bipartite Quantum MaxCut values]
\label{alg:qmc}
For fixed degree and upper weight bounds and fixed rational
\(\varepsilon>0\), the optimal values of \(P_G\) on arbitrary
graphs and of \(Q_G\) on bipartite graphs admit deterministic
classical additive-\(\varepsilon |V|\) approximation in
polynomial time.
For both unweighted classes and each fixed rational \(0<\varepsilon<1\),
it returns a rational value \(L\) satisfying
\[
 (1-\varepsilon)\operatorname{OPT}\le L\le\operatorname{OPT},
\]
where \(\operatorname{OPT}\) denotes the largest eigenvalue of the
respective objective. The output is a value, rather than a state
attaining it.
\end{corollary}

\begin{proof}
The negative EPR objective is Suzuki--Fisher up to a scalar:
\begin{equation}
 -P_G
  =-\frac14\sum_{\{i,j\}\in\mathcal E}w_{ij}
       \bigl(\Id+X_iX_j-Y_iY_j+Z_iZ_j\bigr).
 \label{alg:epr-expansion}
\end{equation}
On each edge the transverse matrix has operator norm
\(w_{ij}/4\), equal to its longitudinal interaction
coefficient.  Subtract the known scalar
\(-\frac14\sum_{\{i,j\}}w_{ij}\) when applying
Theorem~\ref{alg:value}, and restore it afterward.
This proves the assertion for \(P_G\).

If \(G\) is bipartite with parts \(A,B\), let
\(U=\prod_{i\in A}Y_i\).
On every edge, \(U\) maps the singlet vector to the EPR vector
up to a scalar of modulus one.  Hence
\[
 UQ_GU^*=P_G.
 \]
The spectra agree, proving the additive assertion for
\(Q_G\).  This is the same local change of basis used in
\cite[Corollary~2]{BGLW26}.

For the unweighted assertions, remove isolated vertices and let
\(N,m\) be the remaining vertex and edge counts. If \(m=0\), return
zero. Otherwise \(N\le2m\). For EPR, the product vector \(|0^N\rangle\)
has expectation \(m/2\). For bipartite Quantum MaxCut, the product
vector which is \(0\) on one part and \(1\) on the other has the same
expectation. In either case \(\operatorname{OPT}\ge m/2\ge N/4\).
Run the additive algorithm with error
\(\eta=\varepsilon N/8\), obtaining \(\widehat v\), and return
\(L=\max(0,\widehat v-\eta)\). Then
\[
 L\le\operatorname{OPT},\qquad
 L\ge\operatorname{OPT}-2\eta
   \ge(1-\varepsilon)\operatorname{OPT}.
\]
For each fixed \(\varepsilon\), the running time is polynomial.
\end{proof}

The additive statement for bounded weights does not automatically
give a relative guarantee when positive weights may be arbitrarily
small.  Nor is the value scheme a state-producing approximation
algorithm of the type studied, for example, in \cite{ALMPSS26}.
The inverse-polynomial absolute-accuracy conjectures for general
EPR problems in \cite{MarwahaSud26} remain outside its scope.
Indeed, introducing a field of order \(\delta/n\) would make
the factor \(c/\rho\) in \eqref{alg:order-bound} grow with \(n\).
The resulting coefficient count need no longer be logarithmic,
and the algorithm proved here need no longer run in polynomial
time.


\subsection{Stable tests of principal eigenvectors}
\label{alg:states-subsection}

We now pass from one eigenvalue to queries about its principal
eigenvectors.  The input remains \eqref{alg:local-input}--\eqref{alg:promise}.
For the binary coefficient space, the transpose exchanges the two
groups of variables in a matrix kernel.  Proposition~\ref{spec:perron}
therefore gives strictly disk-stable right eigenvectors of both
\(A=A_0-2\mu E\) and \(A^{\mathsf T}\).  Normalize them by
\[
 v_\varnothing=w_\varnothing=1.
\]
The transpose, rather than the adjoint, makes the dependence on a
complex parameter holomorphic.  A \emph{stable product test} is
\(B=\bigotimes_{i=1}^n B_i\), where \(B_i\in M_2(\C)\),
\((B_i)_{00}=1\), and
\[
 K_{B_i}(z,w)=
 (B_i)_{00}+(B_i)_{10}z+(B_i)_{01}w+(B_i)_{11}zw
\]
does not vanish in \(\D^2\).  The matrices \(B_i\) need not be
Hermitian, positive, or invertible.

\begin{theorem}[Classical product tests of principal states]
\label{alg:product-test}
Fix \(d,J,\mu\) as in Theorem~\ref{alg:complex}, and suppose that
\(A_0\) satisfies that theorem's input assumptions.  Given a stable
product test with Gaussian-rational entries, both
\[
 G_B=w^{\mathsf T}Bv,\qquad
 \Tr(B\Pi)=\frac{w^{\mathsf T}Bv}{w^{\mathsf T}v},
 \qquad \Pi=\frac{vw^{\mathsf T}}{w^{\mathsf T}v},
\]
are nonzero.  For every rational \(0<\epsilon<1\), a deterministic
classical algorithm returns Gaussian-rational estimates of these
two quantities with relative complex error at most \(\epsilon\).
Its bit complexity is
\[
 \operatorname{poly}_{d,J,\mu}
       (n,\epsilon^{-1},L_{\mathrm{in}}),
\]
where \(L_{\mathrm{in}}\) includes the local generator, the test
matrices, and the accuracy.  Relative complex error means
\(|\widehat z/z-1|\le\epsilon\).
\end{theorem}

This is query access to specified scalar tests, not a listing of an
exponentially long vector.  No polynomial bound on \(\|\Pi\|\) is
assumed or obtained.  The proof instead constructs a nonvanishing
scalar function with a controlled logarithm in the reciprocal-field
parameter.

\subsection{A uniform domain for normalized eigenvectors}

Write \(V=-A_0\), \(H(x)=2E+xV\), and set
\begin{equation}\label{alg:state-radii}
 r_0=\frac{2^{-18}}{(d+1)J},\qquad \rho_s=\frac{r_0}{8}.
\end{equation}
The following use of the Kirkwood--Thomas expansion concerns its
creation coefficients, not just the energy estimate of
Proposition~\ref{alg:bdl}.
For a nonempty set \(M\subseteq[n]\), put
\[
 a_M^\dagger=\prod_{i\in M}|1\rangle\langle0|_i .
\]
These operators commute, and products with overlapping supports vanish.
The vacuum-normalized eigenvector germ has a unique expression
\[
 v(x)=\exp\left(-\sum_{\varnothing\ne M\subseteq[n]}
                   C(M;x)a_M^\dagger\right)\Omega,
 \qquad C(M;x)=\sum_{k\ge1}C_k(M)x^k.
\]
Uniqueness follows by taking the finite logarithm in the nilpotent
commutative algebra generated by the \(a_i^\dagger\).

\begin{lemma}[Creation coefficients and a zero-free vacuum disk]
\label{alg:creation-disk}
For every local complex \(V\) with the bounds above,
\begin{equation}\label{alg:creation-bound}
 \max_i\sum_{M\ni i}|C_k(M)|\le2^{-15}r_0^{-k},
 \qquad C_k(M)=0\quad\hbox{if }|M|>k+1.
\end{equation}
The normalized right eigenvector germs of \(H(x)\) and
\(H(x)^{\mathsf T}\) extend holomorphically to \(|x|<\rho_s\).
Their coefficient polynomials are nonzero on the closed polydisk
\(\{|z_i|\le2\}\) throughout this parameter disk.
\end{lemma}

\begin{proof}
For Hermitian perturbations with local norm bound \(L\) on a graph
of maximum degree \(D\), Lemma~5, equation~(32), and Lemma~7,
equation~(47), of \cite{BDL08}, with unperturbed gap \(2\), give
\[
 \max_i\sum_{M\ni i}|C_k(M)|\le2^{-15}R_L^{-k},
 \qquad R_L=2^{-16}/(DL),
\]
and say that a nonzero \(C_k(M)\) is supported inside a connected
set of at most \(k+1\) vertices.  In particular \(|M|\le k+1\).
Use the private auxiliary qubits in \eqref{alg:auxiliary} for
one-site terms, with \(D=d+1\).  The normalized eigenvector germ
of the augmented operator is exactly the physical germ tensored
with the auxiliary vacuum.  Its creation coefficients involving
an auxiliary vertex vanish; all other coefficients agree with
the physical ones.

The recursion \eqref{alg:formal-energy}--\eqref{alg:formal-vector}
and the finite logarithm show that \(C_k(M)\) is a homogeneous
polynomial of degree \(k\) in the local entries of \(V\), without
their conjugates.  With \(W_\theta\) from the proof of
Lemma~\ref{alg:complexification},
\[
 C_k(M;V)=\frac{2^k}{2\pi}\int_0^{2\pi}
              e^{ik\theta}C_k(M;W_\theta)\,d\theta.
\]
Taking absolute values, summing over \(M\ni i\), and applying the
Hermitian estimate with \(L=2J\) proves
\eqref{alg:creation-bound}.  The support conclusion is inherited
by the same coefficient identity.  Transposition gives the
identical estimates for the left germ.
For \(|x|<r_0/8\), it follows that
\begin{equation}\label{alg:weighted-creation}
 \max_i\sum_{M\ni i}4^{|M|}|C(M;x)|
 \le 2^{-15}\sum_{k\ge1}4^{k+1}(|x|/r_0)^k
 <2^{-13}.
\end{equation}
All these series converge locally uniformly.

We give the elementary polymer argument that turns this estimate
into zero-freeness.  For weights \(u_M\) on nonempty subsets of a
finite set, define
\[
 Z(S)=\sum_{\substack{\mathcal A\text{ a collection of}\\
              \text{pairwise disjoint nonempty subsets of }S}}
                       \prod_{M\in\mathcal A}u_M .
\]
Suppose that
\(\sup_i\sum_{M\ni i}|u_M|2^{|M|-1}\le1/2\).
Induction on \(|S|\) proves
\begin{equation}\label{alg:polymer-ratio}
 Z(S)\ne0,\qquad
 \left|\frac{Z(S)}{Z(S\setminus\{i\})}-1\right|\le\frac12
 \quad(i\in S).
\end{equation}
Indeed,
\[
 Z(S)=Z(S\setminus\{i\})+
           \sum_{\substack{M\subseteq S\\i\in M}}u_M Z(S\setminus M).
\]
By the induction hypothesis, deleting the \(|M|-1\) additional
vertices gives
\(\left|Z(S\setminus M)/Z(S\setminus\{i\})\right|
 \le2^{|M|-1}\).
The claimed ratio estimate follows, and its disk does not contain
zero.  The empty set starts the induction.

The polynomial of \(v(x)\) is this polymer sum with
\(u_M=-C(M;x)z^M\).  The factorials in the exponential cancel
the orderings of pairwise disjoint sets.  For \(|z_i|\le2\),
the left side of the polymer condition is at most half the
left side of \eqref{alg:weighted-creation}, hence is less than
\(1/2\).  This proves the asserted zero-freeness.

The convergent creation series and its finite exponential give
holomorphic vectors on \(|x|<\rho_s\).  Their formal eigenvalue
is the convergent energy series of
Lemma~\ref{alg:complexification}.  The eigenvector equations hold
near zero and then throughout the disk by the identity theorem.
No assertion of simplicity on the entire disk is needed.
\end{proof}

\begin{lemma}[Nonvanishing principal tests]\label{alg:test-domain}
The vectors in Lemma~\ref{alg:creation-disk} and the normalized
principal vectors of \(A(1/x)\) glue to holomorphic vectors
\(v(x),w(x)\) on
\[
 \Omega_{\rho_s}=\{\Re x>0\}\cup\{|x|<\rho_s\}.
\]
For every stable product test,
\begin{equation}\label{alg:test-function}
 G_B(x)=w(x)^{\mathsf T}Bv(x)
\end{equation}
is nonvanishing on this domain, \(G_B(0)=1\), and
\begin{equation}\label{alg:test-log-bound}
 |G_B(x)|\le4^n,\qquad
 \Re F_B(x)\le n\log4,\quad F_B=\log G_B,\quad F_B(0)=0 .
\end{equation}
\end{lemma}

\begin{proof}
The half-plane argument in Lemma~\ref{alg:half-plane} applies to
the principal eigenvectors as well.  Local analytic eigenvectors
exist for a simple eigenvalue; their vacuum coordinates are nonzero
because their coefficient polynomials are strictly stable.
Vacuum normalization makes these local choices agree on overlaps.
Use \(A(h)^{\mathsf T}\) for the left vectors.
On a sufficiently small positive real interval they coincide
with the vacuum vectors, by the same Riesz-projection argument
as in Lemma~\ref{alg:half-plane}.  That interval may depend on \(n\);
it is used only to identify the branches, not to choose
\(\rho_s\).  The identity theorem on the connected overlap
gives the required gluing.

In the half-plane, \(v(x)\) and \(w(x)\) each have a strict
polydisk radius greater than one.  In the small disk they have
radius at least two by Lemma~\ref{alg:creation-disk}.
Contract their legs against the two legs of every \(K_{B_i}\).
The strict version of the disk contraction rule, equivalently
polarized apolarity, applies because each contracted pair has
product of radii greater than one; see
Lemma~\ref{pre:grace}.
The resulting scalar is \eqref{alg:test-function}, hence is
nonzero.  This argument allows singular tests and needs no
uniform lower bound for the strict radius.

For completeness, a multiaffine stable polynomial
\(f(z)=\sum_S f_Sz^S\) with \(f_\varnothing=1\) satisfies
\(|f_S|\le1\).  Set variables outside \(S\) to zero and all
remaining variables to \(t\).  If \(f_S\ne0\), the resulting
polynomial has degree \(|S|\), constant term one, leading
coefficient \(f_S\), and all roots of modulus at least one.
Their product proves the bound.  The case \(f_S=0\) is immediate.
Apply this observation to \(v,w\) and each \(K_{B_i}\).
The expansion of \(w^{\mathsf T}Bv\) has \(4^n\) terms of modulus
at most one.  Also \(G_B(0)=\prod_i(B_i)_{00}=1\).
The domain is star-shaped about zero, so its nonvanishing
holomorphic function has a unique logarithm vanishing there.
Taking its real part gives \eqref{alg:test-log-bound}.
\end{proof}

\begin{lemma}[Continuation with a one-sided logarithmic bound]
\label{alg:log-continuation}
Use the parameters \eqref{alg:map-parameters} with
\(\rho=\rho_s\) and \(c=1/\mu\).  Write
\(F_B(x)=\sum_{j\ge1} f_jx^j\) and
\(F_B(\phi(u))=\sum_{k\ge1}g_ku^k\).
The transform \eqref{alg:coefficient-transform} remains valid,
and
\begin{equation}\label{alg:log-tail}
 |g_k|\le4n,\qquad
 \left|F_B(c)-\sum_{k=1}^p g_kq^k\right|
       \le \frac{4nq^{p+1}}{1-q}.
\end{equation}
In particular \(|F_B(c)|\le Kn\), where \(K=4q/(1-q)\).
\end{lemma}

\begin{proof}
The map lemma puts \(\phi(\D)\) in \(\Omega_{\rho_s}\).
The holomorphic function
\(U(u)=2n-F_B(\phi(u))\) has nonnegative real part and \(U(0)=2n\).
On a circle of radius \(r<1\), the positive function
\(\Re U(re^{it})\) has average \(2n\).
Its \(k\)-th Fourier coefficient, for \(k\ge1\), equals
\(-g_kr^k/2\), so \(|g_k|r^k\le4n\).
Let \(r\uparrow1\) and sum the geometric tail.
The coefficient transform is a formal identity already proved
in Lemma~\ref{alg:map}.
\end{proof}

\subsection{Connected coefficients and finite bit complexity}

We compute the coefficients of \(F_B\) without expanding the
full eigenvectors.  This is an instance of the connected-logarithm
method for multiplicative analytic graph functions; compare
\cite{PatelRegts17,YYZ22}.  The needed multiplicativity is formal
at the vacuum and does not require a preservation promise on
induced subgraphs.

For \(U\subseteq[n]\), keep just the local terms supported inside
\(U\), obtaining \(V_U\), and let \(B_U=\bigotimes_{i\in U}B_i\).
Define the vacuum-normalized formal vectors \(v_U,w_U\) for
\(2E_U+xV_U\) and its transpose.  Put
\[
 F_{B,U}(x)=\log(w_U(x)^{\mathsf T}B_Uv_U(x)),\qquad
 \Phi_{B,U}(x)=\sum_{W\subseteq U}(-1)^{|U|-|W|}F_{B,W}(x).
\]
For the empty subsystem all logarithms are zero.

\begin{lemma}[Local computation of test logarithms]
\label{alg:connected-tests}
For \(k\ge1\),
\[
 [x^k]F_B(x)=
 \sum_{\substack{U\subseteq[n]\ \mathrm{connected}\\ |U|\le k+1}}
                       [x^k]\Phi_{B,U}(x).
\]
For Gaussian-rational input, the first \(p\) coefficients can
be computed exactly with
\(n\exp(O_d(p))\operatorname{poly}(L_{\mathrm{in}},n,p)\)
bit operations.
\end{lemma}

\begin{proof}
If the induced graph on \(U\) splits into components, the
normalized formal vectors tensor across components and the
test does likewise.  Thus \(F_{B,U}\) is the sum of the
component logarithms.  Subset inversion shows that
\(\Phi_{B,U}=0\) when \(U\) is disconnected.

To obtain the order bound as well, give each local perturbation
term its own parameter.  An order-\(k\) coefficient is a
homogeneous polynomial of degree \(k\) in these parameters,
by the unique formal eigenvector recursion.  If a parameter
monomial has disconnected union of supports, set all unused
parameters to zero.  Tensor factorization makes the logarithm
additive across those components, so that monomial has
coefficient zero.  A connected union of at most \(k\) one- or
two-site supports has at most \(k+1\) vertices.
Finally the coefficient of a monomial with support union \(T\)
is unchanged in every induced subsystem containing \(T\).
The alternating sum over \(W\subseteq U\) retains it only if
\(T=U\).  This proves the formula, including both vanishings.

Here is a direct finite-input implementation.  In each subsystem
of size \(m\le p+1\), use the exact recursions
\eqref{alg:formal-energy}--\eqref{alg:formal-vector}, and their
transposes, through order \(p\).
If
\[
 a_k=\sum_{i+j=k}w_{U,i}^{\mathsf T}B_Uv_{U,j},
 \qquad \log\left(1+\sum_{k\ge1}a_kx^k\right)
                  =\sum_{k\ge1}f_kx^k,
\]
then coefficient comparison after logarithmic differentiation gives
\begin{equation}\label{alg:formal-test-log}
 f_k=a_k-\frac1k\sum_{j=1}^{k-1}j f_j a_{k-j}.
\end{equation}
Every vector has dimension at most \(2^{p+1}\).
Even dense matrix arithmetic and enumeration of every
\(W\subseteq U\) cost only \(\exp(O(p))\operatorname{poly}(p)\)
operations per \(U\).  Connected subsets of size at most \(p+1\)
can be enumerated by traversals of their spanning trees, with
deduplication, in \(n\exp(O_d(p))\) time; a traversal has length
at most \(2p\).  The case \(d=0\) consists just of singletons.

Appendix~\ref{alg:test-bit-details} bounds all denominators and
intermediate numerators by polynomial bit lengths, proving the
asserted bit complexity.
\end{proof}

\begin{proof}[Proof of Theorem~\ref{alg:product-test}]
Use \(c=1/\mu\), \(\rho_s\) from \eqref{alg:state-radii}, and the
rational map parameters of Lemma~\ref{alg:map}.
Given \(0<\tau<1\), choose the least \(p\ge0\) for which
\[
 4nq^{p+1}\le\tau(1-q).
\]
Rational multiplication and comparison find \(p\);
\eqref{alg:contraction-parameter} gives
\(p=O_{d,J,\mu}(\log(2+n/\tau))\).
Use Lemma~\ref{alg:connected-tests}, the triangular transform,
and \eqref{alg:log-tail} to compute a Gaussian rational \(S_B\)
with \(|S_B-F_B(c)|\le\tau\).  Also
\(|S_B|\le Kn+1\).  The preceding lemma and the rational
transform give polynomial bit cost.

We make explicit that exponentiation does not introduce an
unbounded-precision oracle.  Let \(T_0\) be an integer upper
bound on \(Kn+1\), and choose an integer \(\ell\) with
\(2^{-\ell}\le\tau\).
Truncate the Taylor series of \(e^{S_B}\) at degree
\[
 M=\left\lceil10(2T_0+\ell+1)\right\rceil
\]
by taking \(\sum_{j=0}^{M}S_B^j/j!\).  The remainder is at most
\(e^{T_0}T_0^{M+1}/(M+1)!\).
Using \(r!\ge(r/3)^r\), it is less than
\(\tau e^{-T_0}\).
Since \(|e^{S_B}|\ge e^{-T_0}\), this is relative error
less than \(\tau\).  All terms are Gaussian rational,
\(M=O_{d,J,\mu}(n+\log\tau^{-1})\), and the numerators and
denominators have polynomial bit length.
Furthermore
\(|e^{S_B}/G_B(c)-1|\le e^\tau-1\).
Choosing \(\tau=\epsilon/32\) supplies the required margin.
Repeat for \(B=I\), and divide the two nonzero rational
estimates to obtain the estimate of \(G_B/G_I\).
For errors at most \(\epsilon/8\) in each factor, the ratio
error is at most
\((\epsilon/4)/(1-\epsilon/8)<\epsilon\).

The nonvanishing assertions follow from
Lemma~\ref{alg:test-domain}; \(I\) is a stable product test
because its local kernel is \(1+zw\).
For a simple eigenvalue the spectral projection is
\(vw^{\mathsf T}/(w^{\mathsf T}v)\), so this quotient is
indeed \(\Tr(B\Pi)\).
\end{proof}

\subsection{Product-state queries and equatorial measurement sampling}
\label{alg:measurements-subsection}

For \(s\in\C\) with \(|s|\le1\), define
\[
 |s\rangle=\frac{|0\rangle+s|1\rangle}{\sqrt{1+|s|^2}} .
\]
These are the pure states in the closed northern hemisphere
of the Bloch sphere in our fixed basis.  The amplitude and
probability queries extend to this whole hemisphere.  Complete
two-outcome projective measurement sampling is stated on the
equator, where the two states \(|s\rangle,|-s\rangle\) are
orthogonal and both satisfy the condition.

\begin{theorem}[Classical access to Suzuki--Fisher ground states]
\label{alg:sf-state}
Fix \(d,J,\mu\) as in Theorem~\ref{alg:sf}.
Let \(H\) be a rational Suzuki--Fisher Hamiltonian on \(n\ge1\)
qubits satisfying that theorem's hypotheses, and let
\(0<\epsilon<1\) be rational.
Its unique unit ground vector \(\psi\) can be normalized by
\(\langle0^n|\psi\rangle>0\).  The query algorithms in (1)--(2)
have bit complexity
\(\operatorname{poly}_{d,J,\mu}(n,\epsilon^{-1},L_{\mathrm{in}})\).
\begin{enumerate}
\item Given \(s_1,\ldots,s_n\in\mathbb Q(i)\) with
      \(|s_i|\le1\), approximate the nonzero complex amplitude
      \(\langle s_1\cdots s_n|\psi\rangle\) with relative
      complex error at most \(\epsilon\).
\item Given \(S\subseteq[n]\) and \(s_i\in\mathbb Q(i)\),
      \(|s_i|\le1\) for \(i\in S\), approximate the positive
      probability
      \[
      \left\langle\psi\left|
        \bigotimes_{i\in S}|s_i\rangle\langle s_i|
                \otimes I_{S^c}\right|\psi\right\rangle
      \]
      by a positive rational with relative error at most
      \(\epsilon\).
\item Given a polynomial-time adaptive strategy that measures
      each site at most once, in a basis
      \(\{|s\rangle,|-s\rangle\}\) with
      \(s\in\mathbb Q(i)\), \(|s|=1\), sample its outcome
      record by a classical randomized algorithm with total
      variation error at most \(\epsilon\).
\end{enumerate}
The strategy in (3) has a finite description, takes previous
outcomes as input, and outputs the next unmeasured site and a
phase.  More precisely, if its total evaluation time along
any record is at most \(T_{\rm strategy}\), and all generated
phases have bit length at most \(L_{\rm phase}\), the sampler
uses at most
\[
 T_{\rm strategy}+
 \operatorname{poly}_{d,J,\mu}
       (n,\epsilon^{-1},L_{\rm in}+L_{\rm phase})
\]
bit operations.  In particular the sampler is polynomial-time
for each fixed polynomial-time strategy.
\end{theorem}

\begin{proof}
Use \(A_0=-(H+\mu\sum_iZ_i)\) as in
\eqref{alg:sf-shift}.  Its local norm bound is a fixed rational
\(J'\) depending on the Pauli bound \(J\) and on \(\mu\);
all radii in the preceding lemmas are evaluated with this
bound, not with an unadjusted \(J\).
At positive real \(x\), \(H(x)\) is Hermitian, so its
vacuum-normalized transpose eigenvector is \(w=\bar v\).
Consequently \(G_I=\|v\|^2>0\) and
\(\psi=v/\sqrt{G_I}\).  Along the positive real path the
normalized logarithm \(F_I\) is the real logarithm, because
it starts at zero and \(G_I\) stays positive.

For amplitudes take
\[
 B_i^{\mathrm{amp}}=
 \begin{pmatrix}1&\bar s_i\\0&0\end{pmatrix}.
\]
Its kernel \(1+\bar s_iw\) is disk-stable, and
\[
 \langle s_1\cdots s_n|\psi\rangle
 =\left(\prod_i(1+|s_i|^2)^{-1/2}\right)
       \exp(F_{B^{\mathrm{amp}}}(c)-F_I(c)/2).
\]
For the event in (2), take \(B_i=I\) off \(S\) and
\[
 B_i^{\mathrm{prob}}=
 \begin{pmatrix}1&\bar s_i\\s_i&|s_i|^2\end{pmatrix}
 =(1+|s_i|^2)|s_i\rangle\langle s_i|
 \quad(i\in S).
\]
Its kernel factors as \((1+s_iz)(1+\bar s_iw)\).
Thus
\[
 P_S=\left(\prod_{i\in S}(1+|s_i|^2)^{-1}\right)
                    \exp(F_B(c)-F_I(c)).
\]
Nonvanishing follows from Lemma~\ref{alg:test-domain};
positivity of the projector expectation then gives \(P_S>0\).

Compute the logarithms with a fixed fraction of the requested
error, using the constructive proof of
Theorem~\ref{alg:product-test}.  For (2), take the real part
of the logarithmic difference before exponentiating; its error
does not increase because the exact difference is real on
the positive path.  For either logarithmic difference use
\(T_0=\lceil2Kn+2\rceil\) in that finite Taylor procedure.
In case (2), the resulting rational estimate is real and
positive when its relative error is less than one.
The probability normalization is rational and exact.
For (1), the product \(R_s=\prod_i(1+|s_i|^2)\) is a positive
rational between \(1\) and \(2^n\).  Rational bisection for
\(\sqrt{R_s}\), followed by inversion, computes its reciprocal
square root to any required relative error with
\(O(n+\log\epsilon^{-1}+L_{\mathrm{in}})\) bits and
polynomial arithmetic cost.  Splitting the error budget
between the two logarithms, exponentiation and this factor
proves (1) and (2), including the phase assertion.

For (3), fix an already generated history.  Its two extensions
have true probabilities \(p_+,p_->0\).  Part (2) computes
positive estimates \(\widehat p_\pm\) with relative error at
most \(\eta=\epsilon/(16n)\).
Normalizing their sum gives a conditional probability with
error at most \(2\eta/(1-\eta)<4\eta\), regardless of the
probability of the parent history.  Round it to a dyadic
probability within \(\epsilon/(4n)\) and draw the corresponding
Bernoulli variable using
\(O(\log(n/\epsilon))\) unbiased random bits.
On a fixed history, projectors on distinct sites commute, so
its probability is exactly the product event in (2), even
though the bases were chosen adaptively.
Coupling the two processes at each step bounds their total
variation distance by the sum of the at most \(n\) conditional
errors, which is less than \(\epsilon\).
At most \(2n\) queries are made.  Their accuracy and phase bit
lengths give the displayed polynomial overhead; evaluating the
strategy contributes \(T_{\rm strategy}\).
\end{proof}

The novelty of this conclusion is its uniform local input and its
fixed-parameter polynomial complexity on the full positive-field
Suzuki--Fisher family.  For general Lee--Yang states with a fixed
spatial radius \(r>1\), \cite[Section~4, Theorems~5--6]{WBG26}
already provides state-dependent quasipolynomial circuits for
relative \(X\)-basis amplitudes, with marginal probabilities and phases in
the state-preparation proof; see also \cite{WongThesis26}.
Our proof does not assume a size-independent spatial radius
margin.  It uses the common reciprocal-field domain instead.
Connected coefficient computation and zero-free interpolation
are established methods
\cite{PatelRegts17,YYZ22,WildAlhambra23}; the new input to those
methods is Lemma~\ref{alg:test-domain} for principal-state tests.
In particular \cite[Theorem~15]{YYZ22} already treats tensorized
measurements of thermal partition functions.  Its thermal
zero-free hypothesis does not itself supply the ground-state
parameter domain used here.

The queries above do not include arbitrary product bases:
the projector onto \(|1\rangle\) has no nonzero vacuum entry
and is not an allowed normalized test.
Nor does the fixed-field theorem extend the zero-field
energy-value scheme to state approximation or sampling.
The broad quantum algorithms in \cite{BGLW26,RT26} have
different parameter ranges and outputs; the present classical
result does not subsume them.


\section{All-temperature Lee--Yang Hamiltonians}\label{eq:section}

We now specialize the complex generator classification to Hermitian
operators on qubits.  The resulting inequalities are precisely the
Suzuki--Fisher conditions, including their full transverse operator-norm
form.  Their sufficiency is classical; the generator theorem supplies
necessity and excludes additional interactions of arbitrary support.
Positive longitudinal fields then improve the zero-free domain of the
actual Gibbs tensor, at every positive temperature.

\subsection{From matrix tensors to polynomial generators}

Fix the computational basis \(\{|a\rangle:a\in\{0,1\}^n\}\),
with \(Z|0\rangle=|0\rangle\).  Associate to a matrix \(M\) the
polynomial
\begin{equation}\label{eq:matrix-tensor}
 F_M(z,w)=\sum_{a,b\in\{0,1\}^n}M_{ab}z^aw^b.
\end{equation}
We say that \(M\) has the tensor Lee--Yang property at radius \(r\)
if \(F_M\) is nonzero on \(\D_r^n\times\D_r^n\).  The two
groups of variables in this definition are independent complex variables.
In particular, \eqref{eq:matrix-tensor} is not the sesquilinear form
\(\langle z|M|z\rangle\).

For a Hermitian Hamiltonian \(H\), the Gibbs operator is
\(M_\beta=e^{-\beta H}\), and its normalized density matrix is
\(\rho_\beta=M_\beta/\Tr M_\beta\).  Since
\(\Tr e^{-\beta H}>0\) for real \(\beta\), the two matrices have
exactly the same polynomial zeros.

The following bridge fixes both the tensor and operator conventions.
It is the binary instance, with a unitary scalar normalization, of the
Cayley correspondence \eqref{pre:cayley}.

\begin{lemma}\label{eq:cayley-bridge}
Let
\[
 C_1=\frac1{\sqrt2}\begin{pmatrix}i&-i\\1&1\end{pmatrix},
 \qquad C=C_1^{\otimes n},\qquad D_1=\diag(-1,1).
\]
For an invertible matrix \(M\), its tensor has the unit-disk
Lee--Yang property if and only if \(CMC^{-1}\), acting on
multiaffine coefficient vectors, preserves upper-half-plane stability.
Consequently, the Gibbs family corresponds to the generator
\begin{equation}\label{eq:generator}
 A=-CHC^*.
\end{equation}
Here operator conjugation is a similarity, whereas transforming both
groups of tensor variables gives the congruence \(CMC^T\).
\end{lemma}
\begin{proof}
The coefficient action of \(C\) is
\[
 (Cf)(s)=2^{-n/2}\prod_i(s_i+i)
 f\!\left(\frac{s_1-i}{s_1+i},\ldots,
           \frac{s_n-i}{s_n+i}\right).
\]
The fractions map the upper half-plane onto the unit disk, and the
prefactors are nonzero there.  Thus \(C\) identifies the two
stability classes.  By the binary form of the contraction theorem,
or Lemma~\ref{spec:kernel-criterion}, the disk-stable tensor of an
invertible \(M\) is equivalent to preservation of disk-stable
coefficient vectors.  This proves the equivalence by similarity.

For completeness, the reverse tensor conversion can also be checked
directly.  The unitary matrix \(C_1\) satisfies \(C_1C_1^T=D_1\),
so the coefficient matrix of
\[
 \prod_i(s_it_i-1)
\]
is \(D_1^{\otimes n}=CC^T\).  Each factor is stable in the two
upper-half-plane variables: if \(s_it_i=1\), then
\(t_i=1/s_i\) has negative imaginary part.  Acting on its first
group of variables by \(N=CMC^{-1}\) gives coefficient matrix
\(NCC^T=CMC^T\).  For fixed \(t\in\HH^n\), invertibility
excludes a zero output, so stability preservation gives nonvanishing
jointly in \((s,t)\in\HH^{2n}\).  Finally,
\[
 F_{CMC^T}(s,t)
 =2^{-n}\prod_i(s_i+i)(t_i+i)
 F_M\!\left(\frac{s-i}{s+i},\frac{t-i}{t+i}\right).
\]
This verifies the congruence and the sign of the half-plane kernel.
As \(C^{-1}=C^*\), exponential similarity yields
\(Ce^{-\beta H}C^{-1}=e^{\beta A}\).
\end{proof}

\subsection{The exact Hamiltonian classification}

Write \(X,Y,Z\) for the Pauli matrices, with
\(Y=\left(\begin{smallmatrix}0&-i\\ i&0\end{smallmatrix}\right)\).
For a real \(2\times2\) matrix, \(\|\cdot\|_{\mathrm{op}}\)
denotes its largest singular value.

\begin{theorem}[All-temperature qubit classification]\label{eq:sf}
For a Hermitian operator \(H\) on \((\C^2)^{\otimes n}\), the
following conditions are equivalent:
\begin{enumerate}
\item \(e^{-\beta H}\) has the tensor Lee--Yang property at radius
one for every \(\beta\ge0\).
\item The same property holds for every \(\beta\) in some interval
\([0,\varepsilon)\), where \(\varepsilon>0\).
\item There are real coefficients such that
\begin{align}
 H={}&h_0 I+\sum_i(h_i^xX_i+h_i^yY_i+h_i^zZ_i)\notag\\
 &+\sum_{i<j}\bigl(a_{ij}X_iX_j+b_{ij}X_iY_j
           +c_{ij}Y_iX_j+d_{ij}Y_iY_j+e_{ij}Z_iZ_j\bigr),
 \label{eq:sf-form}
\end{align}
where
\begin{equation}\label{eq:sf-inequality}
 h_i^z\le0,\qquad
 e_{ij}\le-
 \left\|\begin{pmatrix}a_{ij}&b_{ij}\\c_{ij}&d_{ij}\end{pmatrix}
 \right\|_{\mathrm{op}}.
\end{equation}
There is no restriction on \(h_0,h_i^x,h_i^y\).
\end{enumerate}
Equivalently, in the ferromagnetic convention
\begin{equation}\label{eq:ferromagnetic-convention}
 H=h_0I-\sum_i\sum_{a=x,y,z}\mu_i^a\sigma_i^a
 -\sum_{i<j}\left(J_{ij}^{zz}Z_iZ_j+
       \sum_{a,b=x,y}J_{ij}^{ab}\sigma_i^a\sigma_j^b\right),
\end{equation}
the conditions are
\(\mu_i^z\ge0\) and
\(J_{ij}^{zz}\ge\|J_{ij}^{\perp}\|_{\mathrm{op}}\), where the
rows and columns of \(J_{ij}^{\perp}\) are ordered \(x,y\).
\end{theorem}

\begin{proof}
The first condition implies the second.  By
Lemma~\ref{eq:cayley-bridge}, the second says that \(e^{tA}\)
preserves upper-half-plane stability for small \(t\ge0\).
The semigroup identity extends preservation to all \(t\ge0\):
for any \(t\), choose an integer \(m\) with \(t/m<\varepsilon\)
and compose \(e^{(t/m)A}\) \(m\) times.  We may therefore apply
Theorem~\ref{gen:main-complex} on the binary capacity box.

We first spell out how differential order becomes physical support.
On one binary site a basis of endomorphisms is
\[
 I,\qquad J^- =\partial_s,\qquad
 J^0=s\partial_s-\tfrac12,\qquad
 J^+=s^2\partial_s-s.
\]
The last three operators are traceless; their first-order coefficients
are the linearly independent polynomials \(1,s,s^2\).
Expand \(A\) in tensor products of this basis and choose a support
set \(S\) of maximal cardinality occurring with a nonzero coefficient.
The canonical coefficient of \(\partial_S=\prod_{i\in S}\partial_i\)
is then the sum of its exact-support terms, because no strictly larger
support occurs.  Those terms have distinct principal monomials
\(\prod_{i\in S}s_i^{r_i}\), with \(r_i\in\{0,1,2\}\).
The order-two conclusion of Theorem~\ref{gen:main-complex} forces
this coefficient to vanish if \(|S|\ge3\).  Linear independence
then eliminates every term of that support.  Descending in support
eliminates all interactions on three or more sites.  Local similarity
preserves the identity and traceless subspaces at each site, so the
same support conclusion holds for \(H\).

Next write \(A=R+iS\) with entrywise real matrices \(R,S\).
The imaginary-part conclusion of Theorem~\ref{gen:main-complex}
has the form
\begin{equation}\label{eq:imaginary-generator}
 S=c_0I+\sum_iJ_i(q_i),\qquad
 q_i(s)=\alpha_i+\eta_i s+\gamma_i s^2\ge0\quad(s\in\R),
\end{equation}
where \(J(q)=q\partial-q'/2\).  On coefficients \((1,s)\),
\[
 J(q)=\begin{pmatrix}-\eta/2&\alpha\\-\gamma&\eta/2\end{pmatrix}.
\]
Since \(A\) is Hermitian, \(R^T=R\) and \(S^T=-S\).
Taking the trace in \eqref{eq:imaginary-generator} first gives
\(c_0=0\).  Independence of the remaining one-site traceless
components shows that each \(J(q_i)\) is skew-symmetric.  Thus
\(\eta_i=0\) and \(\gamma_i=\alpha_i\); nonnegativity gives
\(\alpha_i\ge0\).  In particular,
\begin{equation}\label{eq:imaginary-local}
 q_i(s)=\alpha_i(1+s^2),\qquad
 iS=-\sum_i\alpha_iY_i.
\end{equation}

The local Pauli conjugations are
\begin{equation}\label{eq:pauli-cayley}
 C_1XC_1^*=-Z,\qquad C_1YC_1^*=X,\qquad C_1ZC_1^*=-Y.
\end{equation}
Therefore the physical field
\(h_i^xX_i+h_i^yY_i+h_i^zZ_i\) contributes
\(h_i^xZ_i-h_i^yX_i+h_i^zY_i\) to \(A\).
Comparison with \eqref{eq:imaginary-local} gives
\(h_i^z=-\alpha_i\le0\).  A two-site physical Pauli term
containing exactly one \(Z\) becomes a term containing exactly one
\(Y\), and is entrywise imaginary.  The four such terms on each
pair are linearly independent.  Since \(S\) has no two-site
component, all physical \(XZ,YZ,ZX,ZY\) coefficients vanish.
This leaves exactly the form \eqref{eq:sf-form}.

It remains to determine the pair inequality.  Suppress the pair indices.
A physical term \(aXX+bXY+cYX+dYY+eZZ\) contributes
\[
 -aZZ+bZX+cXZ-dXX-eYY
\]
to \(A\).  The first-order coefficients of the coefficient operators
\(X,Z,Y\), respectively, are
\[
 p_X(s)=1-s^2,\qquad p_Z(s)=-2s,\qquad
 p_Y(s)=-i(1+s^2).
\]
Hence the full coefficient of \(\partial_i\partial_j\) is
\begin{align}
 q_{ij}(s,t)={}&e(1+s^2)(1+t^2)-4ast
       -2bs(1-t^2)\notag\\
       &-2ct(1-s^2)-d(1-s^2)(1-t^2).
 \label{eq:mixed-coefficient}
\end{align}
There is no factor of two in this expression: it is the full mixed
coefficient in Theorem~\ref{gen:main-complex}.
Define
\[
 u(s)=\left(\frac{2s}{1+s^2},\frac{1-s^2}{1+s^2}\right),
 \qquad J=\begin{pmatrix}a&b\\c&d\end{pmatrix}.
\]
Then
\begin{equation}\label{eq:opnorm-identity}
 \frac{q_{ij}(s,t)}{(1+s^2)(1+t^2)}=e-u(s)^TJu(t).
\end{equation}
The set \(u(\R)\) is dense in the unit circle.  Consequently
the required inequality \(q_{ij}(s,t)\le0\) for all real \(s,t\)
is equivalent to
\[
 e\le\min_{\|u\|_2=\|v\|_2=1}u^TJv=-\|J\|_{\mathrm{op}}.
\]
This proves necessity of \eqref{eq:sf-inequality}.

Conversely, suppose \eqref{eq:sf-form}--\eqref{eq:sf-inequality}
hold.  The displayed conjugations show that \(A\) has canonical
differential order at most two.  Its mixed coefficients are real and
nonpositive by \eqref{eq:opnorm-identity}, and its imaginary part is
\eqref{eq:imaginary-generator} with
\(q_i(s)=-h_i^z(1+s^2)\ge0\).  Its real part therefore satisfies
the real criterion in Theorem~\ref{gen:main-complex}, and its
imaginary part satisfies that theorem's one-site condition.  The
theorem proves preservation for all \(\beta\ge0\); the bridge
then gives the tensor Lee--Yang property of \(e^{-\beta H}\).
Finally, changing the signs of the coefficients converts
\eqref{eq:sf-inequality} into \eqref{eq:ferromagnetic-convention}.
\end{proof}

\begin{remark}\label{eq:sf-attribution}
The sufficient Suzuki--Fisher family originates in
\cite{SuzukiFisher71,AsanoJapanese70,AsanoPRL70}. The full transverse
operator-norm condition is already contained, in equivalent algebraic
form, in \cite[equation~(2.47)]{SuzukiFisher71} and
\cite[Theorem~2, equation~(8)]{Dunlop79}; see also the modern notation
in \cite[Definition~1 and Fact~3]{BGLW26}. Indeed, for a real matrix
$J=\left(\begin{smallmatrix}a&b\\c&d\end{smallmatrix}\right)$,
\[
 2\|J\|_{\rm op}
 =\sqrt{(a-d)^2+(b+c)^2}+\sqrt{(a+d)^2+(b-c)^2}.
\]
To verify the identity, the squared radicals are
$\|J\|_F^2-2\det J$ and $\|J\|_F^2+2\det J$.
If the singular values are $s_1\geq s_2\geq0$, the two radicals
are therefore $s_1+s_2$ and $s_1-s_2$, in some order.
Full matrix generating functions and their
composition are also explicit in Asano's contemporaneous account for
the XXZ subfamily \cite[equations~(1.15), (3.7)--(3.10)]{AsanoJapanese70}.
Thus a distinction between classical traces and modern full tensors
alone would not identify the contribution. Theorem~\ref{eq:sf}
adds the converse among all finite Hermitian qubit Hamiltonians,
including the exclusion of higher-support interactions.  The proof
also shows that testing only a sufficiently small interval of inverse
temperatures is enough.
\end{remark}

\begin{corollary}\label{eq:complex-support}
Let \(H\) be an arbitrary complex qubit matrix.  If
\(e^{-\beta H}\) has the unit-disk tensor Lee--Yang property for
all sufficiently small real \(\beta\ge0\), then its Pauli
expansion contains no term supported on three or more sites.
\end{corollary}
\begin{proof}
The Cayley bridge and the semigroup identity still apply.  The complex
generator theorem bounds the canonical order by two.  The maximal-support
basis argument in the proof of Theorem~\ref{eq:sf} does not use
Hermiticity and therefore gives the assertion.
\end{proof}

\begin{remark}\label{eq:trace-caveat}
For a complex Hamiltonian, tensor stability does not justify trace
normalization.  For example, \(H=\diag(0,i\pi)\) has Gibbs
polynomial \(1+e^{-i\pi\beta}zw\), which is nonzero on the open
unit bidisk for every real \(\beta\), but
\(\Tr e^{-H}=0\).  The support conclusion in
Corollary~\ref{eq:complex-support} applies to this complex setting;
the normalized Gibbs-state statements use Hermiticity.
\end{remark}

\subsection{Strict positive fields and the actual Gibbs kernel}

\begin{theorem}[Strict-field Gibbs radius]\label{eq:strict-field}
Let \(H\) have the form \eqref{eq:ferromagnetic-convention}, with
\(J_{ij}^{zz}\ge\|J_{ij}^{\perp}\|_{\mathrm{op}}\) and
\(\mu_i^z>0\) at every site.  Then, for every \(\beta>0\),
\(F_{e^{-\beta H}}\) is nonzero on the closed unit polydisk in
its full \(2n\) variables.  In particular, it has some radius
\(r_\beta>1\).

The ground eigenspace of \(H\) is one-dimensional.  Its nonzero
coefficient vector \(\psi\) defines a polynomial
\(f_\psi(z)=\sum_a\psi_a z^a\) of radius greater than one.
For every \(\beta_0>0\), one radius \(r>1\) works for all
Gibbs kernels at \(\beta\ge\beta_0\), and also for the image
under each \(e^{-\beta H}\) of every nonzero unit-disk-stable
coefficient vector.
\end{theorem}
\begin{proof}
Let \(\mu=\min_i\mu_i^z>0\), and write
\(H=H_0-\mu\sum_iZ_i\).  The Hamiltonian \(H_0\) satisfies
the nonstrict Suzuki--Fisher conditions.  By Theorem~\ref{eq:sf}
and the disk version of Lemma~\ref{eq:cayley-bridge}, its
coefficient operator \(A_0=-H_0\) generates a disk-stability
preserver.  In the multiaffine coefficient basis,
\(Z_i=I-2E_i\).  Thus
\[
 -H=A_0-2\mu E+\mu nI.
\]
The harmless scalar factor \(e^{\beta\mu n}\) reduces every
kernel and evolution assertion to Theorem~\ref{spec:strict}.

Proposition~\ref{spec:perron} supplies an algebraically simple
dominant eigenvalue of \(-H\).  Since \(H\) is Hermitian,
this is exactly the negative of its ground energy, so the ground
space is one-dimensional.  The same proposition gives the strict
radius of its right eigenvector.  Equivalently, one may fix a
Gibbs kernel radius \(r>1\) and normalize its powers by their
largest eigenvalue.  These powers retain radius \(r\) by
contraction, and converge to \(|\psi\rangle\langle\psi|\).
Their nonzero limiting polynomial is
\[
 f_\psi(z)\overline{f_\psi(\bar w)}.
\]
Hurwitz gives nonvanishing on \(\D_r^{2n}\), and therefore
nonvanishing of each factor on \(\D_r^n\).
\end{proof}

\begin{remark}\label{eq:gap-attribution}
The spectral conclusion
\[
 E_1(H)-E_0(H)\ge2\min_i\mu_i^z
\]
is the Hermitian specialization of Theorem~\ref{spec:gap} and
was already proved for this family in
\cite[Theorem~1]{BGLW26}.  Here the additional conclusion is
strict nonvanishing of the actual Gibbs tensor at every positive
inverse temperature, and its transfer to the ground-state polynomial.
No numerical lower bound for \(r_\beta-1\), uniform in system
size, is asserted.  This distinction also separates the actual
kernel from the strict kernels used in a product-formula approximation.
\end{remark}

\subsection{A qualitative radius for deformed EPR ground states}

Consider a finite graph without isolated vertices, with positive edge
weights \(w_e\).  On an edge \(e=\{i,j\}\), fix
\(0\le s_e<1\) and \(\theta_e\in\R\), and define
\[
 |\eta_e\rangle=|00\rangle+s_e e^{i\theta_e}|11\rangle,
 \qquad H_{\mathrm{EPR}}=-\sum_e w_e
 |\eta_e\rangle\langle\eta_e|_{ij}.
\]
The edge vectors here are unnormalized; replacing them by normalized
vectors simply rescales the positive weights.

\begin{corollary}\label{eq:epr}
The ground space of \(H_{\mathrm{EPR}}\) is one-dimensional, and
its ground-state polynomial is zero-free on a polydisk of some radius
\(r>1\).
\end{corollary}
\begin{proof}
On one edge, direct expansion in Pauli matrices gives
\begin{align*}
 |\eta_e\rangle\langle\eta_e|
 ={}&\frac{1+s_e^2}{4}(I+Z_iZ_j)
    +\frac{1-s_e^2}{4}(Z_i+Z_j)\\
 &+\frac{s_e}{2}\bigl[
      \cos\theta_e(X_iX_j-Y_iY_j)
      +\sin\theta_e(X_iY_j+Y_iX_j)\bigr].
\end{align*}
Its ferromagnetic coefficients are therefore
\[
 J_e^{zz}=\frac{w_e(1+s_e^2)}4,\qquad
 J_e^\perp=\frac{w_es_e}{2}
 \begin{pmatrix}\cos\theta_e&\sin\theta_e\\
                 \sin\theta_e&-\cos\theta_e\end{pmatrix}.
\]
The displayed real matrix is orthogonal, so
\(\|J_e^\perp\|_{\mathrm{op}}=w_es_e/2\) and
\(J_e^{zz}-\|J_e^\perp\|_{\mathrm{op}}
 =w_e(1-s_e)^2/4\ge0\).  The longitudinal field at site \(i\) is
\[
 \mu_i^z=\sum_{e\ni i}\frac{w_e(1-s_e^2)}4>0.
\]
All hypotheses of Theorem~\ref{eq:strict-field} follow, including
strictness because every vertex has an incident edge.  Apply that
theorem to obtain the stated uniqueness and radius.
\end{proof}

For a common deformation parameter \(s<1\), the corollary gives
the qualitative strict-radius conclusion asked for in
\cite[Conjecture~1]{WBG26}.  It does not assert a prescribed sharp
radius as a function of \(s\).  Edge phases may be chosen
independently; no simultaneous gauge removal is used in the proof.


\section{All-temperature Hamiltonians at higher spins}\label{spin:section}

The finite-box generator classification also determines the coherent Gibbs kernels
at arbitrary finite local spins. We retain the binomial normalization throughout.

\subsection{Coherent normalization and the Hamiltonian cone}

At site $i$, fix a positive integer $\kappa_i$ and put $j_i=\kappa_i/2$.
Let $\mathcal H_{\kappa_i}=\C^{\kappa_i+1}$ have orthonormal basis
$|a_i\rangle$, $0\leq a_i\leq\kappa_i$.  For multi-indices $a$ set
\begin{equation}\label{spin:weights}
 b_a=\prod_i\binom{\kappa_i}{a_i},\qquad
 |a\rangle\longleftrightarrow \sqrt{b_a}\,z^a.
\end{equation}
Thus the monomial inner product is
$\langle z^a,z^b\rangle=\delta_{ab}/b_a$.  All adjoints in this section
refer to the physical orthonormal basis, or equivalently to this weighted
polynomial inner product.  The coherent kernel of a matrix $M$ is
\begin{equation}\label{spin:kernel}
 F_M(z,w)=\sum_{a,b}\sqrt{b_a b_b}\,M_{ab}z^a w^b.
\end{equation}
In particular, $F_I(z,w)=\prod_i(1+z_iw_i)^{\kappa_i}$.
We call $M$ coherently Lee--Yang if $F_M$ has no zero when all its variables
belong to $\D$.  Both weights in \eqref{spin:kernel} are essential.

The spin matrices have the following differential realization:
\begin{equation}\label{spin:differential}
 S_i^x=\frac{(1-z_i^2)\partial_i+\kappa_i z_i}{2},\qquad
 S_i^y=\frac{(1+z_i^2)\partial_i-\kappa_i z_i}{2i},\qquad
 S_i^z=\frac{\kappa_i}{2}-z_i\partial_i.
\end{equation}
They satisfy
$(S_i^x)^2+(S_i^y)^2+(S_i^z)^2=j_i(j_i+1)I$.
Write $S_i^\perp=(S_i^x,S_i^y)^{\mathsf T}$.

\begin{theorem}[Higher-spin Hamiltonians]\label{spin:hamiltonian}
For a Hermitian operator $H$ on $\bigotimes_i\mathcal H_{\kappa_i}$, the
following conditions are equivalent:
\begin{enumerate}
\item $e^{-\beta H}$ is coherently Lee--Yang for every $\beta\geq0$;
\item this holds for every $0\leq\beta<\varepsilon$, for some
      $\varepsilon>0$;
\item $H$ admits the expression
\begin{align}\label{spin:hamiltonian-form}
 H={}&h_0 I+\sum_i(h_i^xS_i^x+h_i^yS_i^y+h_i^zS_i^z)
       +\sum_{i:\,\kappa_i\geq2}(S_i^\perp)^{\mathsf T}B_iS_i^\perp
       \notag\\
    &+\sum_{i<j}\bigl((S_i^\perp)^{\mathsf T}J_{ij}S_j^\perp
                         +e_{ij}S_i^zS_j^z\bigr),
\end{align}
where all coefficients are real, $B_i=B_i^{\mathsf T}$, and
\begin{equation}\label{spin:hamiltonian-cone}
 h_i^z\leq0,\qquad B_i\succeq0,\qquad
 e_{ij}\leq-\|J_{ij}\|_{\mathrm{op}}.
\end{equation}
\end{enumerate}
The coefficient $(B_i)_{12}$ multiplies
$S_i^xS_i^y+S_i^yS_i^x$.  At a capacity-one site all single-site
quadrupoles reduce to scalars or zero and are omitted from
\eqref{spin:hamiltonian-form}.
\end{theorem}

We give the normalization and coordinate computations in the proof, so that
the finite-capacity infinitesimal theorem can be applied without changing
the physical adjoint.

\begin{lemma}[Polarization and the coherent Cayley transform]
\label{spin:cayley}
Define the Dicke isometry
\begin{equation}\label{spin:dicke}
 W_i|a\rangle=\binom{\kappa_i}{a}^{-1/2}
                  \sum_{|x|=a}|x\rangle,\qquad W=\bigotimes_iW_i.
\end{equation}
The binary matrix kernel of $WMW^*$ is the normalized polarization of
$F_M$, separately in each row and column block.  Moreover, let
\begin{equation}\label{spin:cayley-matrix}
 C_1=\frac1{\sqrt2}\begin{pmatrix}i&-i\\1&1\end{pmatrix},\qquad
 C_i=W_i^*C_1^{\otimes\kappa_i}W_i,\qquad C=\bigotimes_iC_i.
\end{equation}
For Hermitian $H$, coherent stability of $e^{-tH}$ for all sufficiently
small $t\geq0$ is equivalent to the condition that
$\exp(-tCHC^*)$ is a complete upper-half-plane stability preserver for
every $t\geq0$.
\end{lemma}

\begin{proof}
At binary words of weights $a,b$, the matrix entry of $WMW^*$ is
$M_{ab}/\sqrt{b_ab_b}$.  Identifying the variables in their respective
blocks multiplies this coefficient by $b_ab_b$, giving exactly
\eqref{spin:kernel}.  The disk polarization theorem therefore identifies
stability of the two kernels; see \cite[Theorem~2.1]{BB10II}.
Polarization, binary tensor contraction, and diagonal specialization show
that a coherently stable matrix sends each stable coherent vector to a
stable vector or zero.  An invertible matrix cannot produce the latter
alternative.  This argument does not require $WMW^*$ to be invertible
outside the symmetric subspace.

The matrix $C_i$ is unitary.  On polynomials it acts by
\begin{equation}\label{spin:cayley-action}
 f(z_i)\longmapsto
 \left(\frac{s_i+i}{\sqrt2}\right)^{\kappa_i}
 f\left(\frac{s_i-i}{s_i+i}\right),
\end{equation}
which identifies disk stability with upper-half-plane stability.
Consequently coherent stability of $e^{-tH}$ implies ordinary stability
preservation by $e^{-tCHC^*}$.  Preservation for a sufficiently small
time interval implies preservation for all times by composition.
Theorem~\ref{gen:main-complex} then gives complete preservation.

For the converse, put
$D_i=\diag_{0\leq a\leq\kappa_i}((-1)^{\kappa_i-a})$ and
$D=\bigotimes_iD_i$.  Symmetric powers of $C_1C_1^{\mathsf T}$ give
$CC^{\mathsf T}=D$, and
\begin{equation}\label{spin:universal-halfplane}
 F_D(s,t)=\prod_i(s_it_i-1)^{\kappa_i}
\end{equation}
is upper-half-plane stable.  Indeed $st=1$ is impossible for two points
of the upper half-plane.  A complete preserver
$N=e^{-tCHC^*}=Ce^{-tH}C^*$ applied in the $s$ variables to this kernel
produces the coherent coefficient matrix
$ND=Ce^{-tH}C^{\mathsf T}$.  This is precisely the separate Cayley
transform of both groups of variables in $F_{e^{-tH}}$.  It is nonzero
because $N$ is invertible.  Applying the inverse transforms proves the
required disk stability.
\end{proof}

\begin{proof}[Proof of Theorem~\ref{spin:hamiltonian}]
Put $A=-CHC^*$.  By Lemma~\ref{spin:cayley},
Theorem~\ref{gen:main-complex} applies to $A$.  We recall the parts of its conclusion
used here.  The generator belongs to the degree-at-most-two part of the
enveloping algebra of the local $\mathfrak{sl}_2$ actions.  Its
second-order differential coefficients are real and nonpositive on the
appropriate real lines or planes.  If $A=R+iS$ in the real monomial
basis, then its imaginary part has the form
\begin{equation}\label{spin:imaginary-wedge}
 S=cI+\sum_iJ_{i,\kappa_i}(q_i),\qquad
 J_{i,\kappa_i}(q)=q\partial_i-\frac{\kappa_i}{2}q',\qquad
 q_i\in\R[z_i]_{\leq2},\quad q_i\geq0\text{ on }\R.
\end{equation}
Conversely these conditions, together with preservation of the finite
degree box, characterize the generators.

The degree-two operator space of Proposition~\ref{gen:box-decomposition}
decomposes as the direct sum of the scalar
space, the local spin spaces, the local traceless symmetric quadrupoles
when $\kappa_i\geq2$, and the spaces spanned by
$S_i^aS_j^b$ for $i<j$.  Its dimension is
\begin{equation}\label{spin:degree-two-dimension}
 1+3n+5\#\{i:\kappa_i\geq2\}+9\binom n2.
\end{equation}
For clarity, the local assertion follows from the usual decomposition of
the square of the three-dimensional spin representation into its scalar,
antisymmetric, and traceless symmetric parts: the antisymmetric part is
the spin commutator, the scalar part is the Casimir, and the rank-two
part is nonzero precisely for $\kappa_i\geq2$.  Distinct sites give
independent tensor factors.  This also proves the stated directness.
Since the Cayley action is a local spin rotation, it preserves this
operator space.  Thus an arbitrary Hermitian $H$ satisfying condition (2)
already has total spin degree at most two, with the stated capacity-one
exceptions.

Hermiticity of $A$ and the real diagonal monomial inner product give
$R^*=R$ and $S^*=-S$.  For one site write
\[
 J^-=\partial,\qquad J^0=z\partial-\kappa/2,\qquad
 J^+=z^2\partial-\kappa z.
\]
In the weighted inner product,
$(J^-)^*=-J^+$ and $(J^0)^*=J^0$.
Consequently skew-adjointness of
$aJ^-+bJ^0+cJ^+$, for real $a,b,c$, forces $b=0$ and $a=c$.
The scalar in \eqref{spin:imaginary-wedge} vanishes, and hence
\begin{equation}\label{spin:imaginary-fields}
 q_i(s)=\alpha_i(1+s^2),\quad\alpha_i\geq0,
 \qquad iJ_{i,\kappa_i}(q_i)=-2\alpha_iS_i^y.
\end{equation}
The Cayley rotations are
\begin{equation}\label{spin:rotation}
 CS_i^xC^*=-S_i^z,\qquad CS_i^yC^*=S_i^x,\qquad
 CS_i^zC^*=-S_i^y.
\end{equation}
Thus the original longitudinal field contributes $h_i^zS_i^y$ to $A$,
and \eqref{spin:imaginary-fields} gives $h_i^z=-2\alpha_i\leq0$.
Each mixed longitudinal--transverse quadratic term in $H$ gives an
imaginary second-order coefficient in $A$.  The direct decomposition
above, or the linearly independent principal polynomials at each site,
prevents cancellation between distinct such terms.  They all vanish.

It remains to compute the real second-order inequalities.  The principal
coefficients of $(S^x,S^y,S^z)$ in \eqref{spin:differential} are
$((1-s^2)/2,-i(1+s^2)/2,-s)$.  Put
\[
 u(s)=\left(\frac{2s}{1+s^2},\frac{1-s^2}{1+s^2}\right)^{\mathsf T};
\]
its range is dense in the real unit circle.  For the pair term
$(S_i^\perp)^{\mathsf T}J_{ij}S_j^\perp+e_{ij}S_i^zS_j^z$,
the coefficient $q_{ij}$ of $\partial_i\partial_j$ in $A$ satisfies
\begin{equation}\label{spin:pair-principal}
 \frac{4q_{ij}(s,t)}{(1+s^2)(1+t^2)}
       =e_{ij}-u(s)^{\mathsf T}J_{ij}u(t).
\end{equation}
Its nonpositivity is therefore equivalent to
$e_{ij}\leq-\|J_{ij}\|_{\mathrm{op}}$.
Before fixing the local Casimir gauge, write a surviving local quadratic
term as $(S_i^\perp)^{\mathsf T}K_iS_i^\perp+c_i(S_i^z)^2$, with
$K_i$ real symmetric.  For $\kappa_i\geq2$ its coefficient $q_i$ of
$\partial_i^2$ satisfies
\begin{equation}\label{spin:local-principal}
 \frac{4q_i(s)}{(1+s^2)^2}=c_i-u(s)^{\mathsf T}K_iu(s).
\end{equation}
Thus $c_i\leq\lambda_{\min}(K_i)$.  The Casimir identity rewrites this
term as $(S_i^\perp)^{\mathsf T}B_iS_i^\perp$ plus a scalar, where
$B_i=K_i-c_iI_2\succeq0$.  This proves necessity.

Conversely, \eqref{spin:hamiltonian-cone} makes
\eqref{spin:pair-principal} and \eqref{spin:local-principal} real and
nonpositive and makes the imaginary part satisfy
\eqref{spin:imaginary-wedge}.  The expression preserves the degree box
because it is built from the spin operators.
Theorem~\ref{gen:main-complex} and Lemma~\ref{spin:cayley} give coherent stability at
every time.  This proves (3)$\Rightarrow$(1); (1)$\Rightarrow$(2) is
immediate.
\end{proof}

\begin{corollary}[Exact symmetric-qubit realization]\label{spin:qubit-realization}
Every Hamiltonian in Theorem~\ref{spin:hamiltonian} has a Hermitian qubit
extension $\widehat H$, invariant under permutations within the capacity
blocks, whose pair coefficients satisfy the qubit Lee--Yang cone and such
that
\begin{equation}\label{spin:intertwining}
 \widehat HW=WH,\qquad e^{-\beta\widehat H}W=We^{-\beta H}
 \quad(\beta\in\R).
\end{equation}
\end{corollary}

\begin{proof}
For each $B_i\succeq0$, let $\lambda_i=\lambda_{\max}(B_i)$ and use
the exact identity
\begin{equation}\label{spin:casimir-extension}
 (S_i^\perp)^{\mathsf T}B_iS_i^\perp
 =(S_i^\perp)^{\mathsf T}(B_i-\lambda_iI_2)S_i^\perp
   -\lambda_i(S_i^z)^2+\lambda_i j_i(j_i+1)I.
\end{equation}
Replace each spin in this expression and in
\eqref{spin:hamiltonian-form} by
$\widehat S_i^a=\frac12\sum_{r=1}^{\kappa_i}\sigma_{ir}^a$.
Retain the last term of \eqref{spin:casimir-extension} as the indicated
scalar.  Interblock pairs have transverse matrix $J_{ij}/4$ and
longitudinal coefficient $e_{ij}/4$.  Intrablock pairs have transverse
matrix $(B_i-\lambda_iI_2)/2$ and longitudinal coefficient
$-\lambda_i/2$.  If $b_i=\lambda_{\min}(B_i)$, then
\[
 \|B_i-\lambda_iI_2\|_{\mathrm{op}}=\lambda_i-b_i\leq\lambda_i.
\]
All pairs therefore satisfy the qubit cone.  Identical-slot symmetric
off-diagonal products vanish by Pauli anticommutation; the remaining
identical-slot terms are scalars.  Each collective spin preserves the
symmetric sector and intertwines the physical spin through $W_i$.
This proves both identities in \eqref{spin:intertwining}.
\end{proof}

The reduction of higher spins to spin one-half and the corresponding
anisotropic sufficient Lee--Yang theorems are classical; see, in
particular, \cite[Theorem~2 and equations~(2.38)--(2.44)]{SuzukiFisher71}.
The corollary supplies explicit coefficients
for every Hamiltonian in the necessary-and-sufficient cone, including
single-ion quadrupoles.  It is an invariant-sector identity and does not
assert a spectral comparison with the other qubit spin sectors.



\appendix
\section{Exact arithmetic and bit bounds}\label{arith:section}

The algorithms in Section~\ref{alg:section} use rational recursions.
We supply the denominator and intermediate-magnitude bounds that
convert their arithmetic-operation counts into bit-complexity bounds.

\subsection{Eigenvalue continuation}\label{alg:energy-bit-details}

Use the notation of Section~\ref{alg:section} and the coefficient
order \(p\) chosen in the proof of Theorem~\ref{alg:complex}.
Let \(B_0\) bound the bit length of each numerator and denominator
in the original local matrices.
For \(p>0\), the interpolation nodes
\(s_\ell=-1+2\ell/p\) have \(O(\log(p+1))\) bits.
Taking local Hermitian parts and forming \(B+s_\ell C\)
therefore produces Gaussian-rational inputs with bit lengths
polynomial in \(B_0\) and \(\log(p+1)\).

Consider one Hermitian coefficient computation at these nodes.
Let \(Q_0\) be the product of the positive denominators of the
real and imaginary parts of its local entries.  There are
\(O_D(N)\) such entries, so \(\log Q_0\) is polynomial in the
input length.  Set
\begin{equation}
 K_p=48(p+1)!,\qquad
 D_j=(Q_0K_p)^j\quad(0\le j\le p).
 \label{alg:common-denominator}
\end{equation}
We recall the properties of the coefficient recursion used in
Sections~3--5 of \cite{BDL08}.  An order-\(j\) creation
coefficient \(C_j(M)\) is zero unless \(M\) is contained in a
connected set of at most \(j+1\) vertices.  Each nonzero term in
its recursion contains one local matrix entry, previous creation
coefficients of total order \(j-1\), and a divisor
\[
 k!\,E_0(M),\qquad k\le4,\qquad E_0(M)=2|M|.
 \]
The basis creation matrices
\(\prod_{u\in M}|1\rangle\langle0|_u\) have entries zero or one;
expanding their nested commutators therefore produces signed sums
of local perturbation entries.
The corresponding energy recursion has the same form without
the outer divisor \(E_0(M)\).
These are the explicit recursions used to implement the published
coefficient algorithm, rather than an assumption about arbitrary
symbolic perturbation expressions.

For \(j\le p\), every product \(k!E_0(M)\) just displayed
divides \(K_p\).  Induction in \(j\) consequently shows that
\begin{equation}
 D_j C_j(M)\in\mathbb Z[i],
 \qquad D_j e_j(B+s_\ell C)\in\mathbb Z[i].
 \label{alg:denominator-induction}
\end{equation}
Indeed, the previous coefficients contribute denominator
dividing \((Q_0K_p)^{j-1}\), the new local entry contributes
a factor dividing \(Q_0\), and the remaining divisor divides
\(K_p\).  All terms at a fixed order may be accumulated over
this common denominator.  Its bit length is
\[
 O\bigl(p\log Q_0+p^2\log(p+1)\bigr).
 \]

Intermediate numerators have polynomial bit length as well.
Here a deliberately loose bound suffices.
Fully unfold a term of perturbation order at most \(p\).
It has at most \(p\) local insertions and bounded-arity recursion
nodes.  Choose its local edges in at most \((ND)^{O(p)}\)
ways.  Their union has \(O(p)\) vertices, and all intermediate
creation sets lie in this union.  Choosing these sets at
\(O(p)\) nodes, together with the tree shape, order partitions,
and bounded-depth commutator terms, gives at most
\[
 \exp\bigl(O_D(p^2+p\log(N+1))\bigr)
 \]
fully unfolded summands.  Each summand contains at most \(p\)
bounded local entries; the positive integer energy and factorial
divisors cannot increase its absolute value.
Its magnitude is therefore at most
\(\max(1,2J)^{O(p)}\) times a fixed exponential factor.
After introducing \eqref{alg:common-denominator}, this bounds
the logarithmic magnitude of every intermediate numerator by
a polynomial in \(p,N,B_0\).
This argument concerns bit sizes only.  The algorithm does not
unfold these trees: its sharper operation count is the inherited
\(N\exp(O_D(p))\) count.

Exact interpolation also has polynomial bit cost.
For the nodes \(s_\ell=-1+2\ell/p\), its weights at \(i\)
can be written
\begin{equation}
 L_\ell(i)=
 \frac{\displaystyle
       \prod_{\substack{0\le r\le p\\r\ne\ell}}
           (pi+p-2r)}
      {2^p(-1)^{p-\ell}\ell!(p-\ell)!}.
 \label{alg:interpolation-weights}
\end{equation}
Both numerator and denominator have
\(O(p\log(p+1))\) bits.
The coefficients \(e_k(V)\) are obtained by multiplying these
weights by the computed node values and adding \(p+1\) terms.
The rational powers and binomial coefficients in
\eqref{alg:coefficient-transform}, and the powers in
\eqref{alg:stopping}--\eqref{alg:output}, increase bit length
only polynomially.
Thus a safe total bit bound is
\[
 N\exp(O_D(p))\operatorname{poly}(N,B_0,p,L_{\mathrm{in}}),
 \]
where \(L_{\mathrm{in}}\) is the total input bit length,
including the requested accuracy.
This proves the polynomial bit overhead used in Theorem~\ref{alg:complex}.

\subsection{Connected principal-state tests}\label{alg:test-bit-details}

We use the subsystem coefficients and notation from the proof of
Lemma~\ref{alg:connected-tests}.  Let \(Q_0\) be a common positive denominator for all local
entries of \(V\), let \(Q_B\) be the product of the denominators
of the entries of the \(B_i\), and put
\[
 D_p=2(p+1)!,\qquad T=Q_0D_p .
 \]
In a subsystem of size at most \(p+1\), each nonzero diagonal
entry of the reduced inverse \(S_U\) is \(1/(2r)\),
\(1\le r\le p+1\).  Induction in the vector recursion gives
denominators dividing \(T^{2k}\) for \(v_{U,k},w_{U,k}\)
and \(T^{2k-1}\) for the energy coefficient.
The denominator of \(a_k\) divides \(Q_BT^{2k}\).
Equation~\eqref{alg:formal-test-log} then bounds the denominator
of \(f_k\) by \(k!Q_B^kT^{2k}\).
At each order, accumulate the sums over these common
denominators; their logarithms are polynomial in
\(p\) and the total input length.

Intermediate numerators also have polynomial bit length.
For a subsystem \(U\), \(\|V_U\|\le m(d+1)J\).
On \(|x|\le r_U=[16(1+m(d+1)J)]^{-1}\), the resolvent
Neumann series on the circle \(|z|=1\) bounds the spectral
projection at the vacuum and its difference from the
unperturbed projection.  In particular, the projection has
norm at most \(16/15\), its vacuum entry differs from one by
at most \(1/15\), and both vacuum-normalized eigenvectors have
norm less than two.  Cauchy's estimate gives
\(\|v_{U,k}\|,\|w_{U,k}\|\le2r_U^{-k}\).
As \(\|B_U\|\le2^m\),
\[
 |a_k|\le4(k+1)2^m r_U^{-k}.
\]
For \(k\le p\), choose \(A_p=4(p+1)2^{p+1}\).
The finite formula
\(\log(1+a)=\sum_{\ell=1}^p(-1)^{\ell+1}a^\ell/\ell
 \pmod{x^{p+1}}\)
bounds \(|f_k|\) by \((2A_p)^kr_U^{-k}\).
Matrix products, vector sums and the partial sums in
\eqref{alg:formal-test-log} have the same type of bound, up to
factors \(2^{O(p)}\operatorname{poly}(p)\).
Subset sums add at most \(2^{p+1}\) terms, and the final
connected sum has at most \(n\exp(O_d(p))\) terms.
Combining these magnitude bounds with the common denominators
controls every stored numerator, including before cancellation.
This proves the bit complexity.


\section{Connected preservers and the obstruction to flow generation}
\label{glob:section}

The exact infinitesimal cone does not generate every connected static
preserver, even after taking limits of products.  We prove this on the
fixed multiaffine space
$V_n=\C[z_1,\ldots,z_n]_{\mathrm{ma}}\simeq(\C^2)^{\otimes n}$.
Let $S_n$ be the semigroup of invertible complex-linear operators that
preserve disk stability.  For its ordinary monomial coefficient matrix
$M$, write
\begin{equation}\label{glob:symbol}
 G_M(z,w)=\sum_{a,b\in\{0,1\}^n}M_{ab}z^aw^b.
\end{equation}
For \(M\in\operatorname{GL}(V_n)\), the symbol criterion is
\begin{equation}\label{glob:symbol-equivalence}
 M\in S_n\quad\Longleftrightarrow\quad
 G_M\text{ is stable on }\D^{2n}.
\end{equation}
For necessity, fix $w\in\D^n$ and apply $M$ to the stable polynomial
$\prod_i(1+w_iz_i)$.  Invertibility excludes the zero output.
Conversely, tensor contraction of $G_M$ with a stable input proves
preservation, again excluding zero by invertibility.  This is the
multiaffine disk-symbol convention of the preceding sections.
Hurwitz's theorem shows that $S_n$ is closed relative to
$\operatorname{GL}(V_n)$: an invertible limit cannot have identically
zero symbol.  Nonzero scalar multiplication preserves $S_n$.

Define
\begin{equation}\label{glob:semigroups}
 \begin{split}
 W_n&=\{A:e^{tA}\in S_n\text{ for every }t\geq0\},\\
 E_n&=\langle e^A:A\in W_n\rangle,\qquad
 T_n=\overline{E_n}^{\,\operatorname{GL}(V_n)},\qquad
 H_n=S_n\cap S_n^{-1}.
 \end{split}
\end{equation}
Here $E_n$ consists of finite products.  The closure $T_n$ is a
different semigroup.  The group $H_n$ is the group of reversible
preservers, and $H_n^0$ denotes its identity component.

\begin{theorem}[Global connectedness and a strict flow obstruction]
\label{glob:main}
For every $n\geq1$, $S_n$ is path connected.  Nevertheless,
\begin{equation}\label{glob:reversible-parts}
 E_n\cap H_n=H_n^0,
 \qquad T_n\cap T_n^{-1}=H_n^0.
\end{equation}
For $n\geq2$, every nontrivial coordinate permutation $\pi$ satisfies
\begin{equation}\label{glob:permutation-obstruction}
 \pi\in S_n\setminus T_n.
\end{equation}
Thus a positive-radius matrix neighborhood of $\pi$ contains no finite
product of admissible one-parameter flows, even with unbounded product
lengths and generator norms.
\end{theorem}

\begin{proof}[Path connectedness]
Let $D_rz^a=r^{|a|}z^a$, $0<r\leq1$.  It is an invertible preserver,
and
\begin{equation}\label{glob:sandwich}
 G_{D_rMD_r}(z,w)=G_M(rz,rw).
\end{equation}
For $M\in S_n$, its constant coefficient $M_{00}$ is nonzero.
A path of nonzero scalars normalizes it to one.  With this normalization,
$D_rMD_r$ converges as $r\downarrow0$ to the vacuum projector
$|0\rangle\langle0|$.

Consider the convex open set in the complex affine hyperplane $X_{00}=1$
given by
\begin{equation}\label{glob:coefficient-ball}
 \mathcal B=\left\{X:X_{00}=1,
        \ \sum_{(a,b)\ne(0,0)}|X_{ab}|<1\right\}.
\end{equation}
Every member has a kernel nonvanishing on the closed polydisk, by strict
constant-term domination.  For small positive $r,s$, both $D_rMD_r$
and $D_s$ belong to $\mathcal B$ and are invertible.  Indeed the
nonconstant coefficient sum for $D_s$ is $(1+s)^n-1$.

The invertible part of $\mathcal B$ is path connected.  To see this,
take invertible $X,Y\in\mathcal B$ and set
$L(\zeta)=(1-\zeta)X+\zeta Y$.  Its determinant is a nonzero
polynomial in $\zeta$, since it is nonzero at zero.  The real segment
$L([0,1])$ is compactly contained in $\mathcal B$, so a sufficiently
thin complex neighborhood of $[0,1]$ still maps into $\mathcal B$.
A path from $0$ to $1$ in this neighborhood avoids the finitely many
determinant zeros.  Its image connects $X$ to $Y$ through invertible
members of $\mathcal B$.
Now connect $M$ to $D_rMD_r$ using \eqref{glob:sandwich}, connect this
point to $D_s$ inside $\mathcal B$, and connect $D_s$ to $I$ by
increasing $s$ to one.  Every point is in $S_n$.  The singular vacuum
projector is used only as a limit to enter $\mathcal B$; it is never a
point of the path.
\end{proof}

For reference, there is also an explicit two-variable path.  On
$\operatorname{span}(z_1,z_2)$ let $P_\pm$ be the symmetric and
antisymmetric projections, and let $U_\theta$ act as
$P_++e^{i\theta}P_-$ there and as the identity on
$\operatorname{span}(1,z_1z_2)$.  For $r=1/4$,
\begin{align}\label{glob:swap-path}
 G_{D_rU_\theta}
 ={}&1+ra_\theta(z_1w_1+z_2w_2)
       +rb_\theta(z_1w_2+z_2w_1)
       +r^2z_1z_2w_1w_2,\notag\\
 a_\theta={}&(1+e^{i\theta})/2,\qquad
 b_\theta=(1-e^{i\theta})/2.
\end{align}
Its nonconstant coefficient sum is at most
$2\sqrt2r+r^2<1$, and its determinant is $r^4e^{i\theta}\ne0$.
The paths $I\to D_r$, $D_rU_\theta$ for $0\leq\theta\leq\pi$,
and $D_s\pi$ for $r\leq s\leq1$ connect the identity to the
coordinate exchange.  These paths are not asserted to have preserving
forward increments.

\begin{lemma}[Reversible factors]\label{glob:unit-factors}
If $A,B\in S_n$ and $AB\in H_n$, then $A,B\in H_n$.
Moreover every element of $H_n^0$ is a nonzero scalar times a tensor
product of one-site invertible matrices, whereas no nontrivial coordinate
permutation has this form.
\end{lemma}

\begin{proof}
The identities $A^{-1}=B(AB)^{-1}$ and $B^{-1}=(AB)^{-1}A$ prove
the first assertion.  The group $H_n$ is a closed subgroup of
$\operatorname{GL}(V_n)$ and hence a real Lie group.  Its Lie algebra
is $W_n\cap(-W_n)$.  By Corollary~\ref{gen:gram-cones}, in
half-plane coordinates this algebra is
\begin{equation}\label{glob:unit-algebra}
 \C I+\bigoplus_i\mathfrak{sl}_2(\R)_i.
\end{equation}
The fixed tensor Cayley transformation preserves products and
coordinate permutations.  A connected Lie group is generated by the
exponentials of its Lie algebra.  Exponentiating
\eqref{glob:unit-algebra} and taking finite products therefore gives
only scalar multiples of local tensor products.
A nontrivial coordinate permutation cannot have this form: varying one
input factor varies a different output factor, whereas an invertible
local tensor product varies only the corresponding output factor.
Equivalently, for a two-factor exchange, an identity
$Ax\otimes By=c\,y\otimes x$ with fixed nonzero $x$ and varying
independent $y$ would force the fixed vector $Ax$ to be proportional
to every $y$.
\end{proof}

\begin{proof}[The exact-product assertion in Theorem~\ref{glob:main}]
If $e^{A_m}\cdots e^{A_1}\in H_n$, with every $A_j\in W_n$,
Lemma~\ref{glob:unit-factors} makes every exponential factor a unit.
For $0\leq s\leq1$, factor
$e^{A_j}=e^{(1-s)A_j}e^{sA_j}$.  Both factors preserve stability,
so the same lemma gives $e^{sA_j}\in H_n$.  Thus each factor belongs
to $H_n^0$, and so does their product.  Conversely $H_n^0$ is
generated by exponentials of its Lie algebra, which is contained in
$W_n$.  This proves $E_n\cap H_n=H_n^0$.
\end{proof}

For the closure assertion we use Neeb's classical connected-unit theorem:
if $G$ is a connected Lie group and $T\subseteq G$ is a closed
submonoid satisfying
\[
 T=\overline{\langle\exp L(T)\rangle}^{\,G},\qquad
 L(T)=\{A:\exp(tA)\in T\text{ for every }t\geq0\},
\]
then $T\cap T^{-1}$ is connected. This is
\cite[Definition~II.1(3) and Proposition~III.2, pp.~171, 178]{Neeb92};
see also \cite[Section~6.1]{SanMartinSantana02}.

\begin{proof}[The closure assertion in Theorem~\ref{glob:main}]
The ambient real Lie group $G=\operatorname{GL}_{2^n}(\C)$ is connected.
By definition $T_n$ is a closed submonoid and $T_n\subseteq S_n$.
Every $A\in W_n$ generates a one-parameter semigroup in $E_n$, while
$L(T_n)\subseteq L(S_n)=W_n$. Hence
\[
 L(T_n)=W_n,\qquad
 T_n=\overline{\langle\exp L(T_n)\rangle}^{\,G}.
\]
Neeb's theorem therefore makes $T_n\cap T_n^{-1}$ a connected
subgroup of $H_n$, so it is contained in $H_n^0$.
The reverse inclusion follows because the exponentials of both signs of
\eqref{glob:unit-algebra} belong to $T_n$ and generate $H_n^0$.

If a nontrivial coordinate permutation $\pi$ belonged to $T_n$, its
finite order would imply $\pi^{-1}\in T_n$.  It would therefore be
a unit of $T_n$, and hence belong to $H_n^0$, contrary to
Lemma~\ref{glob:unit-factors}.  This proves
\eqref{glob:permutation-obstruction}.  Relative closure in the general
linear group and ordinary matrix-norm closure agree at the invertible
point $\pi$, proving the neighborhood statement.
\end{proof}

\begin{corollary}[Exact local divisibility]\label{glob:local-divisibility}
Call $M\in S_n$ exactly locally divisible if, for every identity
neighborhood $U$ in $\operatorname{GL}(V_n)$, it is a finite product
of members of $S_n\cap U$.  The exactly locally divisible units of
$S_n$ are precisely $H_n^0$.  In addition, if a continuous path
$M:[0,1]\to S_n$ starts at $I$ and satisfies
$M(t)M(s)^{-1}\in S_n$ for $s\leq t$, then a unit endpoint must
belong to $H_n^0$.
\end{corollary}

\begin{proof}
The identity component of a Lie group is open, so choose $U$ with
$H_n\cap U\subset H_n^0$.  Any exact factorization of a unit into
$S_n\cap U$ consists entirely of units by
Lemma~\ref{glob:unit-factors}, hence has endpoint in $H_n^0$.
Conversely, subdividing a finite exponential representation of a member
of $H_n^0$ puts all its factors in any prescribed identity neighborhood.
For the path assertion, write
$M(1)=[M(1)M(t)^{-1}]M(t)$.  A unit endpoint makes $M(t)$ a unit
for every $t$, so the continuous path stays in $H_n$ and ends in
$H_n^0$.
\end{proof}

\begin{remark}\label{glob:scope}
The identity for the closure in \eqref{glob:reversible-parts} is
$T_n\cap T_n^{-1}=H_n^0$.  It does not assert
$T_n\cap H_n=H_n^0$ for arbitrary units of the larger semigroup;
finite order supplies the inverse inclusion needed for permutations.
The result concerns flows of the infinitesimal cone.  Static apolar maps
may lie on other branches: for $F(s,t)=s-t$,
$B_F=\pi-I$, so $I+B_F=\pi$, whereas
\[
 e^{tB_F}=\frac{1+e^{-2t}}2I+\frac{1-e^{-2t}}2\pi.
\]
Allowing such a static permutation as an additional primitive changes
the factorization question.  No description of the remaining
noninvertible boundary or the nonunit part of $T_n$ is claimed.
\end{remark}


\section{An exact counterexample to a stronger ground-state radius}\label{cx:section}

The qualitative strict radius in Theorem~\ref{eq:strict-field} cannot
in general be replaced by the proposed explicit radius for deformed
EPR ground states. We address the stronger clause of Conjecture~1 in
\cite[Section 5.4]{WBG26}, with its stated unweighted-graph hypothesis.

\begin{theorem}\label{cx:epr}
Let $G$ have vertices $\{0,1,2,3,4,5\}$ and edges
\[
 \{01,02,03,14,15\}.
\]
Put
\[
 H=-\sum_{ij\in E(G)}
 (|00\rangle+\tfrac12|11\rangle)
 (\langle00|+\tfrac12\langle11|)_{ij}.
\]
The ground state $\psi$ is unique up to scalar. Its multiaffine
polynomial $f_\psi(z)=\sum_b\psi_bz^b$ has a zero satisfying
\[
 |z_i|<\frac{499}{500}\sqrt2\qquad(0\leq i\leq5).
\]
Thus the proposed radius $s^{-1/2}$ fails at $s=1/2$ on an
unweighted six-vertex tree.
\end{theorem}

\begin{proof}
We give rational spectral and zero certificates. All inequalities
below are verified by the finite procedure at the end of this
appendix, using only integer arithmetic and exact fractions.

Write $K=-H$. Index the computational basis by integers
$b\in\{0,\ldots,63\}$, with bit $b_i$ at vertex $i$.
For each edge $ij$ and each input with $b_i=b_j=a$, changing both
bits to $c\in\{0,1\}$ contributes $2^{-(a+c)}$ to the corresponding
entry of $K$. All other entries from that edge are zero.
This specifies the symmetric rational matrix $K$ completely.

Consider $M=37I-8K$. The charge
\[
 Q(b)=b_0+b_4+b_5-b_1-b_2-b_3
\]
is conserved. In charge order $-3,-2,\ldots,3$, the block dimensions are
\[
 1,\ 6,\ 15,\ 20,\ 15,\ 6,\ 1.
\]
Order each block by increasing integer $b$.
Successive rational symmetric elimination has pivot
signs as follows:
\[
\begin{array}{c|rrrrrrr}
 Q&-3&-2&-1&0&1&2&3\\ \hline
 \text{negative pivots}&0&0&0&1&0&0&0\\
 \text{positive pivots}&1&6&15&19&15&6&1.
\end{array}
\]
No pivot is zero; the negative pivot is the first in the charge-zero
block. For clarity, the recurrence at step $i$ is
\[
 p_i=M^{(i)}_{ii},\qquad
 M^{(i+1)}_{jk}
 =M^{(i)}_{jk}-M^{(i)}_{ji}M^{(i)}_{ik}/p_i
 \quad(j,k>i).
\]
It proves the stated inertia by exact rational comparison, rather
than taking the pivot signs as external data.
Sylvester's law of inertia shows that $K$ has exactly one
eigenvalue above $37/8$. Since $K_{00}=5$, its largest eigenvalue
is at least $5$ and is simple. Every other eigenvalue is below
$37/8$.

Define a rational vector $v$ as follows. For a bit string use the key
\[
 (a,b,i,j)=(b_0,b_1,b_2+b_3,b_4+b_5).
\]
The amplitude for each individual bit string is the following
numerator divided by $10^{15}$; a missing key has amplitude zero.
\begin{center}
\begin{tabular}{cr}
\toprule
Key & Numerator\\
\midrule
$(0,0,0,0)$ & $1000000000000000$\\
$(0,0,1,1)$ & $8942179087130$\\
$(0,0,2,2)$ & $136748346620$\\
$(0,1,0,1)$ & $171232824560490$\\
$(0,1,1,2)$ & $2452435031309$\\
$(1,0,1,0)$ & $171232824560490$\\
$(1,0,2,1)$ & $2452435031309$\\
$(1,1,0,0)$ & $97255575326164$\\
$(1,1,1,1)$ & $37858301388951$\\
$(1,1,2,2)$ & $1200949534669$\\
\bottomrule
\end{tabular}
\end{center}
These are computational-basis amplitudes, not normalized orbit-basis
coordinates. Put
\[
 \widehat\lambda=\frac{539109343678406}{10^{14}}.
\]
Direct rational summation gives
\[
 \|(K-\widehat\lambda I)v\|_2^2<10^{-24},
 \qquad \widehat\lambda-37/8>1/2.
\]
Let $v_g$ be the orthogonal projection onto the largest eigenspace
of $K$. Expanding in an orthonormal eigenbasis and using the
preceding separation gives
\[
 \|v-v_g\|_2
 \leq2\|(K-\widehat\lambda I)v\|_2<2\cdot10^{-12}.
\]
Since $v_{000000}=1$, $v_g$ is nonzero, hence is a scalar multiple
of the true ground state. Both $v$ and $v_g$ have charge zero;
in particular they have even parity.

For an even-parity vector $w$ put
\[
 \widetilde f_w(x)=f_w(\sqrt2\,x)
        =\sum_b2^{|b|/2}w_bx^b.
\]
The coefficients of $\widetilde f_v$ are rational. Set
\[
 c=\frac{175105+981180\,\mathrm i}{10^6},\qquad
 \ell=\frac{-379834+921467\,\mathrm i}{10^6}.
\]
Exact squaring gives $|c|,|\ell|<499/500$. Fix
$x_0=x_1=c$ and $x_2=x_3=x_4=\ell$. Multiaffinity yields
\[
 \widetilde f_v(c,c,\ell,\ell,\ell,x_5)=A_v+B_vx_5,
 \qquad
 |A_v|^2<(85/1000)^2,\quad |B_v|^2>(852/10000)^2 .
\]
The displayed inequalities are finite sums over the amplitude table.
Each coefficient $A$ or $B$ contains at most $32$ terms. A fixed-variable
monomial has modulus at most one, and its even tilt factor is at most
eight. Therefore
\[
 |A_{v_g}-A_v|,\ |B_{v_g}-B_v|
 <32\cdot8\cdot2\cdot10^{-12}<10^{-8}.
\]
In particular $B_{v_g}\neq0$, and
\[
 x_5=-A_{v_g}/B_{v_g},\qquad
 |x_5|<
 \frac{85/1000+10^{-8}}{852/10000-10^{-8}}<499/500.
\]
The last inequality is rational. This gives a zero of
$\widetilde f_{v_g}$ with all coordinates of modulus below $499/500$.
Rescaling by $\sqrt2$ proves the theorem.
\end{proof}

The certificate locates a zero of the exact ground-state polynomial.
It does not treat a rounded eigenvector as an exact eigenvector.
The existence of some radius greater than one, proved above by
positive-field strictification, is compatible with this example.

\subsection*{Complete finite verification procedure}
The following standard-library Python procedure reproduces all
rational assertions used in the proof. The operations implementing
a Gaussian rational number are written explicitly. Floating-point
arithmetic and numerical eigenvalue computations are not used.
The same source is supplied as
\path{certificates/epr_six_vertex_certificate.py}; a second,
independent rational certificate is included with the source package.
\begin{lstlisting}[basicstyle=\ttfamily\scriptsize,breaklines=true,
 columns=fullflexible,keepspaces=true,showstringspaces=false,
 language=Python,frame=single]
"""Exact standard-library certificate: a zero strictly inside radius sqrt(2).

No floating point or numerical eigensolver is used in any assertion.
K is the negative EPR Hamiltonian at s=1/2 on edges 01,02,03,14,15.
"""
from fractions import Fraction as F

EDGES = [(0, 1), (0, 2), (0, 3), (1, 4), (1, 5)]
N = 64
K = [[F(0) for _ in range(N)] for _ in range(N)]
for b in range(N):
    for i, j in EDGES:
        bi, bj = (b >> i) & 1, (b >> j) & 1
        if bi == bj:
            for out in (0, 1):
                c = (b & ~(1 << i) & ~(1 << j)) | (out << i) | (out << j)
                K[c][b] += F(1, 2) ** (bi + out)
assert all(K[i][j] == K[j][i] for i in range(N) for j in range(N))

# Exact inertia proves that all but the largest eigenvalue are below 37/8.
def charge(b):
    return sum((b >> j) & 1 for j in (0, 4, 5)) - sum((b >> j) & 1 for j in (1, 2, 3))

assert all(not K[i][j] or charge(i) == charge(j) for i in range(N) for j in range(N))

negative = 0
pivot_signs = []
for q in range(-3, 4):
    sector = [b for b in range(N) if charge(b) == q]
    A = [[F(37 if i == j else 0) - 8 * K[i][j] for j in sector] for i in sector]
    signs = []
    for i in range(len(sector)):
        pivot = A[i][i]
        assert pivot != 0
        negative += pivot < 0
        signs.append('+' if pivot > 0 else '-')
        for j in range(i + 1, len(sector)):
            for k in range(j, len(sector)):
                A[k][j] -= A[j][i] * A[k][i] / pivot
                A[j][k] = A[k][j]
    pivot_signs.append((q, ''.join(signs)))
assert negative == 1
assert K[0][0] == 5

states = [(0,0,0,0),(0,0,1,1),(0,0,2,2),(0,1,0,1),(0,1,1,2),
          (1,0,1,0),(1,0,2,1),(1,1,0,0),(1,1,1,1),(1,1,2,2)]
numerators = [1000000000000000,8942179087130,136748346620,171232824560490,
              2452435031309,171232824560490,2452435031309,97255575326164,
              37858301388951,1200949534669]
amps = {state:F(num,10**15) for state,num in zip(states,numerators)}
v = []
for b in range(N):
    state = ((b >> 0)&1,(b >> 1)&1,((b >> 2)&1)+((b >> 3)&1),
             ((b >> 4)&1)+((b >> 5)&1))
    v.append(amps.get(state,F(0)))
assert v[0] == 1
lam = F(539109343678406,10**14)
r = [sum(K[i][j] * v[j] for j in range(N)) - lam*v[i] for i in range(N)]
assert sum(x*x for x in r) < F(1,10**24)
assert lam - F(37,8) > F(1,2)

# Therefore the orthogonal projection v_g onto the true top eigenspace
# obeys ||v-v_g||_2 < 2*10^-12. In particular v_g is nonzero.

def add(x,y): return x[0]+y[0],x[1]+y[1]
def mul(x,y): return x[0]*y[0]-x[1]*y[1],x[0]*y[1]+x[1]*y[0]
def scale(c,x): return c*x[0],c*x[1]
def norm2(x): return x[0]*x[0]+x[1]*x[1]

zc=(F(175105,10**6),F(981180,10**6))
zl=(F(-379834,10**6),F(921467,10**6))
certified_radius=F(499,500)
assert norm2(zc)<certified_radius**2 and norm2(zl)<certified_radius**2
fixed=[zc,zc,zl,zl,zl]
A=(F(0),F(0));B=(F(0),F(0))
for b in range(N):
    if not v[b]: continue
    degree=b.bit_count()
    assert degree%2==0
    term=(v[b]*2**(degree//2),F(0))
    for j in range(5):
        if (b>>j)&1: term=mul(term,fixed[j])
    if (b>>5)&1:B=add(B,term)
    else:A=add(A,term)

# These wide rational bounds leave much more room than the eigenvector error.
assert norm2(A) < F(85,1000)**2
assert norm2(B) > F(852,10000)**2

# At the five fixed coordinates each monomial has modulus<1, and tilting
# by sqrt(2) multiplies each even coefficient by at most8. Each linear
# coefficient A or B involves at most32 terms. Thus replacing v by v_g
# changes A and B by less than 32*8*2e-12 < 1e-8.
err=F(1,10**8)
assert 32*8*F(2,10**12) < err
assert F(85,1000)+err < F(852,10000)-err
assert F(85,1000)+err < certified_radius*(F(852,10000)-err)

print('PASS: exact inertia has one negative pivot:',pivot_signs)
print('PASS: squared eigenvector residual < 10^-24')
print('PASS: |A_v|<85/1000 and |B_v|>852/10000')
print('PASS: true ground polynomial has all |z_i|<499/500 after sqrt(2) scaling')
\end{lstlisting}


\section*{Disclosure of artificial intelligence use}
Most of the mathematical content, proofs, computational certificates,
and text in this article were generated with substantial assistance from
OpenAI GPT-6 Pro in ChatGPT chat mode and GPT-6 Astra in Codex, both
using the ultra reasoning setting. These were the only AI tools used.
The author organized and directed the research through repeated
interactions, including questions, choices of research directions,
requests for extensions and critical review, and selection and
organization of the material. AI assistance was also used for
literature searches, proof audits, exact computational checks, and
manuscript revision. These checks do not constitute complete
proof-assistant verification, which is not claimed.
The author takes full responsibility for the accuracy, originality,
attribution, and conclusions of the article.


\makeatletter\def\@biblabel#1{#1.}\makeatother
\providecommand{\bysame}{\leavevmode\hbox to3em{\hrulefill}\thinspace}
\providecommand{\MR}{\relax\ifhmode\unskip\space\fi MR }
% \MRhref is called by the amsart/book/proc definition of \MR.
\providecommand{\MRhref}[2]{%
  \href{http://www.ams.org/mathscinet-getitem?mr=#1}{#2}
}
\providecommand{\href}[2]{#2}
\begin{thebibliography}{10}

\bibitem{ALMPSS26}
Anuj Apte, Eunou Lee, Kunal Marwaha, Ojas Parekh, Lennart Sinjorgo, and James
  Sud, \emph{A 0.8395-approximation algorithm for the {EPR} problem}, 34th
  Annual European Symposium on Algorithms ({ESA} 2026), Leibniz International
  Proceedings in Informatics (LIPIcs), vol. 388, Schloss
  Dagstuhl--Leibniz-Zentrum f{\"u}r Informatik, 2026, pp.~144:1--144:13.

\bibitem{APS26}
Anuj Apte, Ojas Parekh, and James Sud, \emph{Conjectured bounds for 2-local
  {Hamiltonians} via token graphs}, arXiv:2506.03441v2, 2026, Version 2, 27 May
  2026; first version 3 June 2025.

\bibitem{AsanoPRL70}
Taro Asano, \emph{{Lee--Yang} theorem and the {Griffiths} inequality for the
  anisotropic {Heisenberg} ferromagnet}, Phys. Rev. Lett. \textbf{24} (1970),
  1409--1411.

\bibitem{AsanoJapanese70}
\bysame, \emph{Rigorous theorems concerning {Heisenberg} ferromagnets}, Bussei
  Kenkyu \textbf{14} (1970), no.~2, 72--96, In Japanese; title translated from
  the original.

\bibitem{BBKL26}
Ainesh Bakshi, Arpon Basu, Pravesh Kothari, and Anqi Li, \emph{Sharp bounds on
  the eigenvalues of {Kikuchi} graphs and applications to {Quantum Max Cut}},
  arXiv:2605.14994v1, 2026, Version 1, 14 May 2026.

\bibitem{BasuPollackRoy06}
Saugata Basu, Richard Pollack, and Marie-Fran{\c c}oise Roy, \emph{Algorithms
  in real algebraic geometry}, 2 ed., Algorithms and Computation in
  Mathematics, vol.~10, Springer, Berlin, 2006.

\bibitem{BB09}
Julius Borcea and Petter Br{\"a}nd{\'e}n, \emph{The {Lee--Yang} and
  {P{\'o}lya--Schur} programs. {I}. {Linear} operators preserving stability},
  Invent. Math. \textbf{177} (2009), no.~3, 541--569.

\bibitem{BB10II}
\bysame, \emph{The {Lee--Yang} and {P{\'o}lya--Schur} programs. {II}. {Theory}
  of stable polynomials and applications}, Comm. Pure Appl. Math. \textbf{62}
  (2009), no.~12, 1595--1631.

\bibitem{BB09Circular}
\bysame, \emph{{P{\'o}lya--Schur} master theorems for circular domains and
  their boundaries}, Ann. of Math. (2) \textbf{170} (2009), no.~1, 465--492.

\bibitem{BB10}
\bysame, \emph{Multivariate {P{\'o}lya--Schur} classification problems in the
  {Weyl} algebra}, Proc. Lond. Math. Soc. (3) \textbf{101} (2010), no.~1,
  73--104.

\bibitem{Branden10Discrete}
Petter Br{\"a}nd{\'e}n, \emph{Discrete concavity and the half-plane property},
  SIAM Journal on Discrete Mathematics \textbf{24} (2010), no.~3, 921--933.

\bibitem{BDL08}
Sergey Bravyi, David~P. DiVincenzo, and Daniel Loss, \emph{Polynomial-time
  algorithm for simulation of weakly interacting quantum spin systems}, Comm.
  Math. Phys. \textbf{284} (2008), 481--507.

\bibitem{BG17}
Sergey Bravyi and David Gosset, \emph{Polynomial-time classical simulation of
  quantum ferromagnets}, Phys. Rev. Lett. \textbf{119} (2017), 100503.

\bibitem{BGLW26}
Sergey Bravyi, David Gosset, Yinchen Liu, and Benjamin Wong, \emph{Efficient
  quantum algorithm for {Heisenberg} spin systems}, arXiv:2607.14401, 2026,
  Version 1, July 15, 2026.

\bibitem{Burgarth19}
Daniel Burgarth, Paolo Facchi, Hiromichi Nakazato, Saverio Pascazio, and Kazuya
  Yuasa, \emph{Generalized adiabatic theorem and strong-coupling limits},
  Quantum \textbf{3} (2019), 152.

\bibitem{BurkeOverton01}
James~V. Burke and Michael~L. Overton, \emph{Variational analysis of the
  abscissa mapping for polynomials}, SIAM J. Control Optim. \textbf{39} (2001),
  no.~6, 1651--1676.

\bibitem{CravenCsordas94}
Thomas Craven and George Csordas, \emph{Differential operators of infinite
  order and the distribution of zeros of entire functions}, J. Math. Anal.
  Appl. \textbf{186} (1994), no.~3, 799--820.

\bibitem{DalfoEtAl21}
C.~Dalf{\'o}, F.~Duque, R.~Fabila-Monroy, M.~A. Fiol, C.~Huemer, A.~L.
  Trujillo-Negrete, and F.~J. Zaragoza~Mart{\'i}nez, \emph{On the {Laplacian}
  spectra of token graphs}, Linear Algebra Appl. \textbf{625} (2021), 322--348.

\bibitem{Dubois09}
Lo{\"i}c Dubois, \emph{Projective metrics and contraction principles for
  complex cones}, J. Lond. Math. Soc. (2) \textbf{79} (2009), no.~3, 719--737.

\bibitem{Dunlop79}
Fran{\c c}ois Dunlop, \emph{Analyticity of the pressure for {Heisenberg} and
  plane rotator models}, Comm. Math. Phys. \textbf{69} (1979), 81--88.

\bibitem{GaubertQu15}
St{\'e}phane Gaubert and Zheng Qu, \emph{{Dobrushin}'s ergodicity coefficient
  for {Markov} operators on cones}, Integral Equations Operator Theory
  \textbf{81} (2015), no.~1, 127--150.

\bibitem{HMS20}
Aram~W. Harrow, Saeed Mehraban, and Mehdi Soleimanifar, \emph{Classical
  algorithms, correlation decay, and complex zeros of partition functions of
  quantum many-body systems}, Proceedings of the 52nd Annual {ACM SIGACT}
  Symposium on Theory of Computing, ACM, 2020, pp.~378--386.

\bibitem{MarwahaSud26}
Kunal Marwaha and James Sud, \emph{A complexity phase transition at the {EPR
  Hamiltonian}}, arXiv:2604.13026, 2026.

\bibitem{Neeb92}
Karl-Hermann Neeb, \emph{On the foundations of {Lie} semigroups}, J. Reine
  Angew. Math. \textbf{431} (1992), 165--189.

\bibitem{PatelRegts17}
Viresh Patel and Guus Regts, \emph{Deterministic polynomial-time approximation
  algorithms for partition functions and graph polynomials}, SIAM J. Comput.
  \textbf{46} (2017), no.~6, 1893--1919.

\bibitem{PolyaSchur14}
Georg P{\'o}lya and Issai Schur, \emph{{\"U}ber zwei arten von faktorenfolgen
  in der theorie der algebraischen gleichungen}, J. Reine Angew. Math.
  \textbf{144} (1914), 89--113.

\bibitem{PRSS04}
Victoria Powers, Bruce Reznick, Claus Scheiderer, and Frank Sottile, \emph{A
  new approach to {Hilbert}'s theorem on ternary quartics}, C. R. Math. Acad.
  Sci. Paris \textbf{339} (2004), no.~9, 617--620.

\bibitem{Purbhoo18}
Kevin Purbhoo, \emph{Total nonnegativity and stable polynomials}, Canad. Math.
  Bull. \textbf{61} (2018), no.~4, 836--847.

\bibitem{RaghavendraRyderSrivastava17}
Prasad Raghavendra, Nick Ryder, and Nikhil Srivastava, \emph{Real stability
  testing}, 8th Innovations in Theoretical Computer Science Conference, Leibniz
  International Proceedings in Informatics, vol.~67, Schloss
  Dagstuhl--Leibniz-Zentrum f{\"u}r Informatik, 2017, pp.~5:1--5:15.

\bibitem{RT26}
Chaithanya Rayudu and Jun Takahashi, \emph{Spectral gap of {Lee--Yang
  Hamiltonians}}, arXiv:2607.10765, 2026, Version 1, July 12, 2026.

\bibitem{ReyesDalfoFiol24}
M.~A. Reyes, C.~Dalf{\'o}, and M.~A. Fiol, \emph{On the spectra and spectral
  radii of token graphs}, Bol. Soc. Mat. Mex. \textbf{30} (2024), Paper No. 11.

\bibitem{Ruelle10}
David Ruelle, \emph{Characterization of {Lee--Yang} polynomials}, Ann. of Math.
  (2) \textbf{171} (2010), no.~1, 589--603.

\bibitem{Rugh10}
Hans~Henrik Rugh, \emph{Cones and gauges in complex spaces: Spectral gaps and
  complex {Perron--Frobenius} theory}, Ann. of Math. (2) \textbf{171} (2010),
  no.~3, 1707--1752.

\bibitem{SanMartinSantana02}
Luiz A.~B. San~Martin and Alexandre~J. Santana, \emph{The homotopy type of
  {Lie} semigroups in semi-simple {Lie} groups}, Monatsh. Math. \textbf{136}
  (2002), 151--173.

\bibitem{SuzukiFisher71}
Masuo Suzuki and Michael~E. Fisher, \emph{Zeros of the partition function for
  the {Heisenberg}, ferroelectric, and general {Ising} models}, J. Math. Phys.
  \textbf{12} (1971), no.~2, 235--246.

\bibitem{Turbiner92}
Alexander Turbiner, \emph{{Lie}-algebraic approach to the theory of polynomial
  solutions. {I}. {Ordinary} differential equations and finite-difference
  equations in one variable}, arXiv:hep-th/9209079, 1992.

\bibitem{Turbiner94}
\bysame, \emph{{Lie}-algebras and linear operators with invariant subspaces},
  Lie Algebras, Cohomologies and New Findings in Quantum Mechanics (Niky Kamran
  and Peter~J. Olver, eds.), Contemporary Mathematics, vol. 160, American
  Mathematical Society, Providence, RI, 1994, Revised author version:
  arXiv:funct-an/9301001v2 (1998), pp.~263--310.

\bibitem{WeiQuantum2026}
Dongsheng Wei, \emph{{Lee--Yang} preserving quantum {Markov} semigroups and
  channels}, Companion manuscript, 2026.

\bibitem{WeiPositive2026}
\bysame, \emph{Positive stability-preserving semigroups and particle
  capacities}, Companion manuscript, 2026.

\bibitem{WildAlhambra23}
Dominik~S. Wild and {\'A}lvaro~M. Alhambra, \emph{Classical simulation of
  short-time quantum dynamics}, PRX Quantum \textbf{4} (2023), 020340.

\bibitem{WongTalk23}
Benjamin Wong, \emph{Quantum {Max-Cut} on bipartite graphs}, URA Seminar,
  University of Waterloo, August 14, 2023, Public talk abstract.

\bibitem{WongThesis26}
\bysame, \emph{{Lee--Yang} tensors in quantum information and algorithms for
  local {Hamiltonian} problems}, Master's thesis, University of Waterloo, 2026.

\bibitem{WBG26}
Benjamin Wong, Sergey Bravyi, David Gosset, and Yinchen Liu, \emph{{Lee--Yang}
  tensors and {Hamiltonian} complexity}, arXiv:2602.03605, 2026, Version 1,
  February 3, 2026.

\bibitem{YYZ22}
Penghui Yao, Yitong Yin, and Xinyuan Zhang, \emph{Polynomial-time approximation
  of zero-free partition functions}, 49th International Colloquium on Automata,
  Languages, and Programming ({ICALP} 2022) (Dagstuhl, Germany) (Miko{\l}aj
  Boja\'{n}czyk, Emanuela Merelli, and David~P. Woodruff, eds.), Leibniz
  International Proceedings in Informatics, vol. 229, Schloss Dagstuhl --
  Leibniz-Zentrum f{\"u}r Informatik, 2022, pp.~108:1--108:20.

\end{thebibliography}

\enlargethispage{3\baselineskip}
\end{document}

